%% ============================================================================
%% Appendix: Solutions, Model Answers & Grading Rubrics
%% ============================================================================
\chapter{Solutions, Model Answers \& Grading Rubrics}
\label{ch:solutions}

This appendix collects model answers for every self-assessment exercise in
Chapters \ref{ch:introduction}--\ref{ch:synthesis} and grading rubrics for all
110 capstone projects. Three rules govern its use.

\begin{keyconceptbox}
\textbf{Attempt first, then compare.} Every answer here is
\emph{reference-grade}, not \emph{exclusive}: an alternative answer that is
correct, justified, and evidence-backed deserves full credit. Rubrics are
grading contracts --- they define what full marks require, not the only path
to them.
\end{keyconceptbox}

\begin{warningbox}
\textbf{Rubrics are operational, not cosmetic.} For senior and staff
capstones, correctness alone never earns full marks: the weights deliberately
reserve 30--45\% for tests, determinism, observability, and rollback
discipline. Grade against the criteria, not against your first impression of
the code.
\end{warningbox}

\paragraph{Numbering conventions.} Exercise numbering follows each chapter's
own scheme (e.g.\ E0.1, E14.2, E-3, (7.1), Exercise 7.1). A few chapters
number their exercises or capstones without a chapter prefix (``Capstone
1--5''); in this appendix every capstone is cited with a chapter prefix
(C6.1, C12.3, \ldots) for unambiguous reference, and a note flags each such
case. Chapter 17 numbers its exercises ``Exercise 7.1--7.8'' (a template
artifact); this appendix preserves that numbering. Where a chapter splits
questions into named lists (Exercises, Deep-Dive Questions, Concept Review,
Drills), the same split and order is kept here.

\begin{examplebox}
\textbf{How to grade with a rubric.} Score each criterion independently at
0\%, 50\%, or 100\% of its weight (no interpolation needed for classroom
use), multiply by the weight, and sum. A submission missing any
``full marks when\ldots'' condition drops that criterion to 50\% only if the
gap is demonstrably minor (e.g.\ one untested edge case); systemic gaps
(no test suite at all) score 0\% for that criterion.
\end{examplebox}

%% ============================================================================
\section{Chapter 0: Orientation, Failure Taxonomies \& Engineering Judgment}
\label{sec:sol-ch0}

\subsection{Self-Assessment Model Answers}

The readiness table (Table~\ref{tab:ch0-readiness}) is self-scored and has no
per-row key: award yourself a 2 only if you can perform the skill unprompted
(e.g.\ compute a percentile by hand on ten numbers, or explain the bytes
between \texttt{urlopen} and the response). Any honest 0 names your first
study topic; the target total $\geq 14$ stands.

\begin{enumerate}[label=\textbf{E0.\arabic*.}, leftmargin=3.2em]
  \item (a)~502 during a deploy: \emph{availability} failure in the
    operational layer; detect via per-status-code rate monitoring correlated
    with deployment markers. (b)~field type change: \emph{contract}
    (compatibility) failure; detect with schema-diff and contract tests in CI.
    (c)~wrong currency with 200: \emph{semantic} failure --- syntactically
    valid, wrong meaning; detect with invariant checks, canary metric
    comparison, and consumer-side validation. (d)~daily p99 spike:
    \emph{performance} (tail-latency) failure correlated with the batch
    window; detect with percentile alerts plus event correlation.
    (e)~ID-increment invoice access: \emph{security} failure (broken
    object-level authorization); detect with authorization regression tests
    and access-pattern anomaly audits.
  \item Sorted: 4, 5, 5, 6, 7, 9, 12, 18, 40, 220. Nearest-rank
    $p_{50}=7$\,ms (rank 5), $p_{90}=40$\,ms (rank 9), mean $=32.6$\,ms.
    Put $p_{99}$ (or $p_{90}$ with this little data) on the dashboard: the
    mean (32.6) is 4$\times$ the median and describes no real request ---
    pain concentrates in the tail. Always annotate sample size and offered
    load.
  \item Whitelist the query-key set (\texttt{in\_stock} only); any other key
    returns 400 with an \texttt{application/problem+json} body
    \cite{rfc9457} naming the offending parameter. Parse the value strictly
    (\texttt{true}/\texttt{false}), filter the catalog in the handler, keep
    type hints and stdlib-only dependencies, and extend the self-test to
    cover both branches.
  \item A 99.95\% monthly SLO allows $0.0005 \times 43{,}200 = 21.6$ minutes
    of downtime per 30-day month, i.e.\ $0.0005 \times 10^6 = 500$ failed
    responses per million requests. The per-request unit alerts better: it
    is traffic-normalized and computable on rolling windows, whereas
    ``minutes down'' requires an arbitrary definition of \emph{down}.
  \item Interleave the three payload sizes from a seeded PRNG (e.g.\
    \texttt{random.Random(42)}) and bucket samples per size before computing
    percentiles. Seeding makes the load mix identical across runs, so
    differences between runs reflect the system, not sampling noise;
    unseeded randomness makes results irreproducible and runs incomparable.
  \item The mini API scores roughly L1 (works, partial tests, no formal
    contract). Cheapest three upgrades: (1)~explicit timeouts on every
    outbound call; (2)~structured request logging with a correlation ID;
    (3)~a checked-in OpenAPI document with one generated contract test in
    CI. Together they reach L2 without architectural change.
  \item Timestamp at socket accept, handler entry, and response flush;
    aggregate 500 requests into per-stage p50/p99 columns and show the
    stages sum to the end-to-end budget of
    Equation~\eqref{eq:ch0-budget}. Tighten the stage with the largest
    share \emph{and} variance first --- typically the downstream-dependency
    wait, which is also where a timeout buys the most.
  \item Little's Law: $\lambda_{\max} = C/L = 128/30 \approx 4.3$ rps with
    30-second stalls; with a 2-second timeout, $\lambda_{\max} = 128/2 = 64$
    rps. A single configuration change buys $\sim$15$\times$ capacity ---
    the cheapest scaling decision in the book.
  \item \textbf{S}: no authentication --- ignored. \textbf{T}: no integrity
    controls beyond transport --- partially ignored. \textbf{R}: no audit
    log --- ignored. \textbf{I}: error bodies may leak internals ---
    partially controlled. \textbf{D}: no rate limiting, no timeouts on
    inbound work --- ignored. \textbf{E}: input validation on paths and
    bodies --- partially controlled. The code today controls only fragments
    of I and E; the exercise's point is how much a ``working'' API leaves
    open.
  \item Fail the build on: missing/invalid OpenAPI document, contract-test
    failure, breaking-change diff without a major bump, missing error-schema
    conformance (RFC~9457), absent timeouts in client code. Warn on:
    coverage regression, missing runbook link, undocumented error codes.
    Override: platform-team approval with a time-boxed, logged exception
    (compare the addendum's CI quality-gate template --- the gate fails
    closed, the exception process is the only escape valve).
\end{enumerate}

\subsubsection*{Deep-Dive Questions}

\begin{enumerate}[label=\textbf{D0.\arabic*.}, leftmargin=3.2em]
  \item \emph{For:} retryability is failure semantics --- part of the
    contract by Definition~\ref{def:api}; a fleet that starts retrying POST X
    multiplies write traffic, and if the endpoint is not idempotency-safe the
    provider has silently accepted duplicate-mutation risk.
    \emph{Against:} it is additive guidance; no schema changes and clients
    may ignore it. Verdict: treat it as breaking \emph{unless} the endpoint
    is provably idempotent (keys or natural idempotence) --- then it is a
    minor contract expansion.
  \item Example: a platform with rigorous design governance (catalog,
    linting, review boards) whose services run single-region with no
    timeouts, bulkheads, or breakers. The risk is compounded: governance
    certifies the APIs, so consumers trust them, so the first real
    infrastructure incident cascades through a fleet of ``compliant'' but
    fragile services --- and the governance program absorbs the blame.
  \item A retry storm violates stationarity and independence: arrivals are
    no longer exogenous (retries are generated by failures), and $L$
    inflates as queues build. Applying $C=\lambda L$ with the original
    offered rate then \emph{understates} true concurrency --- the law
    reassures you precisely when the system has left the regime where it
    applies.
  \item Count shed load as failure at the edge: SLI $=$
    (requests served correctly within the latency bound) $/$
    (all offered requests), measured \emph{before} load shedding, with a
    separate alert on the shed ratio itself. Fast 503s then burn budget
    instead of gaming the percentile.
  \item Offline labs cannot teach: real TLS/certificate failure modes, CDN
    and proxy behavior, provider-specific rate-limit shapes, DNS and network
    partitions, genuine latency distributions, and identity-provider edge
    cases. Compensate with sandbox accounts, recorded-traffic replay, chaos
    shims that mimic provider faults, provider status-page archaeology, and
    small-blast-radius canaries before depending on any real provider.
\end{enumerate}

\subsection{Capstone Grading Rubrics}

\subsubsection*{Capstone 0.1 --- Contract Archaeology [junior]}
\begin{center}\small
\begin{tabularx}{\textwidth}{l c X}
\toprule
\textbf{Criterion} & \textbf{Weight} & \textbf{Full marks when\ldots} \\
\midrule
Contract coverage & 35\% & every path, method, status, and error body of the mini API appears in the written contract. \\
Failure semantics & 25\% & each operation documents its error codes, retryability, and idempotence posture. \\
Methodology evidence & 20\% & probe transcripts/logs are included and reproducible by re-running the probe script. \\
Reflection & 20\% & the writeup names at least two things black-box probing cannot discover (rate limits, auth rules, SLAs, deprecation policy). \\
\bottomrule
\end{tabularx}
\end{center}
\textbf{Reference approach.} Build a probe matrix: enumerate routes from
hints and 404 behavior, exercise happy path plus malformed, unknown-ID, and
wrong-method cases per route, capture raw request/response pairs, then
tabulate the contract from evidence only. The stretch (a hand-written
OpenAPI~3.1 document validated offline \cite{openapi31}) converts the table
into a machine-checkable artifact.

\subsubsection*{Capstone 0.2 --- Honest Benchmarks [mid]}
\begin{center}\small
\begin{tabularx}{\textwidth}{l c X}
\toprule
\textbf{Criterion} & \textbf{Weight} & \textbf{Full marks when\ldots} \\
\midrule
Experimental design & 30\% & 3 load shapes $\times$ 3 repetitions, seeded, with warmup discarded and raw CSVs committed. \\
Statistical analysis & 30\% & per-config p50/p95/p99 show where the tail detaches, attributed to a budget stage. \\
Reproducibility & 20\% & one command regenerates every table and figure bit-for-bit. \\
Interpretation & 20\% & the writeup explains quantitatively why the mean understated user pain. \\
\bottomrule
\end{tabularx}
\end{center}
\textbf{Reference approach.} Extend the chapter benchmark harness with
seeded load shapes (steady, bursty, ramp); write one CSV row per request;
analyze offline with a small script that computes percentiles per
configuration and reconciles stage timings with the latency-budget equation.
The stretch fits $L = a + b\cdot\text{bytes}$ from a fourth variable.

\subsubsection*{Capstone 0.3 --- Failure Taxonomy Field Guide [mid]}
\begin{center}\small
\begin{tabularx}{\textwidth}{l c X}
\toprule
\textbf{Criterion} & \textbf{Weight} & \textbf{Full marks when\ldots} \\
\midrule
Taxonomy coverage & 30\% & five runnable demos, one per failure category, each offline and deterministic. \\
Detection evidence & 30\% & each demo prints the symptom \emph{and} the monitoring signal/alert that would catch it. \\
Playbook quality & 20\% & on-call entries follow the incident format: symptom, cause class, detection, first action. \\
Determinism \& hygiene & 20\% & seeded randomness only; scripts exit non-zero on unexpected outcomes. \\
\bottomrule
\end{tabularx}
\end{center}
\textbf{Reference approach.} One script per category (crash, timeout,
contract drift, semantic error, tail-latency), each driving the mini API
against a seeded chaos stub; every run ends with a symptom $\to$ category
$\to$ detection mapping line. The stretch packages the flaky dependency as a
reusable seeded stub for later chapters' drills.

\subsubsection*{Capstone 0.4 --- Budgets and Breakers for the Mini API [senior]}
\begin{center}\small
\begin{tabularx}{\textwidth}{l c X}
\toprule
\textbf{Criterion} & \textbf{Weight} & \textbf{Full marks when\ldots} \\
\midrule
Protection correctness & 30\% & timeouts, bounded jittered retries, and a three-state circuit breaker each have focused tests \cite{nygard2018releaseit}. \\
Controlled experiment & 30\% & the \emph{same} scripted fault runs with and without protections; p99 and error rate both improve when protected. \\
SLO arithmetic & 20\% & error budget computed from raw samples, not assumed. \\
Operational writeup & 20\% & post-incident report in war-story form: timeline, trigger, mechanism, fix, follow-ups. \\
\bottomrule
\end{tabularx}
\end{center}
\textbf{Reference approach.} Wrap the dependency call in timeout
$\to$ retry-budget $\to$ breaker layers; inject a deterministic stall fault;
graph unprotected vs protected latency distributions; compute budget burn
from the sample CSVs. The stretch adds a bulkhead and shows an unrelated
endpoint surviving the fault.

\subsubsection*{Capstone 0.5 --- API Platform Governance Starter [staff]}
\begin{center}\small
\begin{tabularx}{\textwidth}{l c X}
\toprule
\textbf{Criterion} & \textbf{Weight} & \textbf{Full marks when\ldots} \\
\midrule
Catalog completeness & 25\% & every lab API registered with owner, tier, lifecycle state, and contract link. \\
Automated gates & 30\% & offline linter exits non-zero on violations; gate runs in CI with a documented exception workflow. \\
Policy measurability & 20\% & deprecation policy defines sunset in observable terms (usage at zero for $n$ weeks). \\
Validation \& audit & 25\% & automated maturity report agrees with a manual audit on at least one API; discrepancies analyzed. \\
\bottomrule
\end{tabularx}
\end{center}
\textbf{Reference approach.} Define a YAML/JSON catalog schema; write a
stdlib linter over the registered OpenAPI documents; wire it into CI with a
time-boxed waiver mechanism; write the deprecation policy against telemetry
keys; produce the quarterly maturity report by script. The stretch generates
a static HTML portal from the same catalog --- governance that ships value
gets adopted (Chapter~\ref{ch:lifecycle}).

%% ============================================================================
\section{Chapter 1: HTTP on the Wire}
\label{sec:sol-ch1}

\subsection{Self-Assessment Model Answers}

\begin{enumerate}[label=\textbf{E1.\arabic*}, leftmargin=3.2em]
  \item The server reads exactly \texttt{Content-Length} bytes, leaving one
    stray body byte in the TCP receive buffer; the next response parse on the
    keep-alive connection begins at that byte, so the status line is garbage
    or the framing splits mid-header. This is the desynchronization class
    behind request smuggling: \texttt{Content-Length} is the only framing
    authority, and writer/reader disagreement corrupts every subsequent
    message on the connection \cite{rfc9112}.
  \item The client's \texttt{pending} buffer must accumulate bytes until a
    complete chunk-size line, then a full chunk plus CRLF, have arrived ---
    TCP delivers byte streams, not chunks. Lowering \texttt{MAX\_CHUNK}
    below an advertised chunk size must raise the defensive
    \texttt{ValueError} \emph{before} allocating or reading, bounding memory
    against hostile or buggy peers.
  \item \texttt{POST /robots}: 201 + \texttt{Location} on success; 400
    (malformed body), 422 (semantically invalid), 409 (duplicate serial),
    401/403. \texttt{GET /robots/\{id\}}: 200; 304 with
    \texttt{If-None-Match}; 404. \texttt{PUT /robots/\{id\}/firmware}:
    200/204 (or 202 if the update is deferred); 400/422; 404; 412 on a
    failed \texttt{If-Match} precondition. \texttt{DELETE
    /robots/\{id\}}: 204; 404 on repeat (still idempotent); 409 if a
    referenced resource blocks deletion \cite{rfc9110}.
  \item Default \texttt{SETTINGS\_MAX\_FRAME\_SIZE} is $2^{14}=16{,}384$
    octets. A 70{,}000-octet header block emits $\lceil 70000/16384\rceil =
    5$ frames: one HEADERS (without \texttt{END\_HEADERS}) followed by four
    CONTINUATION frames, the last carrying \texttt{END\_HEADERS} (and
    \texttt{END\_STREAM} if the request has no body)
    \cite{rfc9113}.
  \item HTTP/1.1 resends the full header block every request: at roughly
    400~bytes each, 100 requests cost $\approx$40\,KB. With HPACK, request
    1 encodes literally and populates the dynamic table; requests 2--100
    encode as indexed references plus one literal \texttt{x-trace-id}
    ($\sim$20 bytes), totaling $\approx$2--3\,KB --- a $>$90\% reduction,
    assuming a 4{,}096-octet dynamic table that never evicts the reusable
    entries \cite{rfc7541}.
  \item Only DATA frames consume flow-control credit, so HEADERS, SETTINGS,
    PING, RST\_STREAM, PRIORITY, WINDOW\_UPDATE, and GOAWAY still flow ---
    legal and deliberate, so the connection stays diagnosable and
    recoverable. A server that never sends \texttt{WINDOW\_UPDATE} leaves
    the client with a permanently zero window: requests are accepted,
    responses never start, and users see hangs ending in client timeouts.
  \item Cleartext: \texttt{ClientHello}, \texttt{ServerHello}. Encrypted
    under handshake keys: \texttt{EncryptedExtensions},
    \texttt{Certificate}, \texttt{CertificateVerify}, \texttt{Finished}
    \cite{rfc8446}. On resumption, 0-RTT early data can be sent
    \emph{immediately after the ClientHello} (encrypted under PSK-derived
    early traffic keys), before any server message arrives.
  \item CL.TE layout: the front-end trusts \texttt{Content-Length: 6} and
    forwards six body bytes; the back-end trusts
    \texttt{Transfer-Encoding: chunked}, sees a terminating zero chunk, and
    treats the leftover bytes (the smuggled \texttt{G\ldots} prefix) as the
    start of the \emph{next} request. Mitigations: (1)~the front-end must
    reject or normalize any request carrying both headers; (2)~use HTTP/2
    end-to-end, whose binary length-prefixed framing removes the ambiguity
    --- and never downgrade h2$\to$h1 without strict normalization.
  \item On a lossy link (1\% loss, 50\,ms RTT), one dropped segment stalls
    \emph{all ten} HTTP/2 streams behind TCP's ordered delivery while
    retransmission takes a full RTT; six HTTP/1.1 connections lose only the
    unlucky connection's request. h2 loses exactly when the multiplexing
    gain is smaller than the HOL stall probability. HTTP/3 differs because
    QUIC recovers loss per stream: a lost packet stalls only its own stream
    \cite{rfc9000, rfc9114}.
  \item 0-RTT data is encrypted under resumption keys with no server
    round-trip, so there is no freshness exchange --- a network attacker can
    capture and replay it verbatim. Replayed \texttt{POST /payments} means
    duplicate charges. Defences: restrict early data to safe/idempotent
    methods via an allowlist, use single-use session tickets, keep
    anti-replay windows tight, and require application-level idempotency
    keys for anything money-moving (Chapter~\ref{ch:idempotency}).
\end{enumerate}

\subsection{Capstone Grading Rubrics}

\subsubsection*{Capstone 1.1 --- From-Scratch HTTP/1.1 Server [junior]}
\begin{center}\small
\begin{tabularx}{\textwidth}{l c X}
\toprule
\textbf{Criterion} & \textbf{Weight} & \textbf{Full marks when\ldots} \\
\midrule
Wire correctness & 35\% & \texttt{curl -v} shows byte-correct status line, CRLF headers, and exact \texttt{Content-Length} framing. \\
Keep-alive & 20\% & ten sequential requests on one connection all succeed. \\
Error semantics & 20\% & malformed request-line $\to$ 400, unknown path $\to$ 404, POST $\to$ 405, each byte-correct. \\
HEAD correctness & 15\% & HEAD returns GET-identical headers with zero body bytes. \\
Code quality & 10\% & typed, stdlib-only, readable parser state machine. \\
\bottomrule
\end{tabularx}
\end{center}
\textbf{Reference approach.} A \texttt{socketserver} request handler loops:
read until CRLFCRLF, parse request line and headers, dispatch
\texttt{/parts} and \texttt{/parts/\{id\}} from a route table, serialize the
JSON body once to compute \texttt{Content-Length}, honor
\texttt{Connection: close} by exiting the loop. The stretch adds
\texttt{ETag}/\texttt{If-None-Match} $\to$ 304.

\subsubsection*{Capstone 1.2 --- Chunked Streaming Proxy [mid]}
\begin{center}\small
\begin{tabularx}{\textwidth}{l c X}
\toprule
\textbf{Criterion} & \textbf{Weight} & \textbf{Full marks when\ldots} \\
\midrule
Relay fidelity & 30\% & decoded body is byte-identical to a direct fetch, re-chunked in 256-octet frames. \\
Defensive parsing & 25\% & oversized or malformed chunks rejected with 502 and a log line naming the violation. \\
Header hygiene & 20\% & hop-by-hop headers stripped per RFC~9110 \S7.6.1 \cite{rfc9110}. \\
Offline determinism & 15\% & deterministic test script runs the whole matrix with no network. \\
Observability & 10\% & per-request byte counts and timings via \texttt{logging}. \\
\bottomrule
\end{tabularx}
\end{center}
\textbf{Reference approach.} Threaded TCP proxy: forward the request head
verbatim, then run a bounded chunk-decoding state machine over origin bytes,
re-emitting fixed-size chunks to the client. Stretch: trailer support and
bounded buffers for backpressure.

\subsubsection*{Capstone 1.3 --- HPACK-Lite Header Encoder [mid]}
\begin{center}\small
\begin{tabularx}{\textwidth}{l c X}
\toprule
\textbf{Criterion} & \textbf{Weight} & \textbf{Full marks when\ldots} \\
\midrule
Integer codec & 25\% & prefix integer coding round-trips every corpus value including prefix overflow \cite{rfc7541}. \\
Indexing \& dynamic table & 30\% & a repeated header block encodes as a single indexed reference; table never exceeds its size bound (eviction tested). \\
Compression ratio & 25\% & $\geq$60\% byte reduction on the repeat-heavy corpus, measured in the test. \\
Test quality & 20\% & golden hex vectors plus property-style round-trips, fully offline. \\
\bottomrule
\end{tabularx}
\end{center}
\textbf{Reference approach.} Implement the integer prefix primitive first
with golden vectors, then a static-table dictionary, then a dynamic table
with FIFO eviction; the encoder prefers indexed $\to$
literal-with-incremental-indexing. Stretch: static Huffman and a decoder
that round-trips the encoder.

\subsubsection*{Capstone 1.4 --- Protocol Fingerprinting \& Frame Dissection [senior]}
\begin{center}\small
\begin{tabularx}{\textwidth}{l c X}
\toprule
\textbf{Criterion} & \textbf{Weight} & \textbf{Full marks when\ldots} \\
\midrule
Classification accuracy & 30\% & every fixture classified (h1 text, h2 preface, TLS records) with zero false positives. \\
Anomaly detection & 25\% & planted CL.TE desync and truncated frame both flagged with byte offsets. \\
Output contract & 20\% & structured report plus \texttt{--json} mode; deterministic across runs. \\
Engineering & 15\% & fixture generator and tests committed; CLI usable without reading the source. \\
Depth of dissection & 10\% & frame-level decode (type, flags, stream id) rather than heuristic string matching. \\
\bottomrule
\end{tabularx}
\end{center}
\textbf{Reference approach.} Layered recognizers: TLS record-header check
first, then the h2 client preface, then an HTTP/1.1 token grammar; each
recognized stream feeds a frame dissector emitting structured records.
Stretch: HPACK state reconstruction across HEADERS/CONTINUATION boundaries.

\subsubsection*{Capstone 1.5 --- HTTP Version Latency Profiler [staff]}
\begin{center}\small
\begin{tabularx}{\textwidth}{l c X}
\toprule
\textbf{Criterion} & \textbf{Weight} & \textbf{Full marks when\ldots} \\
\midrule
Harness rigor & 30\% & one-command offline run, seeded randomness, warmup excluded, raw samples kept. \\
Measurement validity & 25\% & warm-vs-cold keep-alive benefit quantified per payload size; delay shim demonstrates HOL growth under injected latency. \\
Reporting & 20\% & output table conforms to the addendum's benchmark template with units and $n$. \\
Extensibility & 15\% & pluggable transport interface so a real h2 client drops in without harness changes. \\
Operational notes & 10\% & README states noise controls, known confounders, and how results should drive decisions. \\
\bottomrule
\end{tabularx}
\end{center}
\textbf{Reference approach.} Insert a delay/loss shim between client and
server sockets; sweep payload sizes $\times$ connection reuse; summarize
with percentiles, not means. The staff signal is the analysis: results must
explain \emph{which} budget stage each configuration moves and when the
conclusion would flip on a lossy network.

%% ============================================================================
\section{Chapter 2: Consuming APIs in Python}
\label{sec:sol-ch2}

\subsection{Self-Assessment Model Answers}

\emph{Numbering note: the chapter uses a plain 1--6 list; numbers below
match it.}

\begin{enumerate}[leftmargin=2.2em]
  \item Each un-annotated call leaves at least the connect or read dimension
    unbounded; a silent peer then parks a worker forever, and at scale the
    pool exhausts and the outage cascades --- the Chapter~\ref{ch:http_wire}
    war-story mechanism. For every hit, name which of connect/read/write/pool
    is unbounded and the failure it permits.
  \item Deterministic exponential waits: $0.5, 1, 2, 4$~s $\to$ 7.5~s total.
    Full jitter draws $U[0, \text{cap}_i]$, halving the expectation:
    $E \approx 3.75$~s. Worst case equals the deterministic total (7.5~s);
    the $c=8$~s cap never binds because no individual wait exceeds 4~s.
  \item \texttt{GET /parts/NR-1001}: safe and idempotent --- retry blindly.
    \texttt{POST /orders}: unsafe --- a read timeout is ambiguous (the order
    may exist); retry only with an idempotency key. \texttt{PUT
    \ldots/status}: idempotent --- safe. \texttt{POST /payments} with
    \texttt{Idempotency-Key}: safe \emph{iff} the server honors the key;
    verify that contract before retrying.
  \item Worst case is $32 \times 2 = 64$ simultaneous upstream calls, so
    \texttt{max\_connections=64} (headroom, not average) and
    \texttt{max\_keepalive\_connections} between 32 and 64 so warm
    connections survive between bursts without unbounded idle growth.
  \item \texttt{ApiAuthError} is raised for 401/403 inside
    \texttt{\_raise\_for\_status} \emph{before} retry classification, so the
    retry policy maps it to non-retryable. The MockTransport test returns
    401, asserts the typed exception, and asserts the transport saw exactly
    one request --- auth failures must never be retried.
  \item Loop pages with \texttt{offset += limit}; terminate when a page
    returns fewer than \texttt{limit} items (or is empty); raise on
    \texttt{max\_pages} as a tripwire against infinite loops from a buggy
    server. MockTransport serves scripted page sequences to cover both
    termination branches.
\end{enumerate}

\subsubsection*{Deep-Dive Questions}

\begin{enumerate}[leftmargin=2.2em]
  \item The pool timeout bounds waiting for a \emph{free connection slot};
    the read timeout bounds waiting for \emph{response bytes} after the
    request is sent. \texttt{PoolTimeout} therefore signals pool
    undersizing or upstreams stalling while holding connections --- a
    topology/configuration fact that \texttt{ReadTimeout} can never reveal.
  \item For the individual caller, halved expected wait is a latency
    \emph{saving}. The fleet-level win holds regardless of the mean:
    jitter decorrelates retry instants across clients, preventing the
    synchronized retry herd that re-collapses a recovering service
    \cite{nygard2018releaseit}.
  \item \texttt{Retry-After} is authoritative for the server's state, but
    your own latency budget is authoritative for your caller. On
    irreconcilable conflict (120~s $>$ 20~s budget): stop retrying, fail
    fast with a typed error that preserves the server's hint, and let the
    caller decide --- never silently exceed your budget.
  \item \texttt{requests} is built on urllib3's synchronous HTTP/1.1
    connection core; multiplexed framing would require rewriting that core.
    \texttt{httpx} separated the transport behind an interface, so HTTP/2
    became a pluggable option rather than a fork
    \cite{httpxdocs, requestsdocs}.
  \item Reader fetches page 1 (rows 1--5); a writer inserts one row at the
    front; page 2 at offset 5 re-reads row 5 (duplicate) --- and the mirror
    image, deleting a read row, skips an unread one. A cursor pins the walk
    to key values, so inserts or deletes \emph{before} the cursor cannot
    shift positions (Chapter~\ref{ch:pagination}).
  \item MockTransport legitimately excluded the socket/TLS layer. Minimum
    offline addition: a test that builds the real TLS context and completes
    a handshake against a local self-signed \texttt{ssl} socket server,
    asserting hostname verification behavior --- no network required.
\end{enumerate}

\subsection{Capstone Grading Rubrics}

\subsubsection*{Capstone 2.1 --- Robust CLI Downloader [junior]}
\begin{center}\small
\begin{tabularx}{\textwidth}{l c X}
\toprule
\textbf{Criterion} & \textbf{Weight} & \textbf{Full marks when\ldots} \\
\midrule
Timeout discipline & 25\% & a test asserts no code path issues a request without an explicit timeout. \\
Artifact integrity & 25\% & a killed download leaves no partial file (write to \texttt{.tmp}, atomic rename). \\
Fault-matrix tests & 30\% & 200, 404, 500-then-200, and malformed body covered via MockTransport with the stub stopped. \\
Log hygiene & 20\% & log lines carry method, path, status, latency --- and no headers or tokens. \\
\bottomrule
\end{tabularx}
\end{center}
\textbf{Reference approach.} Thin CLI over an \texttt{httpx} client with a
four-dimension timeout; stream the body to a temp file, fsync, rename;
map status classes to typed errors. Stretch: \texttt{Range} resume and a
\texttt{--dry-run} mode.

\subsubsection*{Capstone 2.2 --- Rate-Limit-Aware Catalog Crawler [mid]}
\begin{center}\small
\begin{tabularx}{\textwidth}{l c X}
\toprule
\textbf{Criterion} & \textbf{Weight} & \textbf{Full marks when\ldots} \\
\midrule
Rate conformance & 30\% & full crawl completes with zero steady-state 429s at $\leq$5 req/s, evidenced by the timestamp CSV. \\
Retry-After correctness & 25\% & a forced-429 MockTransport test proves \texttt{Retry-After} overrides computed jitter. \\
Determinism & 25\% & seeded RNG and injected sleeper; no wall-clock assertions anywhere. \\
Cursor correctness & 20\% & cursor-paginated walk visits every catalog item exactly once. \\
\bottomrule
\end{tabularx}
\end{center}
\textbf{Reference approach.} Client-side token bucket paces requests below
the published limit; on 429, sleep exactly \texttt{Retry-After} (injected
clock); walk cursors until \texttt{next} is absent. Stretch: AIMD bucket
adaptation and a process-wide bucket shared by parallel crawlers
(Chapter~\ref{ch:rate_limiting}).

\subsubsection*{Capstone 2.3 --- Typed SDK Wrapper with Full Test Matrix [mid]}
\begin{center}\small
\begin{tabularx}{\textwidth}{l c X}
\toprule
\textbf{Criterion} & \textbf{Weight} & \textbf{Full marks when\ldots} \\
\midrule
Fault matrix breadth & 30\% & $\geq$15 scenarios spanning every taxonomy branch (transport, 4xx classes, 5xx, malformed bodies). \\
Retry/idempotency policy & 25\% & contract errors never retried; transient errors retried within budget; mutations gated on idempotency keys. \\
Type safety & 20\% & public API fully typed; \texttt{mypy} clean; pydantic models validate at the boundary \cite{pydanticdocs}. \\
Docs \& packaging & 25\% & installable local package; PowerShell-only README covering install, usage, and error semantics. \\
\bottomrule
\end{tabularx}
\end{center}
\textbf{Reference approach.} Layer as models $\to$ transport policy $\to$
resource methods $\to$ pagination iterators; express the retry policy as a
table mapping error class to action. Stretch: generate models from JSON
Schema at build time and add an async facade over \texttt{httpx.AsyncClient}.

\subsubsection*{Capstone 2.4 --- Chaos-Client Harness [senior]}
\begin{center}\small
\begin{tabularx}{\textwidth}{l c X}
\toprule
\textbf{Criterion} & \textbf{Weight} & \textbf{Full marks when\ldots} \\
\midrule
Fault coverage & 30\% & $\geq$8 fault classes, each with an executable expectation (not just an injection). \\
Reproducibility & 25\% & every run replays from its seed with zero wall-clock dependence. \\
Assertion rigor & 25\% & at least one fault provably exhausts the retry budget and fails fast with the correct typed error. \\
Safety guards & 20\% & decompression-bomb test proves the size guard engages. \\
\bottomrule
\end{tabularx}
\end{center}
\textbf{Reference approach.} A chaos transport wraps MockTransport with a
seeded fault schedule (latency spikes, resets, truncation, 5xx bursts,
gzip bombs); the resilient client runs unmodified on top. Stretch: fake-clock
latency injection and chaos schedules replayed from synthetic incident
timelines.

\subsubsection*{Capstone 2.5 --- Async Concurrent Fetcher [staff]}
\begin{center}\small
\begin{tabularx}{\textwidth}{l c X}
\toprule
\textbf{Criterion} & \textbf{Weight} & \textbf{Full marks when\ldots} \\
\midrule
Correctness under load & 25\% & 10{,}000-item queue drains with zero unhandled exceptions and zero pool timeouts. \\
Rate conformance & 25\% & sustained $\leq$50 req/s with $<$0.1\% 429s after warm-up. \\
Cancellation hygiene & 20\% & cancellation test shows no orphan tasks and no partial artifacts. \\
Determinism & 15\% & MockTransport plus injected async sleeps; suite is reproducible. \\
Design memo & 15\% & every knob (pool, concurrency, backoff, budget) quantified with its trade-off. \\
\bottomrule
\end{tabularx}
\end{center}
\textbf{Reference approach.} An \texttt{asyncio} queue feeding a bounded
worker pool; semaphore-limited concurrency; one token bucket for pacing;
structured cancellation via task groups so \texttt{close()} always runs.
Stretch: AIMD adaptive concurrency and per-host fairness buckets.

%% ============================================================================
\section{Chapter 3: The REST Architectural Style}
\label{sec:sol-ch3}

\subsection{Self-Assessment Model Answers}

\emph{Numbering note: the chapter numbers its exercises (7.1)--(7.8); this
appendix preserves that scheme.}

\begin{enumerate}[label=\textbf{(7.\arabic*)}, leftmargin=3.2em]
  \item A lint-clean route table uses plural-noun resources
    (\texttt{/parts}, \texttt{/parts/\{id\}}, \texttt{/orders/\{id\}/lines})
    with methods mapped to semantics (GET read, POST create, PUT replace,
    PATCH partial update, DELETE remove), correct status codes (200/201/204/
    404/409/412), and an explicit cache policy per GET route
    \cite{fielding2000rest}. Controller-style routes (\texttt{/getPart})
    fail the lint rules.
  \item Session cookie: \emph{not} stateless if the server must look up
    server-side session state --- the request alone is not self-contained.
    Bearer token: stateless --- the request carries all context. OAuth token
    validated by introspection: the request itself is self-contained, but
    validation reintroduces a hidden server-to-server dependency; strictly
    stateless per the definition, operationally a coupling to watch.
  \item URI \texttt{/promotions/today} (or date-keyed
    \texttt{/promotions/2026-09-03}); freshness via \texttt{Cache-Control:
    max-age} bounded by time-to-midnight, plus \texttt{ETag} and
    \texttt{Last-Modified} validators. Eight hours stale in a shared cache
    means customers see yesterday's prices --- direct revenue and trust
    damage, caused by a caching-policy bug, not a data bug.
  \item HEAD derives from the GET handler minus the body (strengthening the
    uniform interface and cacheability --- validators become usable);
    \texttt{Idempotency-Key} on \texttt{POST /parts} strengthens the
    stateless uniform contract's reliability by making the one
    non-idempotent method safely retryable (Chapter~\ref{ch:idempotency}).
  \item The JSON~Patch evaluator applies \texttt{add}/\texttt{remove}/
    \texttt{replace}/\texttt{test} atomically (all-or-nothing, with
    \texttt{test} as precondition). Media type
    \texttt{application/json-patch+json} selects it;
    \texttt{application/merge-patch+json} keeps the merge evaluator ---
    content negotiation picks the semantics.
  \item Scenario: pushing firmware to 50{,}000 robots one resource-PUT at a
    time vs one bespoke RPC broadcast. The uniform interface costs
    50{,}000 round trips where one would do; what the inefficiency purchases
    is evolvability --- every intermediary, cache, tool, and new client
    understands the uniform semantics without new code
    \cite{fielding2000rest}.
  \item \emph{For} DELETE on \texttt{/warehouses/\{wid\}/stock/\{sku\}}:
    uniform, automation-friendly, and honest --- 410 can signal deliberate
    retirement. \emph{Against}: it conflates ``zero stock'' (a state of an
    existing record) with ``no record'' (absence), and invites accidental
    data loss. 404 communicates ``currently absent, may return''; 410
    communicates ``gone by policy'' \cite{rfc9110}.
  \item JSON:API decides for you: document shape, URL conventions, error
    objects, sparse fieldsets, and pagination parameters --- entire
    Chapter~\ref{ch:rest_style} bikesheds vanish. It forbids custom
    envelopes and bespoke error formats; its opinion costs flexibility
    exactly where domain-specific representations or non-CRUD operations
    matter.
\end{enumerate}

\subsection{Capstone Grading Rubrics}

\subsubsection*{Capstone 3.1 --- Model the Library Domain [junior]}
\begin{center}\small
\begin{tabularx}{\textwidth}{l c X}
\toprule
\textbf{Criterion} & \textbf{Weight} & \textbf{Full marks when\ldots} \\
\midrule
Resource purity & 35\% & zero lint findings (R001--R005); no controller archetype without a written reification-failure analysis. \\
Model completeness & 25\% & books, copies, members, loans, holds, and fines all modeled with lifecycle states. \\
URI discipline & 20\% & nesting depth $\leq$2; every collection/member relationship explicit. \\
Cache policy & 20\% & every GET route has a stated validator and freshness policy. \\
\bottomrule
\end{tabularx}
\end{center}
\textbf{Reference approach.} List nouns, not verbs; give loans and holds
their own resources (they have lifecycles); express actions as state
transitions on resources. Stretch: sparse fieldsets and a Merge Patch
document renewing a loan.

\subsubsection*{Capstone 3.2 --- Run and Extend the Stdlib Catalog [junior]}
\begin{center}\small
\begin{tabularx}{\textwidth}{l c X}
\toprule
\textbf{Criterion} & \textbf{Weight} & \textbf{Full marks when\ldots} \\
\midrule
Constraint demonstrations & 40\% & seven curl transcripts captured: 200+ETag, 304, 201+Location, 409, 412, 406, CSV via \texttt{Accept}. \\
Extension correctness & 25\% & the ``batteries'' category works end-to-end and the selftest still passes. \\
Analysis & 25\% & the memo correctly identifies statelessness as the first constraint server-side sessions would break. \\
Hygiene & 10\% & transcripts reproducible; no code changes beyond the extension. \\
\bottomrule
\end{tabularx}
\end{center}
\textbf{Reference approach.} Drive the listing with scripted curl; capture
raw exchanges; add the category through the public surface only. Stretch:
HEAD-only verification and Sunset/Deprecation headers on a retired part.

\subsubsection*{Capstone 3.3 --- Productionize the URI Linter [mid]}
\begin{center}\small
\begin{tabularx}{\textwidth}{l c X}
\toprule
\textbf{Criterion} & \textbf{Weight} & \textbf{Full marks when\ldots} \\
\midrule
Exit-code contract & 20\% & documented 0/1/2 semantics (clean / violations / tool error) honored exactly. \\
Waiver mechanism & 25\% & waivers suppress only the cited rule+path and expire; expired waivers fail the build. \\
Rule test coverage & 30\% & every rule has positive and negative deterministic tests. \\
Machine-readable output & 25\% & JSON report validates against a committed schema. \\
\bottomrule
\end{tabularx}
\end{center}
\textbf{Reference approach.} Keep the rule engine pure (routes in, findings
out); layer CLI, waivers file, and JSON serializer around it. Stretch:
OpenAPI~3.1 as an input adapter and a route-table diff classifying deltas
as breaking vs additive (Chapter~\ref{ch:versioning}).

\subsubsection*{Capstone 3.4 --- Port an RPC API to Resource Orientation [senior]}
\begin{center}\small
\begin{tabularx}{\textwidth}{l c X}
\toprule
\textbf{Criterion} & \textbf{Weight} & \textbf{Full marks when\ldots} \\
\midrule
Migration completeness & 30\% & all 12 legacy endpoints mapped to resources, or each exception has a written reification-failure analysis. \\
Compatibility & 30\% & shims preserve behavior for all 12 legacy calls, proven by before/after transcripts. \\
Contract quality & 20\% & design-first OpenAPI lints clean. \\
Migration plan & 20\% & dates, owners, and kill criteria keyed to a synthetic access log; no ``someday'' items. \\
\bottomrule
\end{tabularx}
\end{center}
\textbf{Reference approach.} Design the resource model first, generate the
contract, then implement shims that translate legacy RPC calls onto the new
resources so consumers migrate on their own schedule
\cite{richardson2013restful}. Stretch: hypermedia transition links and a
payload/round-trip delta quantification.

\subsubsection*{Capstone 3.5 --- Design-Review Report on a Public API [staff]}
\begin{center}\small
\begin{tabularx}{\textwidth}{l c X}
\toprule
\textbf{Criterion} & \textbf{Weight} & \textbf{Full marks when\ldots} \\
\midrule
Evidence discipline & 30\% & every claim cites observable evidence (docs, captured exchanges); no speculation presented as fact. \\
Balance & 20\% & at least one finding credits a sophisticated design choice --- no pure hatchet job. \\
Remediation quality & 30\% & the backlog obeys the no-rename-and-restructure-in-one-step doctrine with sequencing. \\
Communication & 20\% & $\leq$8 pages, executive-readable, constraint-by-constraint verdicts. \\
\bottomrule
\end{tabularx}
\end{center}
\textbf{Reference approach.} Score the API against each REST constraint
with captured evidence; classify findings by evolvability impact; propose a
remediation sequence that never breaks existing consumers. Stretch: repeat
the review five years apart from archived docs to test evolvability claims.

%% ============================================================================
\section{Chapter 4: Serialization \& Contracts}
\label{sec:sol-ch4}

\subsection{Self-Assessment Model Answers}

\begin{enumerate}[label=\textbf{E4.\arabic*}, leftmargin=2.2em]
  \item Varints: 1 $\to$ \texttt{01}; 127 $\to$ \texttt{7F}; 128 $\to$
    \texttt{80 01} (continuation bit set on the low seven bits first);
    16{,}384 $\to$ \texttt{80 80 01}. Verify each against the codec's
    \texttt{encode\_varint} \cite{protobufdocs}.
  \item $\mathrm{zigzag}(-64) = 127$ (negative $n$ maps to $-2n-1$);
    $\mathrm{zigzag}^{-1}(127) = -64$. The point: small-magnitude negatives
    get small varints, which plain two's-complement cannot deliver.
  \item \texttt{double} is wire type 1 (fixed 64-bit), \texttt{int64} is
    wire type 0 (varint). Old consumers see tag 4 with an unexpected wire
    type: they cannot interpret the bytes as their known field and skip it
    --- silent data loss. This is why the evolution rules forbid changing a
    field's wire type; deprecate tag 4 and add a new tag instead
    \cite{protobufdocs}.
  \item Add \texttt{if: \{properties: \{status: \{const: "active"\}\},
    required: [status]\}, then: \{required: [qty]\}} to the schema. Validate
    three instances: active-with-qty (pass), active-without-qty (fail),
    inactive-without-qty (pass) \cite{jsonschemadocs}.
  \item \texttt{additionalProperties: false} inside one \texttt{allOf}
    branch is evaluated \emph{per subschema}, so fields defined only in a
    sibling branch look ``additional'' and legitimate extensions are
    rejected. \texttt{unevaluatedProperties: false} at the top level is
    evaluated after all subschemas merge, fixing the composition;
    before/after validator output proves it \cite{jsonschemadocs}.
  \item \texttt{text/html} matches its explicit range at $q=0.7$ (most
    specific wins); \texttt{application/json} matches only \texttt{*/*} at
    $q=0.5$. Server score: html 0.7, json 0.5 --- html is served. The
    negotiation selector must confirm exactly this.
  \item Count records per second over a deterministically generated
    1{,}000{,}000-line file with a buffered reader. The early stop after 100
    records demonstrates generator laziness: the reader never touches line
    102 onward --- streaming gives constant memory and early-exit safety.
  \item Torn stream: a final line truncated mid-JSON by naive concatenation
    (\texttt{\{"a":1\}\textbackslash n\{"b":2}). The chapter's reader yields
    every complete record and treats the trailing partial line as
    incomplete (skip/quarantine), never crashing --- the property that makes
    NDJSON safe for append-oriented shipping.
  \item \texttt{object\_pairs\_hook} receives the raw pairs list; flag
    duplicates when key count $\neq$ unique key count. On
    \texttt{\{"a": 1, "a": 2\}} a first-wins parser returns 1 and a
    last-wins parser returns 2 --- a parser differential that attackers use
    to slip values past one layer of a system \cite{rfc8259}.
  \item \texttt{float(2**53 + 1) == float(2**53 + 2)} evaluates
    \texttt{True} in Python while the integers differ --- both collapse to
    the same IEEE-754 double. This simulates JSON parsers (JavaScript's
    \texttt{JSON.parse}) that decode all numbers as doubles; IDs above
    $2^{53}$ must travel as strings.
  \item Avro resolution permits \texttt{int}$\to$\texttt{long} because every
    32-bit value fits in 64 bits; the reverse is rejected because the
    \emph{value space}, not the current data, defines compatibility ---
    nothing guarantees tomorrow's writer stays in range. Forward
    compatibility is a property of schemas, not of observed datasets.
  \item The missing header is \texttt{Vary: Accept}; without it the CDN
    keys the cache entry on the URI alone and serves the cached v1
    representation to clients that negotiated \texttt{v=2}. The origin must
    emit \texttt{Vary: Accept} on every negotiated response
    \cite{rfc9110}.
\end{enumerate}

\subsection{Capstone Grading Rubrics}

\emph{Numbering note: the chapter labels its capstones ``Capstone 1--5'';
they are cited here as C4.1--C4.5.}

\subsubsection*{Capstone 4.1 --- NDJSON Inventory Log Shipper [junior]}
\begin{center}\small
\begin{tabularx}{\textwidth}{l c X}
\toprule
\textbf{Criterion} & \textbf{Weight} & \textbf{Full marks when\ldots} \\
\midrule
Validation accuracy & 30\% & all 9{,}950 clean records shipped; exactly the poisoned set quarantined --- no more, no fewer. \\
Atomicity & 25\% & a mid-batch kill leaves only a \texttt{.tmp} file; no torn line ever lands in \texttt{.ndjson}. \\
Bounded memory & 20\% & flat RSS across the full run (streaming, not load-all). \\
Reproducibility & 15\% & same seed $\to$ same shipped/quarantined split. \\
Error reporting & 10\% & quarantine summary names line numbers and schema violations. \\
\bottomrule
\end{tabularx}
\end{center}
\textbf{Reference approach.} Stream the input line by line; validate each
against the JSON Schema; write valid lines to a staging \texttt{.tmp},
quarantine invalid ones with reasons; atomically rename on batch
completion. Stretch: gzip rotation and an RFC~9457 problem summary
\cite{rfc9457}.

\subsubsection*{Capstone 4.2 --- Schema-Validation Middleware [mid]}
\begin{center}\small
\begin{tabularx}{\textwidth}{l c X}
\toprule
\textbf{Criterion} & \textbf{Weight} & \textbf{Full marks when\ldots} \\
\midrule
Error contract & 30\% & every invalid body returns 422 \texttt{application/problem+json} with a complete \texttt{errors[]} array of JSON Pointers. \\
Zero false rejections & 25\% & all 400 valid corpus bodies pass byte-identical. \\
Performance & 25\% & schemas compiled once at startup; measured per-request overhead documented. \\
Reusability & 20\% & route-to-schema mapping is configuration, not code edits. \\
\bottomrule
\end{tabularx}
\end{center}
\textbf{Reference approach.} WSGI middleware intercepting request bodies
for mapped routes, running a pre-compiled validator, and short-circuiting
with the problem document on failure. Stretch: content-negotiate the error
body (JSON vs HTML) using the chapter's selector.

\subsubsection*{Capstone 4.3 --- Hand-Rolled Protobuf Codec Extension [senior]}
\begin{center}\small
\begin{tabularx}{\textwidth}{l c X}
\toprule
\textbf{Criterion} & \textbf{Weight} & \textbf{Full marks when\ldots} \\
\midrule
Wire correctness & 30\% & byte-level golden tests (committed hex) for $\geq$5 messages covering packed repeated, enum, and embedded message \cite{protobufdocs}. \\
Unknown-field retention & 25\% & unknown fields survive a decode/encode round-trip byte-identical. \\
Evolution governance & 25\% & reserved-tag rules documented; a lint rejects tag reuse. \\
Engineering & 20\% & codec stays dependency-free; tests deterministic. \\
\bottomrule
\end{tabularx}
\end{center}
\textbf{Reference approach.} Extend the mini codec with a second message
type and nested-message framing (length-prefixed), packed varint runs for
repeated fields, and an unknown-fields bag carried on the decoded object.
Stretch: zigzag \texttt{sint64} deltas and a packed-vs-unpacked benchmark.

\subsubsection*{Capstone 4.4 --- Negotiation-Aware Cache Key Tool [senior]}
\begin{center}\small
\begin{tabularx}{\textwidth}{l c X}
\toprule
\textbf{Criterion} & \textbf{Weight} & \textbf{Full marks when\ldots} \\
\midrule
Detection accuracy & 35\% & every seeded mis-\texttt{Vary}'d response flagged; zero false positives. \\
Key stability & 25\% & variant keys deterministic and canonicalized (case, q-value normalization). \\
Explanation quality & 25\% & report explains each collision via the negotiation algorithm, not just ``mismatch''. \\
Offline rigor & 15\% & recorded-exchange fixtures; one-command run. \\
\bottomrule
\end{tabularx}
\end{center}
\textbf{Reference approach.} Parse recorded exchanges; compute the cache
key as URI plus the \texttt{Vary}-listed request headers (normalized);
simulate lookups across the corpus and report collisions and misses.
Stretch: quantify hit-rate loss from an overly broad \texttt{Vary:
User-Agent}.

\subsubsection*{Capstone 4.5 --- Format-Migration CLI [staff]}
\begin{center}\small
\begin{tabularx}{\textwidth}{l c X}
\toprule
\textbf{Criterion} & \textbf{Weight} & \textbf{Full marks when\ldots} \\
\midrule
Conversion completeness & 25\% & all $4\times4$ format pairs stream offline in bounded memory. \\
Compatibility gate & 25\% & seeded backward-incompatible edit blocked with precise reasons; compatible edit permitted. \\
Semantic verification & 25\% & verify mode proves semantic equality (including float tolerance), not byte equality. \\
Reporting \& ops & 25\% & dry-run impact report matches the benchmark harness within noise; versioned manifest committed; failure behavior documented. \\
\bottomrule
\end{tabularx}
\end{center}
\textbf{Reference approach.} Define a canonical in-memory record model;
write one streaming reader and writer per format; a schema manifest with
explicit compatibility classes drives the gate; verify mode replays both
datasets through the canonical model. Stretch: a machine-readable registry
manifest and a chaos mode that truncates output mid-stream, which verify
must catch.

%% ============================================================================
\section{Chapter 5: Client-Side Security}
\label{sec:sol-ch5}

\subsection{Self-Assessment Model Answers}

\emph{Numbering note: the chapter uses plain 1--7 exercises and 1--5
deep-dive questions; numbering below matches.}

\begin{enumerate}[leftmargin=2.2em]
  \item 22 base64url chars from \texttt{secrets}: $22\times6=132$ bits ---
    acceptable. UUIDv4: 122 bits --- acceptable. 32 hex chars from
    \texttt{random.Random(42)}: zero real entropy (deterministic PRNG, known
    seed) --- reject. \texttt{sha256(hostname+date)}: the input space is
    brute-forceable, so effective entropy is a few dozen bits --- reject.
    Only CSPRNG output counts; hashing low-entropy input does not create
    entropy.
  \item \texttt{canonical\_query} must sort parameters by key, then by
    value, and sign \emph{every} value of a repeated key; keeping only the
    first or last value means client and server sign different queries ---
    an interop break, and worse, a smuggling gap where an attacker appends
    an unsigned second \texttt{filter}. SigV4 requires sorted, fully
    enumerated pairs.
  \item Leeway must exceed worst-case clock skew plus latency: $20 +
    8 \approx 28$~s, so 30--60~s is defensible. Within the leeway window a
    captured token replays successfully; a \texttt{jti} cache with TTL equal
    to the leeway window closes that residual risk \cite{rfc7519}.
  \item Without \texttt{state}, an attacker can inject their own
    authorization code into the victim's redirect (login CSRF/code
    substitution), binding the victim's session to the attacker's account.
    PKCE protects the token \emph{exchange} from interception; it does not
    bind the redirect response to the browser that initiated the flow ---
    that is exactly \texttt{state}'s job \cite{rfc7636, rfc6749}.
  \item \texttt{exp} as a JSON string: reject --- claims must be type-checked
    before comparison. \texttt{aud} as a list containing your audience plus
    a foreign one: safest is reject unless multi-audience is explicitly
    supported (and then require \texttt{azp} binding); accepting blindly
    invites audience confusion. \texttt{crit: ["exp"]}: reject --- an
    unrecognized critical extension must fail closed, never be ignored.
  \item With 64 threads and a 50\,ms fake refresh, instrumentation must show
    exactly one refresh: the double-checked lock admits one thread to
    refresh while the rest find the fresh token. Removing the inner check
    lets every queued thread refresh after acquiring the lock --- a
    thundering herd against the token endpoint.
  \item Effective adversarial lines: secrets split across two log fields,
    URL-encoded tokens, JSON with irregular spacing. Regex redaction misses
    these; defense in depth: structured logging that never emits raw
    headers/bodies, allowlist field logging, and a DLP scan at log-ship
    time.
\end{enumerate}

\subsubsection*{Deep-Dive Questions}

\begin{enumerate}[leftmargin=2.2em]
  \item Signing a payload hash prevents \emph{tampering}, not
    \emph{replay}: a captured signed request can be resent verbatim.
    Timestamp + nonce bounds validity; the server-side replay cache is keyed
    by (signature or nonce) with TTL equal to the allowed clock window, and
    deduplication runs before any side effect.
  \item The front channel (browser redirect) is observable and modifiable
    --- URLs land in history, logs, and referrer headers --- so it must
    never carry tokens. The back channel (server-to-server) is
    authenticated and confidential. The split exists because the browser is
    a confused deputy: it transports, but cannot be trusted with secrets
    \cite{rfc6749}.
  \item Rotation reuse detection: presenting a retired refresh token
    revokes the whole family. False positive: the client crashed after
    receiving the new token but before persisting it, so it replays the old
    one. A robust client persists-then-discards (write the new token before
    deleting the old) and offers a re-auth recovery; a lenient server may
    accept the immediately-previous token once within a short grace window
    and flag it for review.
  \item Stateless JWT verification buys zero-lookup verification,
    horizontal scale, and latency; the price is a revocation gap until
    \texttt{exp} and stale claims (a revoked role lingers). Introspection
    wins when revocation immediacy matters (payments, admin tooling),
    tokens are long-lived, or per-request policy freshness is required ---
    it reintroduces the lookup deliberately.
  \item Auditor answer, in three moves: (1)~a mobile app is a public client
    --- any embedded secret is extractable from the binary regardless of
    obfuscation, so client credentials would be one shared secret whose
    extraction compromises every installation; (2)~authorization code +
    PKCE binds the exchange to the initiating device and user, issuing
    per-user tokens with consent; (3)~the flow preserves user identity for
    audit and per-user revocation, which a shared client secret cannot
    \cite{rfc7636, owaspapisecurity}.
\end{enumerate}

\subsection{Capstone Grading Rubrics}

\subsubsection*{Capstone 5.1 --- HMAC-Signed Webhook Consumer [junior]}
\begin{center}\small
\begin{tabularx}{\textwidth}{l c X}
\toprule
\textbf{Criterion} & \textbf{Weight} & \textbf{Full marks when\ldots} \\
\midrule
Verification correctness & 35\% & 17 valid events accepted; 3 corrupted events rejected with distinct reasons. \\
Replay defense & 25\% & replayed event rejected by the replay cache within the tolerance window. \\
Constant-time comparison & 25\% & \texttt{hmac.compare\_digest} used throughout --- grep finds no \texttt{==} on signatures. \\
Offline determinism & 15\% & byte-reproducible output; no network. \\
\bottomrule
\end{tabularx}
\end{center}
\textbf{Reference approach.} Parse \texttt{t=...,v1=...}, recompute HMAC
over \texttt{t.body} with the shared secret, compare in constant time,
enforce the timestamp tolerance, and record seen event IDs. Stretch:
two-secret rollover window and constant-time measurement.

\subsubsection*{Capstone 5.2 --- Full PKCE CLI [mid]}
\begin{center}\small
\begin{tabularx}{\textwidth}{l c X}
\toprule
\textbf{Criterion} & \textbf{Weight} & \textbf{Full marks when\ldots} \\
\midrule
Flow correctness & 30\% & complete auth-code+PKCE exchange against the stub IdP; \texttt{state} and verifier validated \cite{rfc7636}. \\
Token storage & 20\% & tokens stored owner-only-permission; zero token material in any log. \\
Refresh discipline & 25\% & proactive refresh; exactly one refresh under 8 concurrent invocations (single-flight). \\
Rotation handling & 25\% & retired-token replay triggers family revocation and a clean re-authentication path. \\
\bottomrule
\end{tabularx}
\end{center}
\textbf{Reference approach.} Loopback redirect listener for the code;
verifier/challenge from \texttt{secrets}; a token cache file with 0600
permissions guarded by a process-wide lock for single-flight refresh.
Stretch: device-code mode and OS keychain with file fallback.

\subsubsection*{Capstone 5.3 --- JWT Security Linter [mid]}
\begin{center}\small
\begin{tabularx}{\textwidth}{l c X}
\toprule
\textbf{Criterion} & \textbf{Weight} & \textbf{Full marks when\ldots} \\
\midrule
Rule coverage & 35\% & $\geq$6 anti-pattern categories detected with file/line attribution (algorithm confusion, \texttt{==} comparison, missing \texttt{aud}/\texttt{exp} checks, \ldots). \\
Precision & 25\% & zero false positives on the clean corpus. \\
Traceability & 20\% & each rule cites a chapter subsection or OWASP category \cite{owaspapisecurity}. \\
Determinism & 20\% & sorted, byte-stable output across runs. \\
\bottomrule
\end{tabularx}
\end{center}
\textbf{Reference approach.} Walk the \texttt{ast} for JWT-related call
sites; match call patterns against rule predicates; report through a
stable, sorted finding model. Stretch: auto-fix mode and pre-commit
integration.

\subsubsection*{Capstone 5.4 --- Token-Broker Service [senior]}
\begin{center}\small
\begin{tabularx}{\textwidth}{l c X}
\toprule
\textbf{Criterion} & \textbf{Weight} & \textbf{Full marks when\ldots} \\
\midrule
End-to-end flow & 30\% & login $\to$ broker $\to$ scoped call works; scope escalation rejected. \\
Secret hygiene & 25\% & no long-lived secret in any client process or file. \\
Audit trail & 25\% & audit log carries \texttt{sub}/\texttt{jti}/scope/expiry and zero token material. \\
Failure behavior & 20\% & broker outage yields bounded client failure (one retry, exit 2) --- no hang, no credential fallback. \\
\bottomrule
\end{tabularx}
\end{center}
\textbf{Reference approach.} The broker holds the IdP credential, mints
short-lived task-scoped tokens, and logs issuance; clients authenticate to
the broker with their developer identity. Stretch: a \texttt{jti} blocklist
endpoint with propagation measurement and RS256 issuance.

\subsubsection*{Capstone 5.5 --- Red-Team Your Own Client [staff]}
\begin{center}\small
\begin{tabularx}{\textwidth}{l c X}
\toprule
\textbf{Criterion} & \textbf{Weight} & \textbf{Full marks when\ldots} \\
\midrule
Threat model & 25\% & $\geq$6 abuse cases traced to concrete code locations. \\
Findings \& fixes & 30\% & $\geq$3 genuine findings in the chapter listings, each with a landed fix and regression test. \\
CI gate & 25\% & gate fails when a vulnerability is reintroduced and passes on the fixed tree. \\
Residual-risk report & 20\% & accepted findings quantified; rollback/exception process stated. \\
\bottomrule
\end{tabularx}
\end{center}
\textbf{Reference approach.} Work the attack surface methodically (token
handling, signature verification, storage, logging); convert every exploit
into a failing test, then fix; wire the corpus into CI. Stretch:
differential-fuzz the verifier against a reference library over 10{,}000
seeded tokens.

%% ============================================================================
\section{Chapter 6: Building APIs with FastAPI}
\label{sec:sol-ch6}

\subsection{Self-Assessment Model Answers}

\emph{Numbering note: the chapter labels exercises E1--E8 (no chapter
prefix); numbering below matches.}

\begin{enumerate}[label=\textbf{E\arabic*.}, leftmargin=2.2em]
  \item \texttt{def get\_api\_version(x\_api\_version: str = Header("1"))}
    exposed through a typed alias (\texttt{Annotated[str, Depends(...)]});
    tests assert the default \texttt{"1"} without the header and the
    override with it. The alias keeps every signature self-documenting
    \cite{fastapidocs}.
  \item Store \texttt{(key $\to$ body hash, stored response)}; same key +
    same body replays the original 201 without re-executing; same key +
    different body returns 409 (fingerprint mismatch). This is the
    idempotency contract of Chapter~\ref{ch:idempotency} applied to
    \texttt{POST /orders}.
  \item The middleware schedules \texttt{asyncio.sleep(0)} and measures the
    delay until it runs --- that delay \emph{is} event-loop lag. A
    deliberately blocking endpoint (CPU loop or \texttt{time.sleep}) spikes
    the metric for \emph{all} concurrent requests; moving the work to a
    threadpool (or making it genuinely async) flattens it. Lag is the
    observable; blocking is the cause.
  \item The counter must live on \texttt{app.state} because middleware
    instances are not the canonical shared store --- the application object
    owns lifespan-managed state, and counters on middleware attributes can
    be recreated or fragmented across instances. Pure ASGI middleware keeps
    the count below the framework layer, so it sees every request.
  \item Duplicate-SKU is a persistence invariant: only the repository can
    see all existing rows, and a Pydantic \texttt{model\_validator} couples
    the rule to one transport DTO (a second write path would bypass it).
    Moving it to \texttt{insert} as a domain error mapped to 422 is correct,
    and it implies cross-aggregate invariants (max five open orders per
    customer) also belong in the repository/service layer, never in DTOs.
  \item Paging the repository in limit/offset batches of 100 inside the
    generator bounds memory at one batch and lets the stream observe
    committed writes between batches. Snapshotting prevents iterator
    invalidation and mid-yield mutation, but at the cost of pinning the
    entire result set in memory for the stream's lifetime --- the failure
    mode this exercise trades away deliberately.
  \item Exiting the TestClient context runs the lifespan shutdown, which
    calls the repository's \texttt{close()} --- asserted by a flag on the
    repo. Uvicorn differs by being signal-driven: lifespan shutdown runs on
    SIGTERM during graceful shutdown (documented in uvicorn's lifespan
    section), so the same code path must be verified under a real server
    before production (Chapter~\ref{ch:deployment}).
  \item The test diffs live \texttt{/openapi.json} against a checked-in
    snapshot; renaming a response field changes the document and fails CI.
    The point: your OpenAPI output \emph{is} the contract, and silent drift
    becomes a build failure rather than a customer surprise
    \cite{openapi31}.
\end{enumerate}

\subsection{Capstone Grading Rubrics}

\emph{Numbering note: the chapter labels its capstones ``Capstone 1--5'';
they are cited here as C6.1--C6.5.}

\subsubsection*{Capstone 6.1 --- Todo API with a Full Test Suite [junior]}
\begin{center}\small
\begin{tabularx}{\textwidth}{l c X}
\toprule
\textbf{Criterion} & \textbf{Weight} & \textbf{Full marks when\ldots} \\
\midrule
CRUD correctness & 30\% & all routes behave per contract, including the complete transition, with deterministic IDs. \\
Response discipline & 25\% & every route declares \texttt{response\_model}; no internal field leaks in any response. \\
Contract documentation & 20\% & OpenAPI documents a 422 for every body-accepting route. \\
Test suite & 25\% & pytest passes offline from a fresh venv; in-memory repo behind a protocol. \\
\bottomrule
\end{tabularx}
\end{center}
\textbf{Reference approach.} Repository protocol + in-memory
implementation injected via dependency overrides; TestClient covering every
route and error path. Stretch: correlation-ID middleware with header echo
and an NDJSON streaming endpoint with a streaming test.

\subsubsection*{Capstone 6.2 --- Inventory API with a SQLite Repository [mid]}
\begin{center}\small
\begin{tabularx}{\textwidth}{l c X}
\toprule
\textbf{Criterion} & \textbf{Weight} & \textbf{Full marks when\ldots} \\
\midrule
Repository parity & 30\% & the identical suite passes against \texttt{memory} and \texttt{sqlite} with zero code changes. \\
Conflict semantics & 25\% & oversubscription returns a 409 problem document naming part and shortfall \cite{rfc9457}. \\
Readiness honesty & 25\% & readiness returns 503 when the SQLite path is a directory --- a real failure, not a mock. \\
Settings plumbing & 20\% & backend selection via settings only; no branching in route code. \\
\bottomrule
\end{tabularx}
\end{center}
\textbf{Reference approach.} Rebuild the Orders pattern for parts, stock,
and reservations; the readiness probe opens the configured database for
real. Stretch: optimistic concurrency with a \texttt{version} field and 412
on mismatch (Chapter~\ref{ch:idempotency}).

\subsubsection*{Capstone 6.3 --- Plugin-Style Dependency-Injection Architecture [senior]}
\begin{center}\small
\begin{tabularx}{\textwidth}{l c X}
\toprule
\textbf{Criterion} & \textbf{Weight} & \textbf{Full marks when\ldots} \\
\midrule
Plugin ergonomics & 30\% & a new plugin is one new file plus one registry line --- no route or factory edits. \\
Fail-fast config & 25\% & unknown plugin name fails at startup with an actionable error. \\
Test replaceability & 25\% & \texttt{dependency\_overrides} tests prove per-test plugin swapping. \\
Composition clarity & 20\% & lifespan composes the selected plugins explicitly; startup logs the effective set. \\
\bottomrule
\end{tabularx}
\end{center}
\textbf{Reference approach.} Define protocols for repositories, notifiers,
and pricing policies; a name $\to$ factory registry resolved in the
lifespan; routes depend on protocols only. Stretch: expose plugin versions
in \texttt{/health} and a contract test that every notifier satisfies a
shared behavioral suite.

\subsubsection*{Capstone 6.4 --- Async Migration of a Blocking Service [senior]}
\begin{center}\small
\begin{tabularx}{\textwidth}{l c X}
\toprule
\textbf{Criterion} & \textbf{Weight} & \textbf{Full marks when\ldots} \\
\midrule
Baseline proof & 25\% & loop-lag middleware shows $>$5~s lag under the 200-way concurrency load before the fix. \\
Fix efficacy & 30\% & both fixes (def-threadpool and \texttt{asyncio.to\_thread}) bring lag under 50~ms; p99 improves by an order of magnitude over three reproducible runs. \\
Evidence quality & 25\% & decision memo cites pool-occupancy and lag measurements, not folklore. \\
Reproducibility & 20\% & benchmark harness is deterministic and offline. \\
\bottomrule
\end{tabularx}
\end{center}
\textbf{Reference approach.} Keep the blocking implementation behind the
repository protocol, swap implementations, and let the lag middleware plus
the concurrency benchmark adjudicate. Stretch: a genuinely async store and
threadpool queue-depth instrumentation.

\subsubsection*{Capstone 6.5 --- OpenAPI-First Generator Round-Trip [staff]}
\begin{center}\small
\begin{tabularx}{\textwidth}{l c X}
\toprule
\textbf{Criterion} & \textbf{Weight} & \textbf{Full marks when\ldots} \\
\midrule
Gate correctness & 30\% & seeded compatible and breaking changes classified correctly with readable reasons; running service passes its own gate with zero round-trip diffs. \\
Round-trip fidelity & 25\% & generated typed client works against the service in a TestClient integration test without hand edits. \\
Governance policy & 25\% & written policy states who may bump which version number and how breaking changes ship (Chapter~\ref{ch:versioning}). \\
Operational rollout & 20\% & gate wired into CI with documented override and rollback path. \\
\bottomrule
\end{tabularx}
\end{center}
\textbf{Reference approach.} Hand-author the OpenAPI~3.1 spec as the
source of truth \cite{openapi31}; generate server skeleton and typed
client; implement; build the diff gate over normalized specs; then execute
one governed compatible change and demonstrate one blocked breaking change.
Stretch: diff report as an RFC~9457 problem document and
undocumented-endpoint detection.

%% ============================================================================
\section{Chapter 7: Errors, Validation \& Problem Details}
\label{sec:sol-ch7}

\subsection{Self-Assessment Model Answers}

\emph{Numbering note: the chapter uses a plain 1--10 list; numbering below
matches.}

\begin{enumerate}[leftmargin=2.2em]
  \item Malformed JSON at \texttt{POST /shipments} fails at the
    \emph{transport/syntax} layer of the taxonomy: the body never becomes a
    data structure, so Pydantic is never invoked. The result is 400 (bad
    request --- the bytes are at fault), not 422, and the problem body
    should say the JSON itself could not be parsed \cite{rfc9457}.
  \item 404-for-unauthorized is correct when resource \emph{existence} is
    sensitive (private tenant data --- IDOR probing must learn nothing). It
    still leaks when existence is inferable anyway (enumerable public IDs,
    timing differences between ``exists but forbidden'' and ``absent'', or
    when a 403-only endpoint elsewhere confirms the namespace). Then 404
    only confuses legitimate clients while fooling no one.
  \item Safe retry of a conditionally-retriable 409: (1)~on 409, re-read the
    current resource state; (2)~re-evaluate the precondition locally;
    (3)~if the conflict is semantic (duplicate key), abort --- no retry
    helps; (4)~otherwise resubmit with a fresh precondition
    (\texttt{If-Match}/version); (5)~bound attempts with jittered backoff.
    Pseudo-code must show the re-read --- blind resubmission just re-loses.
  \item \texttt{unit\_weight\_kg} is \texttt{strict=False} because real
    clients send JSON strings like \texttt{"1.5"}; removing the opt-out
    rejects those payloads (a compatibility break). Removing the companion
    validator as well drops the range/unit checks, so negative or absurd
    weights pass --- the two knobs guard different failure classes and both
    are needed.
  \item JSON Pointer: \texttt{/lines/2/dimensions/width} (array indices are
    decimal positions). Per RFC~6901 escaping, \texttt{\textasciitilde}
    becomes \texttt{\textasciitilde 0} and \texttt{/} becomes
    \texttt{\textasciitilde 1}, so a key \texttt{a/b} is addressed as
    \texttt{/a\textasciitilde 1b}. The pointer must survive this escaping
    round-trip exactly.
  \item Misreading breaker-503s as application 500s sent the incident to the
    wrong team with the wrong hypothesis (code bug vs dependency/capacity),
    burning the golden first hour. A metric label carrying the failure
    \emph{origin} (e.g.\ \texttt{error\_source="circuit\_breaker"})
    alongside \texttt{status\_code} prevents recurrence.
  \item \emph{For} localizing \texttt{title}: developer experience for
    global consumers. \emph{Against}: it breaks log aggregation, alerting
    rules, and support search that key on the stable string. The invariant
    that makes localization safe: \texttt{type} (and any \texttt{code}
    extension) remain constant and machine-consumable; only human-readable
    fields localize \cite{rfc9457}.
  \item Apply the scrubber per chunk with a sliding overlap window so
    secrets spanning a chunk boundary are still caught; accept the
    trade-off: boundary guarantees are approximate (window-bound) and
    over-redaction near boundaries is possible --- in exchange for never
    buffering a 50\,MB stream.
  \item \texttt{additionalProperties: false} (or
    \texttt{unevaluatedProperties} across composed schemas) on the problem
    schema makes any undocumented extension member fail validation in CI
    \cite{jsonschemadocs}. Implied workflow: registry-first --- the schema
    and the emitter change land in the same PR, so the gate and the
    contract move together.
  \item Map alert thresholds onto error-budget burn rates
    (Chapter~\ref{ch:observability}): page on fast-burn 5xx, ticket on
    slow-burn. The code that should never page is 4xx (except
    authentication-abuse spikes): client errors are caller-owned, and paging
    on them trains the on-call to ignore pages.
\end{enumerate}

\subsection{Capstone Grading Rubrics}

\emph{Numbering note: the chapter labels its capstones ``Project 1--5'';
they are cited here as C7.1--C7.5.}

\subsubsection*{Capstone 7.1 --- Problem-Details Retrofit [junior]}
\begin{center}\small
\begin{tabularx}{\textwidth}{l c X}
\toprule
\textbf{Criterion} & \textbf{Weight} & \textbf{Full marks when\ldots} \\
\midrule
Uniformity & 30\% & every error --- including framework-generated 404s --- is schema-valid \texttt{application/problem+json} \cite{rfc9457}. \\
Success-path preservation & 25\% & no success-path test changed or broken. \\
Validation errors & 25\% & 422s carry JSON-Pointer-located \texttt{errors[]}. \\
Runnability & 20\% & pytest green from a fresh venv, offline. \\
\bottomrule
\end{tabularx}
\end{center}
\textbf{Reference approach.} Install exception handlers for the framework's
request/validation/HTTP exceptions plus a catch-all; centralize the problem
serializer; prove with a route-by-route error matrix. Stretch:
\texttt{Accept-Language}-localized \texttt{title} for two locales.

\subsubsection*{Capstone 7.2 --- Error-Code Registry \& Documentation Generator [mid]}
\begin{center}\small
\begin{tabularx}{\textwidth}{l c X}
\toprule
\textbf{Criterion} & \textbf{Weight} & \textbf{Full marks when\ldots} \\
\midrule
Registry validation & 30\% & import-time validation catches an injected duplicate code and a missing remediation field. \\
Generated artifacts & 25\% & Markdown docs and \texttt{errors.json} validate against the published schema. \\
Orphan detection & 25\% & the CI check fails on a service emitting an unregistered code. \\
Determinism & 20\% & version-diff output is byte-stable and sorted. \\
\bottomrule
\end{tabularx}
\end{center}
\textbf{Reference approach.} Declarative registry (one record per code:
status, title, remediation, since-version) validated at import; generators
render docs and JSON from the same records; the orphan check runs the
service's error paths against the registry. Stretch: per-code deep links as
\texttt{type} URIs and a mark-never-delete deprecation workflow.

\subsubsection*{Capstone 7.3 --- Redaction Library with Fuzz Corpus [mid]}
\begin{center}\small
\begin{tabularx}{\textwidth}{l c X}
\toprule
\textbf{Criterion} & \textbf{Weight} & \textbf{Full marks when\ldots} \\
\midrule
Corpus coverage & 30\% & 100\% of the 50+ adversarial corpus redacted (split tokens, encodings, odd spacing). \\
False-positive discipline & 25\% & documented bounded FP rate with a CI regression gate. \\
Extensibility & 20\% & new rules added without editing library code. \\
Determinism & 25\% & seeded mutation sweeps; fully offline; reproducible reports. \\
\bottomrule
\end{tabularx}
\end{center}
\textbf{Reference approach.} Harden the chapter scrubber behind a rule
pipeline (pattern + replacement + confidence); attack it with the curated
corpus plus seeded mutation; publish the FP budget. Stretch: streaming
redaction with boundary windows and JSON-aware field-name scrubbing.

\subsubsection*{Capstone 7.4 --- Error-Observability Dashboard Specification [senior]}
\begin{center}\small
\begin{tabularx}{\textwidth}{l c X}
\toprule
\textbf{Criterion} & \textbf{Weight} & \textbf{Full marks when\ldots} \\
\midrule
Incident detection & 30\% & replay of the seeded week detects all three injected incidents with zero false pages on baseline. \\
Declarative rules & 25\% & alert rules are data (thresholds, windows, codes), not code. \\
Traceability & 25\% & every dashboard panel traces to a registry error code. \\
Burn-rate correctness & 20\% & the burn-rate panel matches a hand-computed 100-request example. \\
\bottomrule
\end{tabularx}
\end{center}
\textbf{Reference approach.} Derive the metric set from the error registry
(rate per code, burn-rate windows); replay a seeded log week through the
alert evaluator; render the dashboard spec as data. Stretch: fast/slow
burn-rate windows per Chapter~\ref{ch:observability} and per-code anomaly
baselines.

\subsubsection*{Capstone 7.5 --- Chaos Error-Injection Test Harness [staff]}
\begin{center}\small
\begin{tabularx}{\textwidth}{l c X}
\toprule
\textbf{Criterion} & \textbf{Weight} & \textbf{Full marks when\ldots} \\
\midrule
Chain integrity & 30\% & every fault yields a schema-valid problem response in the correct 5xx family --- never a raw 500 or HTML. \\
Correlation & 25\% & trace IDs join both services' logs for every scenario. \\
Determinism \& speed & 20\% & fixed seed, fully offline, full run under one minute. \\
Coverage matrix & 25\% & a fault $\times$ contract matrix demonstrates no blind spots; rollback of injected state verified between scenarios. \\
\bottomrule
\end{tabularx}
\end{center}
\textbf{Reference approach.} In-process fault injection (timeouts, malformed
upstream payloads, resets, slow leaks) between two toy services; assert the
whole error chain --- status mapping, problem shape, redaction,
correlation, metrics --- under each fault. Stretch: composed faults and a
``contract score'' trending across registry versions.

%% ============================================================================
\section{Chapter 8: Pagination, Filtering \& Sorting}
\label{sec:sol-ch8}

\subsection{Self-Assessment Model Answers}

\begin{enumerate}[label=\textbf{E8.\arabic*.}, leftmargin=2.8em]
  \item Page $k$ costs $c_{\text{seek}} + k\,c_s + p\,c_f$; summing over
    $n/p$ pages gives total row visits
    $\sum_{k} k p \approx p\,\frac{(n/p)(n/p+1)}{2} \approx
    \frac{n^2}{2p}$ --- $\Theta(n^2/p)$. The worst page size from the
    server's viewpoint is $p{=}1$: maximal pages, maximal total work
    ($\Theta(n^2)$).
  \item One valid schedule (page size 5): page 1 emits r1--r5; page 2
    (offset 5) emits r6--r10; then the writer \emph{moves} r7 to the end of
    the table (position 20) and \emph{moves} r20 into position 7 (e.g.\ via
    sort-key updates). Page 3 (offset 10) now reads r11--r15; page 4
    (offset 15) reads r16, r17, r18, r19, and the relocated r7. Final
    multiset: r1--r19 with r7 twice; r20, sitting in the already-read
    region, is never emitted.
  \item Re-reading the previous page's last row detects overlap (duplicates)
    at the boundary, but a skipped row leaves no positional trace --- the
    position stream carries no identity information, so detection without
    recovery is the best any positional scheme can do; identity-bearing
    cursors are the fix.
  \item Embed a SHA-256 hash of the active filter/sort parameters inside the
    signed cursor; a continuation whose presented parameters no longer match
    the hash is rejected with 400. This closes cursor-replay and
    confused-deputy bugs where a cursor minted under one filter is replayed
    against another --- the server would otherwise silently return a
    different result set than the client believes it is walking.
  \item Cap page size and maximum page depth; cache or approximate the
    COUNT with a TTL (exact counts are a separate, expensive product);
    back the grid with a covering index; offer background export for deep
    access. Tell the PM the physics: deep offset pages cost
    $\Theta(n^2/p)$ server work --- page 47 is a filter/search problem, not
    a browsing problem.
  \item By trichotomy: if $a > x$ the tuple is greater regardless of $b$;
    if $a = x$ the comparison reduces to $b > y$; if $a < x$ both sides are
    false --- hence $(a,b)>(x,y) \iff a>x \lor (a=x \land b>y)$. Six-row
    counterexample for the hand-rolled \texttt{a >= x AND b > y} with cursor
    $(1,2)$ over rows $(1,1),(1,2),(1,3),(2,1),(2,2),(2,3)$: the buggy
    predicate drops $(2,1)$ and $(2,2)$ (skipped rows), while the mirror bug
    \texttt{a > x OR b > y} re-admits already-seen rows (duplicates) ---
    the two distinct error modes of offset-free predicates done wrong.
  \item Order by the null-flag first: \texttt{ORDER BY (discontinued\_at IS
    NULL), discontinued\_at, id} and build the keyset predicate over the
    triple --- the boolean null flag becomes the leading tuple element, so
    the disjunctive expansion handles the NULL boundary exactly like any
    tie-break; verify against the filter\_compiler demo database with
    \texttt{EXPLAIN QUERY PLAN}.
  \item Expansion policy: cap depth at 2, batch expansions with per-resource
    fan-out limits (e.g.\ 50), set \texttt{Cache-Control} to
    \texttt{private}/\texttt{no-store} when expansions embed per-user data
    and short \texttt{max-age} otherwise. When selections become arbitrary
    or deeper than the cap, redirect the consumer to the GraphQL gateway ---
    that is the workload GraphQL exists for (Chapter~\ref{ch:graphql}).
  \item Create an index whose column order cannot serve both the range
    predicate and the ORDER BY (e.g.\ \texttt{(id, posted\_at)} where the
    keyset needs \texttt{(posted\_at, id)}), or one leading with a
    low-cardinality column: the plan flips from \texttt{SEARCH ... USING
    INDEX} to \texttt{SCAN} plus a temp B-tree. The planner is saying no
    available index satisfies filter + ordering together --- the defining
    requirement of keyset pagination.
  \item Two architectures: (1)~transactionally maintained count tables
    (failure: drift and write amplification --- mitigate with reconciliation
    jobs); (2)~an asynchronous export/audit service running exact
    \texttt{COUNT(*)} over the indexed predicate with statement timeouts
    and as-of timestamps (failure: latency and staleness --- bounded by
    reporting the count's observation time). Either way, exact totals are
    isolated from the list API so metadata cost never taxes interactive
    reads.
\end{enumerate}

\subsection{Capstone Grading Rubrics}

\emph{Numbering note: the chapter labels its capstones ``Capstone 1--5'';
they are cited here as C8.1--C8.5.}

\subsubsection*{Capstone 8.1 --- Cursor-Paginated Ledger API [junior]}
\begin{center}\small
\begin{tabularx}{\textwidth}{l c X}
\toprule
\textbf{Criterion} & \textbf{Weight} & \textbf{Full marks when\ldots} \\
\midrule
Walk correctness & 30\% & a scripted full walk emits every one of the 100k postings exactly once (set comparison). \\
Index discipline & 25\% & \texttt{EXPLAIN} shows \texttt{SEARCH} using the composite \texttt{(posted\_at, ledger\_id)} index on every page. \\
Cursor security & 25\% & forged, expired, and version-mismatched cursors each receive the same uniform 400 problem body. \\
Tie handling & 20\% & timestamp ties across page boundaries produce no duplicates or skips. \\
\bottomrule
\end{tabularx}
\end{center}
\textbf{Reference approach.} Order by the composite key; the cursor is an
HMAC-signed payload of the last row's key tuple; the predicate is the
disjunctive tuple expansion. Stretch: \texttt{rel="prev"} via reversed-order
query and per-account filtering on an index-led equality prefix.

\subsubsection*{Capstone 8.2 --- Filter Compiler with Property Tests [mid]}
\begin{center}\small
\begin{tabularx}{\textwidth}{l c X}
\toprule
\textbf{Criterion} & \textbf{Weight} & \textbf{Full marks when\ldots} \\
\midrule
Injection safety & 30\% & property tests prove no raw input ever reaches SQL; all 50 attack strings rejected or provably inert when executed. \\
Grammar \& OR support & 25\% & documented grammar; OR-groups compile correctly. \\
Namespace hygiene & 20\% & no public parameter name maps to an unintended internal column. \\
Parameter binding & 25\% & inspection test shows compiled queries use bound parameters exclusively. \\
\bottomrule
\end{tabularx}
\end{center}
\textbf{Reference approach.} Parse to an AST, compile the AST to SQL with
placeholders, never interpolate; the property test mutates inputs and
asserts the compiled SQL text is input-free. Stretch: case-insensitive
\texttt{contains} with correct collation and an EXPLAIN-driven
index-recommendation report.

\subsubsection*{Capstone 8.3 --- Pagination Benchmark Report [mid]}
\begin{center}\small
\begin{tabularx}{\textwidth}{l c X}
\toprule
\textbf{Criterion} & \textbf{Weight} & \textbf{Full marks when\ldots} \\
\midrule
Reproducibility & 25\% & one command regenerates every table from fixtures. \\
Asymptotic evidence & 30\% & offset growth demonstrably linear (fit reported); keyset flat within noise; cursor overhead constant with depth. \\
Index evidence & 25\% & every query named with its \texttt{EXPLAIN} plan. \\
Scale coverage & 20\% & 100k / 200k / 1M rows, p50/p95, cold-cache variant. \\
\bottomrule
\end{tabularx}
\end{center}
\textbf{Reference approach.} Seed the three table sizes; sweep page depths
under each scheme; fit offset latency vs depth to a line and report
residuals. Stretch: a covering-index variant showing the offset penalty
shrink but not vanish, and a PowerShell runner for the matrix.

\subsubsection*{Capstone 8.4 --- Infinite-Scroll Client over Signed Cursors [senior]}
\begin{center}\small
\begin{tabularx}{\textwidth}{l c X}
\toprule
\textbf{Criterion} & \textbf{Weight} & \textbf{Full marks when\ldots} \\
\midrule
Opaque-cursor discipline & 25\% & client source contains no offset arithmetic and no cursor parsing --- it follows \texttt{Link rel="next"} only \cite{rfc8288}. \\
Resume correctness & 30\% & kill-and-resume completes the walk with zero duplicates using on-disk state. \\
Failure drills & 25\% & expired and forged cursors get controlled, user-visible recovery; server-restart drill passes. \\
Model fidelity & 20\% & mid-walk-insert behavior note matches the chapter's drift model. \\
\bottomrule
\end{tabularx}
\end{center}
\textbf{Reference approach.} Terminal browser persisting the last
\texttt{next} link after every page; on restart, resume from it; on cursor
rejection, offer refetch-from-filter with a clear notice. Stretch: per-page
ETag conditional requests and silent offset-to-cursor link negotiation.

\subsubsection*{Capstone 8.5 --- Multi-Tenant Keyset Design [staff]}
\begin{center}\small
\begin{tabularx}{\textwidth}{l c X}
\toprule
\textbf{Criterion} & \textbf{Weight} & \textbf{Full marks when\ldots} \\
\midrule
Isolation & 25\% & cross-tenant cursor replay is impossible, proven by test (tenant id inside the signed cursor and the predicate). \\
Depth independence & 25\% & all tenants (1k / 200k / 5M rows) show depth-independent page latency, measured and tabulated. \\
Resilience & 20\% & whale-tenant export resumes after a simulated crash with no loss. \\
ADR quality & 30\% & every rejected trade-off (offsets, per-page counts, transparent cursors) named with reasons; operational criteria (index ownership, monitoring, rollback) included. \\
\bottomrule
\end{tabularx}
\end{center}
\textbf{Reference approach.} Tenant-led composite index
\texttt{(tenant\_id, sort key\ldots)}; the cursor binds tenant and
position; per-tenant policy table for page size and export quotas; the ADR
carries the decision record for the org. Stretch: row-level security via
views/roles and snapshot-id cursors for compliance exports.

%% ============================================================================
\section{Chapter 9: Versioning \& Evolution}
\label{sec:sol-ch9}

\subsection{Self-Assessment Model Answers}

\begin{enumerate}[label=\textbf{E9.\arabic*.}, leftmargin=2.8em]
  \item Optional response field added: \emph{non-breaking} (additive).
    Optional request field made required: \emph{breaking}. \texttt{price}
    $\to$ \texttt{unit\_price\_cents} (rename + retype): \emph{breaking} ---
    twice. JSON member reordering: \emph{non-breaking} --- object member
    order is insignificant \cite{rfc8259}. Changed 404 \texttt{detail}
    message: formally non-breaking (\texttt{detail} is human-readable), but
    \emph{semantically hazardous} if clients parse it --- a contract-test
    finding, not a schema one. New 410 for retired resources: a
    \emph{behavioral} change clients must handle; announce it like a
    breaking change even though the schema is untouched.
  \item The gate flags as breaking: a newly \emph{added required} path/query
    parameter and any \emph{removed} parameter; fixtures prove both
    directions (old consumers cannot call the new shape; new definition no
    longer accepts old calls). Optional added parameters and added optional
    request body fields pass.
  \item With \texttt{Accept:
    application/vnd.nr.v2+json;q=0.3, application/vnd.nr.*;q=0.9}, the
    higher-q wildcard wins and \texttt{negotiate} returns the default
    version --- surprising but correct q-value semantics. Reverse the
    weights (explicit 0.9, wildcard 0.3) and v2 is chosen: explicit intent
    must outbid the wildcard.
  \item A multi-currency price list cannot be back-translated into v1's
    scalar price. The three honest options: (1)~ship the field in v2 only
    (v1 consumers simply do not get the feature); (2)~expose a documented
    lossy approximation in v1 (e.g.\ base-currency only, flagged);
    (3)~schedule v1 retirement. Expand-and-contract prescribes (1) with a
    managed path to (3): add alongside, never mutate meaning in place, and
    let usage telemetry drive the sunset.
  \item Required response headers: \texttt{Vary: X-API-Version} on every
    response, plus a deliberate \texttt{Cache-Control} per versioned
    resource. If a CDN ignores \texttt{Vary}, it keys entries on URI alone
    and serves a cached v1 representation to v2 requesters --- a contract
    violation that looks, maddeningly, like an application bug
    \cite{rfc9110}.
  \item \emph{Defend} media-type versioning here: the ecosystem is closed
    (you ship the clients), generated models tolerate exact content types,
    and URIs stay stable. The six-month release tail dictates the timeline:
    support N-2 versions for at least 6--9 months, instrument per-version
    traffic, and use a minimum-version killswitch only with a consumer
    notification runway.
  \item Publish a security-exception clause \emph{before} the incident:
    severe vulnerabilities may force immediate breaking change within a
    version, with (a)~simultaneous notification on every channel,
    (b)~minimal blast radius (revoke/compensate rather than redesign),
    (c)~parallel-run or extended overlap wherever technically possible, and
    (d)~a public post-incident review. The clause converts a contract
    violation into a governed emergency procedure.
\end{enumerate}

\subsection{Capstone Grading Rubrics}

\subsubsection*{Capstone 9.1 --- Add v2 to an Existing API with an Adapter Layer [junior]}
\begin{center}\small
\begin{tabularx}{\textwidth}{l c X}
\toprule
\textbf{Criterion} & \textbf{Weight} & \textbf{Full marks when\ldots} \\
\midrule
Dual-version correctness & 30\% & v1 and v2 both respond correctly for every seeded record after the rename+retype. \\
Deprecation signaling & 25\% & v1 responses carry \texttt{Deprecation}, \texttt{Sunset}, and \texttt{Link} headers; v2 carries none \cite{rfc8594}. \\
Adapter purity & 25\% & \texttt{to\_v1} is a pure function; no duplicated business logic; adapter tests deterministic and offline. \\
Test coverage & 20\% & both surfaces covered for identical underlying state. \\
\bottomrule
\end{tabularx}
\end{center}
\textbf{Reference approach.} Evolve the internal model to v2; implement
\texttt{to\_v1} as a pure projection; mount both routers over one
repository. Stretch: a third experimental version behind an opt-in header
and an adapter-overhead timing loop.

\subsubsection*{Capstone 9.2 --- Media-Type Version Negotiation Middleware [mid]}
\begin{center}\small
\begin{tabularx}{\textwidth}{l c X}
\toprule
\textbf{Criterion} & \textbf{Weight} & \textbf{Full marks when\ldots} \\
\midrule
Negotiation correctness & 35\% & every outcome matches the published decision table (explicit, wildcard, q-values, default). \\
Error quality & 20\% & 406 bodies enumerate the supported \texttt{version=N} values. \\
Cache safety & 25\% & \texttt{Vary: Accept} present on every 200, proven by test. \\
Suite health & 20\% & full pytest suite passes offline. \\
\bottomrule
\end{tabularx}
\end{center}
\textbf{Reference approach.} Generalize the chapter's negotiation listing
into middleware that resolves the vendor media type once per request and
dispatches to versioned serializers on a single \texttt{/catalog} URI.
Stretch: negotiate a second wire format (\texttt{+json} vs
\texttt{+cbor}-style suffix) orthogonal to version.

\subsubsection*{Capstone 9.3 --- Production Breaking-Change CI Gate [mid]}
\begin{center}\small
\begin{tabularx}{\textwidth}{l c X}
\toprule
\textbf{Criterion} & \textbf{Weight} & \textbf{Full marks when\ldots} \\
\midrule
Verdict correctness & 35\% & exit code 1 exactly when breaking changes lack a major bump; all five fixture pairs produce expected verdicts. \\
Diff completeness & 25\% & local \texttt{\$ref} resolution and parameter diffing implemented. \\
Reporting & 20\% & JSON report round-trips through \texttt{json.loads}; enum-growth is a warning tier, not a failure. \\
Escape hatch & 20\% & \texttt{--allow-breaking-with=VERSION\_BUMP} works and is audit-logged. \\
\bottomrule
\end{tabularx}
\end{center}
\textbf{Reference approach.} Normalize both specs (resolve local refs),
diff operations/parameters/schemas against the chapter's compatibility
classes, and classify findings into error vs warning tiers. Stretch:
semantic-risk heuristics (format changes, pattern tightening,
\texttt{additionalProperties: false} introductions).

\subsubsection*{Capstone 9.4 --- Deprecation Campaign Plan \& Retirement Analytics [senior]}
\begin{center}\small
\begin{tabularx}{\textwidth}{l c X}
\toprule
\textbf{Criterion} & \textbf{Weight} & \textbf{Full marks when\ldots} \\
\midrule
Policy completeness & 25\% & timeline, SLAs, channels, and exceptions all documented. \\
Projection accuracy & 30\% & per-consumer trend analytics project a zero-crossing date and name the stragglers --- deterministically over the 6-month/12-consumer fixture. \\
Operational checklist & 25\% & go/no-go items are all checkable from telemetry (usage at zero, brownout results), not sentiment. \\
Communication plan & 20\% & outreach keyed to the contact list with brownout simulation evidence. \\
\bottomrule
\end{tabularx}
\end{center}
\textbf{Reference approach.} Parse the synthetic access log into
per-consumer weekly series; extrapolate linearly to a retirement date;
simulate brownouts; gate the kill date on the checklist. Stretch:
``what-if'' mode showing the date shift if the top two consumers migrate
next month (Chapter~\ref{ch:lifecycle}).

\subsubsection*{Capstone 9.5 --- Version Governance at Platform Scale [staff]}
\begin{center}\small
\begin{tabularx}{\textwidth}{l c X}
\toprule
\textbf{Criterion} & \textbf{Weight} & \textbf{Full marks when\ldots} \\
\midrule
Policy codifiability & 25\% & the policy is complete enough to lint mechanically --- no judgment calls left implicit. \\
Linter breadth & 25\% & $\geq$4 violation classes flagged on the seeded 40-internal/6-external registry. \\
Cost model & 25\% & carrying-cost report ranks versions reproducibly offline and identifies when the quadratic term bites. \\
Decision record & 25\% & RFC-style document defends the strategy mix against $\geq$2 alternatives, with rollback/exception governance. \\
\bottomrule
\end{tabularx}
\end{center}
\textbf{Reference approach.} Write the policy as lint rules first, prose
second; build the registry linter over a synthetic catalog; model carrying
cost as teams $\times$ versions $\times$ time; the decision record compares
URI/header/media-type mixes. Stretch: a ten-year portfolio simulation
showing the policy-change trigger point.

%% ============================================================================
\section{Chapter 10: Idempotency, Caching \& Concurrency}
\label{sec:sol-ch10}

\subsection{Self-Assessment Model Answers}

\emph{Numbering note: the chapter uses a plain 1--10 list; numbering below
matches.}

\begin{enumerate}[leftmargin=2.2em]
  \item GET: safe and idempotent. PUT: idempotent, not safe. DELETE:
    idempotent, not safe. POST: neither. PATCH with
    \texttt{\{"op":"set",...,"value":5\}}: idempotent (absolute write), not
    safe. PATCH with \texttt{\{"op":"inc",...\}}: neither --- $n$ retries
    add $n$ times \cite{rfc9110}.
  \item The claim is wrong: \texttt{no-cache} permits storage but mandates
    \emph{revalidation with the origin before every reuse}; it is the
    freshness guarantee, not a storage ban. The never-persist directive is
    \texttt{no-store} --- quoting the wrong one either leaks sensitive data
    or destroys cache efficiency.
  \item With \texttt{max-age=60} and storage at $t{=}0$: $t{=}30$ is a hit
    (fresh). $t{=}60$ is the boundary --- age $\geq$ max-age means stale,
    so revalidate (a 304 can still save the body). $t{=}61$: stale,
    revalidate. Verify all three against Listing 10.5.
  \item First request: processed, response stored under the key. Every
    subsequent request with the same key but a different payload: fingerprint
    mismatch $\to$ 409/422, \emph{no} re-execution. The batch job's logic
    bug (key reuse) is surfaced loudly instead of silently replaying the
    first transfer $n$ times --- exactly-once is preserved, and the client
    is forced to mint one key per logical operation.
  \item \emph{Store/replay:} a deterministic 422 will fail identically
    forever; replaying saves downstream load and gives the client a stable
    answer. \emph{Abandon:} stored errors can fossilize a transient
    validation bug and complicate TTL policy. Stripe's published design
    persists and replays the first result including client errors
    (validation is deterministic) but does not replay transient 5xx ---
    matching Listing 10.1's abandonment of the key row on downstream 5xx.
  \item Extend the cache key with \texttt{Accept-Language} and emit
    \texttt{Vary: Accept-Language}. Timeline of the bug without it: a
    French request primes the shared cache; an English request for the same
    URI is served the French representation --- a cross-user correctness
    leak caused by a missing header, not by application code
    \cite{rfc9110}.
  \item One table: \texttt{idempotency\_keys(key PK, request\_fingerprint,
    status, response\_snapshot, created\_at)} with
    \texttt{status $\in$ \{in\_flight, completed\}}. The business mutation
    and the key-row upsert commit in \emph{one} transaction; a crash rolls
    both back, leaving neither an orphan key nor a partial mutation. A
    concurrent retry seeing \texttt{in\_flight} waits or receives 409 ---
    the row is the serialization point.
  \item At a 95\% 412 rate, expected attempts per success $\approx$ 20, so
    optimistic throughput collapses to $\sim$5\% even with jitter; bounded
    retry only bounds the waste. Pessimistic row locks serialize correctly
    but head-of-line block every reader of the row. A per-SKU command queue
    serializes only the hot key while the rest of the catalog stays
    concurrent --- pick the queue (or locks if simplicity dominates), with
    the 20$\times$ waste figure as the quantitative justification.
  \item Tag every cached representation with surrogate keys of the entities
    it embeds (\texttt{part-123} on the product page, search cards, and
    recommendation slots); a price change on part 123 purges exactly
    \texttt{part-123} across all three surfaces --- no URL-pattern purges,
    no over-invalidation.
  \item \emph{Refuted.} Idempotency keys have a TTL (hours to days); after
    expiry a late retry re-executes and duplicates. Unique constraints are
    the timeless integrity backstop that knows nothing of TTLs; keys are a
    dedup/UX layer at the transport, constraints are the correctness layer
    in storage. Both, always.
\end{enumerate}

\subsection{Capstone Grading Rubrics}

\subsubsection*{Capstone 10.1 --- Idempotency-Key Client SDK [junior]}
\begin{center}\small
\begin{tabularx}{\textwidth}{l c X}
\toprule
\textbf{Criterion} & \textbf{Weight} & \textbf{Full marks when\ldots} \\
\midrule
Exactly-once effect & 35\% & up to 5 retries produce exactly one server-side mutation, proven against Listing 10.1's server. \\
Retry classification & 25\% & 422 never retried; 409 retried with backoff; the policy table is explicit. \\
Key discipline & 25\% & a second logical operation mints a fresh key and executes a second transfer. \\
Determinism & 15\% & tests pass offline, seeded, no wall-clock sleeps. \\
\bottomrule
\end{tabularx}
\end{center}
\textbf{Reference approach.} An \texttt{httpx} wrapper minting one UUID key
per logical operation, attaching \texttt{Idempotency-Key}, and classifying
responses by retryability. Stretch: persist in-flight operations so a
crashed client resumes with the same key.

\subsubsection*{Capstone 10.2 --- Conditional-GET Sync Agent [junior]}
\begin{center}\small
\begin{tabularx}{\textwidth}{l c X}
\toprule
\textbf{Criterion} & \textbf{Weight} & \textbf{Full marks when\ldots} \\
\midrule
Revalidation correctness & 35\% & unchanged polls return 304 and transfer zero representation bytes. \\
State fidelity & 25\% & local state equals the server representation after every poll. \\
Savings evidence & 25\% & bytes-saved report vs naive always-GET, reproducible from the scripted schedule. \\
Scheduling hygiene & 15\% & poll loop deterministic in tests (injected clock/sleeper). \\
\bottomrule
\end{tabularx}
\end{center}
\textbf{Reference approach.} Store the \texttt{ETag} beside each synced
representation; send \texttt{If-None-Match}; on 304 touch nothing, on 200
replace. Stretch: \texttt{Last-Modified}/\texttt{If-Modified-Since}
fallback when no ETag is present.

\subsubsection*{Capstone 10.3 --- Optimistic-Locking Order Service [mid]}
\begin{center}\small
\begin{tabularx}{\textwidth}{l c X}
\toprule
\textbf{Criterion} & \textbf{Weight} & \textbf{Full marks when\ldots} \\
\midrule
Conservation under storm & 30\% & 100-worker storm preserves quantity exactly --- no lost updates. \\
Status semantics & 25\% & blind writes get 428; stale writes get 412; correct retries succeed. \\
Client loop & 25\% & read--decide--write--retry loop with bounded attempts and jitter. \\
Reproducibility & 20\% & contention report identical across runs. \\
\bottomrule
\end{tabularx}
\end{center}
\textbf{Reference approach.} Extend Listing 10.3 so reserve/release are
conditional writes on a version token (\texttt{If-Match}); the client
re-reads on 412 before re-deciding. Stretch: adaptive escalation to a
serialized per-SKU queue past a 412-rate threshold, with the
throughput/latency trade-off quantified.

\subsubsection*{Capstone 10.4 --- Cache-Policy Workbench [senior]}
\begin{center}\small
\begin{tabularx}{\textwidth}{l c X}
\toprule
\textbf{Criterion} & \textbf{Weight} & \textbf{Full marks when\ldots} \\
\midrule
Evaluation correctness & 30\% & candidate policies ranked on origin load, staleness, and error resilience over the synthetic timeline; table bit-reproducible. \\
Outage resilience & 25\% & at least one policy serves zero failed requests during the scripted outage; the memo explains the mechanism (stale-while-revalidate / grace). \\
Safety classification & 25\% & account page correctly assigned \texttt{no-store}/\texttt{private} with justification. \\
Analysis depth & 20\% & trade-offs quantified, not adjectives. \\
\bottomrule
\end{tabularx}
\end{center}
\textbf{Reference approach.} Turn Listing 10.5 into a policy engine: a
scripted request timeline drives each candidate \texttt{Cache-Control}
policy; metrics are counted offline. Stretch: model a shared cache with a
\texttt{Vary} dimension and demonstrate a misconfiguration leaking across
user segments.

\subsubsection*{Capstone 10.5 --- Correctness Audit of a Retry-Capable Gateway [staff]}
\begin{center}\small
\begin{tabularx}{\textwidth}{l c X}
\toprule
\textbf{Criterion} & \textbf{Weight} & \textbf{Full marks when\ldots} \\
\midrule
Audit precision & 25\% & linter flags exactly the seeded violations --- zero false positives, zero misses. \\
Remediation rigor & 30\% & every remediation names its enforcing mechanism (unique constraint, fingerprint, fencing token). \\
Telemetry plan & 20\% & one alert per war-story failure mode, with thresholds. \\
Blast-radius analysis & 25\% & the plan states what happens when the idempotency-key store is down (fail open vs closed, per endpoint) and how it rolls back. \\
\bottomrule
\end{tabularx}
\end{center}
\textbf{Reference approach.} Enumerate every retry source (clients,
gateways, queues); classify each mutation endpoint by idempotency posture;
the linter encodes the posture rules; the report sequences remediations by
risk. Stretch: extend the audit across the async boundary (queue
redelivery plus consumer-side dedup windows, Chapter~\ref{ch:async_apis}).

%% ============================================================================
\section{Chapter 11: Async \& Event-Driven APIs}
\label{sec:sol-ch11}

\subsection{Self-Assessment Model Answers}

\begin{enumerate}[label=\textbf{E11.\arabic*.}, leftmargin=2.8em]
  \item Public stock-ticker widget: \textbf{SSE} --- one-way server push
    over plain HTTP. Two-player browser chess: \textbf{WebSockets} ---
    bidirectional, low latency. ERP notification on invoice approval:
    \textbf{webhook} --- server-to-server with retries and signatures.
    Mobile badge count on a metered connection: \textbf{polling with ETags}
    (or platform push) --- cheapest bytes, no standing connection. SOC
    wallboard behind a corporate proxy: \textbf{SSE} --- survives proxies
    that terminate WebSocket upgrades.
  \item Two explanations fit valid signatures with 40-minute-old timestamps:
    (1)~receiver (or sender) clock skew beyond the tolerance window;
    (2)~genuine delayed/replayed deliveries after a long outage. The
    one-line distinguisher: have the provider include its current timestamp
    (or per-attempt IDs) in the delivery --- or log the receiver's clock at
    rejection --- so skew vs delay is decidable from one log line.
  \item During rotation the deliverer emits \texttt{t=...,v1=sig\_new,
    v1=sig\_old} (two \texttt{v1} values); the receiver accepts if
    \emph{any} \texttt{v1} verifies against a known secret. The test proves
    holders of either secret accept, and a receiver with neither rejects.
  \item Per the WHATWG grammar: a line with no colon is treated as a field
    name with an empty value; \texttt{data:} with empty value appends an
    empty line (the event still dispatches); an \texttt{id} containing NUL
    causes the field to be ignored; \texttt{retry: abc} is ignored because
    the value is not all digits. Encode each as a parser test.
  \item Unmasked final text frame, 60-byte payload: \texttt{81 3C}. Masked
    final binary frame, 300-byte payload: \texttt{82 FE} (FIN+opcode 2;
    mask bit + 126 extended length), then \texttt{01 2C}, then the 4-byte
    masking key. Ping with payload ``hi'': \texttt{8A 02 68 69}. Byte two's
    top bit is the mask flag in each case.
  \item \texttt{0x48 XOR 0x37 = 0x7F} goes on the wire. An in-browser
    attacker cannot choose payload bytes that survive masking unchanged
    because the \emph{browser} generates a fresh masking key per frame and
    the attacker never observes it --- masking exists precisely to stop
    attacker-controlled bytes from confusing intermediaries.
  \item Delays $\min(2^{n-1}, 60)$ for $n=1..8$: $1+2+4+8+16+32+60+60 =
    183$~s worst case. $\pm$20\% jitter prevents synchronized reconnection
    herds after a shared failure (thundering herd). The suite seeds
    \texttt{random} so every delay sequence is reproducible.
  \item With 2{,}000 missed events against a 500-event ring buffer, resume
    is impossible: the server answers the stale \texttt{Last-Event-ID} with
    a \texttt{resync-required} control event; the client drops its state,
    refetches a full snapshot via REST, then resumes the stream fresh. UX:
    a ``reconnecting'' notice followed by a clean full refresh --- degraded,
    never corrupt.
  \item Runbook for ``dead-letter rate $>$ 10/hour'': (1)~DLQ grouped by
    endpoint and error class; (2)~last successful delivery per endpoint;
    (3)~recent deploys/secret rotations. Decision tree: poison payload
    (schema mismatch) $\to$ fix producer, then replay; consumer down $\to$
    wait for recovery, then replay; expired credentials $\to$ rotate, then
    replay; truly obsolete events $\to$ discard with a logged reason.
    Customer communication triggers when a contracted endpoint's freshness
    SLO is breached, not on the alert itself.
  \item The CI script loads old and new AsyncAPI documents, walks
    channels/messages/payload properties, and fails on any removed or
    retyped property --- a \emph{backward-compatibility} gate (the new
    provider must still emit everything old consumers parse)
    \cite{asyncapi}.
\end{enumerate}

\subsection{Capstone Grading Rubrics}

\subsubsection*{Capstone 11.1 --- Signed Webhook Receiver for Robot Alerts [junior]}
\begin{center}\small
\begin{tabularx}{\textwidth}{l c X}
\toprule
\textbf{Criterion} & \textbf{Weight} & \textbf{Full marks when\ldots} \\
\midrule
Verification correctness & 30\% & all 12 valid events processed exactly once under the \texttt{t=...,v1=...} scheme with 300\,s tolerance. \\
Hostile-input handling & 30\% & tampered, stale, wrong-secret, replayed, and missing-header variants each rejected with a 4xx problem body. \\
Idempotent duplicates & 25\% & duplicates return 200 \texttt{\{"status":"duplicate"\}} with no side effects. \\
Determinism & 15\% & fixed injected clock; offline tests. \\
\bottomrule
\end{tabularx}
\end{center}
\textbf{Reference approach.} Verify HMAC over \texttt{t.body} in constant
time, enforce tolerance, dedupe on event ID in a small store. Stretch:
two-secret rotation window and Prometheus-style counters.

\subsubsection*{Capstone 11.2 --- Webhook Delivery Service with DLQ and Replay [mid]}
\begin{center}\small
\begin{tabularx}{\textwidth}{l c X}
\toprule
\textbf{Criterion} & \textbf{Weight} & \textbf{Full marks when\ldots} \\
\midrule
Status classification & 25\% & permanent 422 gets exactly one attempt per event; transient failures retry. \\
Retry \& DLQ correctness & 30\% & bounded jittered backoff asserted with injected clock/sleeper (no test sleeps); always-down endpoint dead-letters everything. \\
Replay operator path & 25\% & \texttt{replay} restores delivery after recovery without duplicates. \\
Documentation \& ops & 20\% & one-page delivery-semantics doc; per-endpoint circuit breaker behavior stated. \\
\bottomrule
\end{tabularx}
\end{center}
\textbf{Reference approach.} Queue $\to$ deliverer with classification
table $\to$ dead-letter store $\to$ operator replay command; the circuit
breaker protects the fleet from a slow consumer. Stretch: SSRF guardrails
on endpoint registration and a per-endpoint metrics endpoint.

\subsubsection*{Capstone 11.3 --- Resumable SSE Fleet Dashboard Backend [mid]}
\begin{center}\small
\begin{tabularx}{\textwidth}{l c X}
\toprule
\textbf{Criterion} & \textbf{Weight} & \textbf{Full marks when\ldots} \\
\midrule
Resume correctness & 30\% & resume within the retention window misses zero events; outside it, a \texttt{resync-required} control event is emitted. \\
Proxy survival & 25\% & through a buffering proxy, events arrive within 1\,s of emission (the \texttt{X-Accel-Buffering} fix demonstrated). \\
Grammar robustness & 25\% & tests cover comments, multi-line data, unknown fields, and empty data. \\
Channel model & 20\% & multi-channel subscription with clean disconnect handling. \\
\bottomrule
\end{tabularx}
\end{center}
\textbf{Reference approach.} A 500-event ring buffer per channel keyed by
monotonic event IDs; \texttt{Last-Event-ID} lookup decides resume vs
resync; heartbeat comments keep intermediaries honest. Stretch: HTTP/2
serving, per-channel retry tuning, and an EventSource demo page.

\subsubsection*{Capstone 11.4 --- Bidirectional Robot Control Channel [senior]}
\begin{center}\small
\begin{tabularx}{\textwidth}{l c X}
\toprule
\textbf{Criterion} & \textbf{Weight} & \textbf{Full marks when\ldots} \\
\midrule
Exactly-once commands & 30\% & commands survive connection kills exactly-once under the seeded chaos harness. \\
Liveness discipline & 20\% & zombie connections reaped within two heartbeat intervals. \\
Graceful shutdown & 20\% & fleet-wide shutdown completes with zero unclean (1006) closes. \\
Auth \& negotiation & 15\% & subprotocol negotiation and first-message ticket auth enforced before any command. \\
Design justification & 15\% & written note defends WebSockets over SSE+REST for this workload. \\
\bottomrule
\end{tabularx}
\end{center}
\textbf{Reference approach.} Ticket-issued first message authenticates;
sequence numbers with server-side dedup give idempotent commands; ping/pong
heartbeats drive reaping; shutdown sends close frames and drains. Stretch:
permessage-deflate and a stubbed backplane proving $N$-node scalability.

\subsubsection*{Capstone 11.5 --- Event-Contract Governance Platform [staff]}
\begin{center}\small
\begin{tabularx}{\textwidth}{l c X}
\toprule
\textbf{Criterion} & \textbf{Weight} & \textbf{Full marks when\ldots} \\
\midrule
Compatibility gate & 30\% & diff gate passes the compatible change and fails the breaking retype, with diagnostics naming the property \cite{asyncapi}. \\
Shared library & 25\% & the consumer-obligation receiver library passes the chapter's full webhook suite unmodified. \\
Governance coverage & 25\% & one-pager covers schema evolution, secret rotation, DLQ ownership, and event retirement. \\
Operational adoption & 20\% & AsyncAPI docs generated in CI; gate wired with a documented exception path. \\
\bottomrule
\end{tabularx}
\end{center}
\textbf{Reference approach.} Treat the AsyncAPI documents as the contract
of record; the diff gate enforces consumer-driven compatibility; the shared
receiver library encodes verification/idempotency obligations so no team
re-implements them. Stretch: CloudEvents envelopes across all three
contracts and a generated Markdown catalog.

%% ============================================================================
\section{Chapter 12: GraphQL}
\label{sec:sol-ch12}

\subsection{Self-Assessment Model Answers}

\emph{Numbering note: the chapter uses one running 1--15 list split into
Recall (1--5), Application (6--10), and Deep-Dive (11--15); numbering below
matches.}

\begin{enumerate}[leftmargin=2.2em]
  \item The eight type-system kinds with catalog examples: Scalar
    (\texttt{SKU: String}), Object (\texttt{Product}), Interface
    (\texttt{Node}), Union (\texttt{SearchResult = Product | Category}),
    Enum (\texttt{StockStatus}), InputObject (\texttt{ProductFilter}),
    List (\texttt{[Product]}), NonNull (\texttt{id: ID!})
    \cite{graphqlspec}.
  \item \texttt{[Review]} (all nullable): an element error nulls that
    element only. \texttt{[Review!]}: an element error nulls the whole
    list. \texttt{[Review]!}: a null list propagates to the parent field.
    \texttt{[Review!]!}: any element error nulls the list \emph{and}
    propagates to the parent --- potentially \texttt{data: null} at the
    root. Nullability \emph{is} error-containment policy.
  \item Phases: (1)~parse/validate --- failures are request errors, no
    execution, no \texttt{data} key; (2)~execution --- per-field failures
    land in \texttt{errors[]} alongside partial \texttt{data};
    (3)~response assembly --- transport-level failures (the channel the
    client must also handle distinctly).
  \item Variables are \emph{values} substituted after parsing; they can
    never alter the query document's structure, so GraphQL-injection in the
    SQL sense is impossible through variables. What remains injectable: the
    \emph{values themselves} when a resolver concatenates them into SQL or
    shell --- injection moved one layer down, not away.
  \item APQ handshake: the client sends only the query hash; on a cache miss
    the server answers \texttt{PersistedQueryNotFound}; the client then
    resends hash + full text, and subsequent requests go hash-only. The
    miss is HTTP 200 because it is a GraphQL-protocol outcome inside a
    valid response envelope, not a transport failure.
  \item Without batching: 1 query for categories plus one per category ---
    $1 + C$ queries. With the DataLoader, the product fetches collapse into
    one batched read: 2 queries total, independent of $C$; the
    instrumentation must show exactly this.
  \item A query with \texttt{first: 1} on a connection whose resolver runs
    a full search-index scan is cheap to the estimator (charged 1) and
    expensive to execute. Fix: per-field base costs plus argument
    multipliers, not multipliers alone. Collateral damage: legitimate
    small-page consumers of genuinely expensive fields get throttled ---
    the cost model must be calibrated against measured resolver latency.
  \item Backward pagination runs the window query in reverse order
    (\texttt{last}/\texttt{before}) and re-reverses the page.
    \texttt{hasPreviousPage} requires probing for rows \emph{before} the
    returned window (a limit+1 fetch in the reverse direction) --- the
    mirror of forward pagination's \texttt{hasNextPage} probe, with cursor
    direction semantics that are easy to invert.
  \item Add \texttt{priceCents}, mark \texttt{unitPriceCents}
    \texttt{@deprecated}, keep resolving both. This exercises continuous
    evolution: additive change plus machine-readable deprecation, never
    in-place removal --- the GraphQL analogue of expand-and-contract
    \cite{graphqlspec}.
  \item If \texttt{errors == [COST\_DATA\_UNAVAILABLE]} and \texttt{data}
    is present, render the card with a fallback cost block (partial data is
    a feature). On \texttt{GraphQLTransportError} (no \texttt{data} at
    all), retry the operation with backoff --- the two channels have
    opposite handling.
  \item The spec can order mutation root fields because they arrive in one
    document on one connection; cross-document ordering would require
    global coordination no transport-neutral spec can mandate. Client
    implication: never depend on ordering between separate operations ---
    sequence dependent mutations in one document or await each response.
  \item \emph{For} \texttt{edges: [ProductEdge!]!}: simpler client
    invariants, no null checks. \emph{Against}: null propagation couples
    error containment to schema design --- one bad edge nulls the whole
    connection, turning a single-row fault into a page-level outage.
    Validating telemetry: per-path field-error rates before and after
    tightening nullability.
  \item Under a full persisted-query allow-list, GraphQL still buys the
    typed schema, developer tooling, deprecation workflow, and
    partial-field selection \emph{within} approved documents. What it
    costs: the open query language's flexibility --- you have built closed
    RPC with a sophisticated type system, which may be exactly right for a
    public API under abuse pressure.
  \item Price each field as $\text{base cost} \times \text{argument
    multiplier}$, with base costs sourced from measured resolver latency
    (search-index hit vs memory hit differ by orders of magnitude). In a
    code-first schema the metadata lives in per-field extensions/directives
    consumed by the static analyzer.
  \item DataLoader batching transfers to cross-service entity resolution
    (\texttt{\_entities} batching per subgraph per request), but three
    things break or degrade: cross-request cache invalidation is now
    distributed; partial subgraph failure mid-batch must degrade per-field;
    and the batching window adds latency to the critical path --- the N+1
    became network fan-out, which is better but not free
    \cite{kleppmann2017ddia}.
\end{enumerate}

\subsection{Capstone Grading Rubrics}

\emph{Numbering note: the chapter labels its capstones ``Capstone 1--5'';
they are cited here as C12.1--C12.5.}

\subsubsection*{Capstone 12.1 --- Catalog Schema from Scratch [junior]}
\begin{center}\small
\begin{tabularx}{\textwidth}{l c X}
\toprule
\textbf{Criterion} & \textbf{Weight} & \textbf{Full marks when\ldots} \\
\midrule
Schema completeness & 30\% & types, enum, input, union mutation result, and a cursor connection all present; generated SDL shows them. \\
Pagination correctness & 25\% & full cursor walk returns every product exactly once. \\
Error discipline & 25\% & forged cursor $\to$ field error with machine-readable \texttt{extensions.code}; out-of-range rating $\to$ typed union payload with \texttt{errors} absent. \\
Repository quality & 20\% & own seeded repository; deterministic offline tests. \\
\bottomrule
\end{tabularx}
\end{center}
\textbf{Reference approach.} Re-implement schema-first from the SDL
outward; keep mutation results as union types so domain errors are typed
payloads, not exceptions. Stretch: \texttt{@deprecated} verified in
introspected SDL and backward pagination.

\subsubsection*{Capstone 12.2 --- N+1 Hunter [mid]}
\begin{center}\small
\begin{tabularx}{\textwidth}{l c X}
\toprule
\textbf{Criterion} & \textbf{Weight} & \textbf{Full marks when\ldots} \\
\midrule
Measured elimination & 35\% & store-query counts independent of page size, verified at two sizes. \\
Dedup correctness & 25\% & a test proves duplicate keys collapse into one batch read. \\
Loader lifecycle & 25\% & regression test fails with a module-level (stale-cache) loader and passes with per-request loaders. \\
Reporting & 15\% & before/after query counts tabulated per operation. \\
\bottomrule
\end{tabularx}
\end{center}
\textbf{Reference approach.} Instrument the store with a counting wrapper;
find every resolver doing per-item fetches; introduce request-scoped
DataLoaders. Stretch: per-tick batch-size metrics in a response
\texttt{extensions} block.

\subsubsection*{Capstone 12.3 --- Cost-Control Gateway [mid]}
\begin{center}\small
\begin{tabularx}{\textwidth}{l c X}
\toprule
\textbf{Criterion} & \textbf{Weight} & \textbf{Full marks when\ldots} \\
\midrule
Gate accuracy & 30\% & every corpus operation accepted/rejected exactly as labeled. \\
Cost model soundness & 25\% & complexity monotonically increasing in \texttt{first} (property-tested). \\
Adversarial robustness & 25\% & fragment cycles rejected with a distinct code --- never hang, never crash. \\
Enforcement & 20\% & per-operation timeout proven to fire. \\
\bottomrule
\end{tabularx}
\end{center}
\textbf{Reference approach.} Schema-agnostic AST walker computing depth and
argument-weighted complexity before execution; reject over budget with a
machine-readable code. Stretch: per-field cost weights from measured
resolver latency and complexity reported in response extensions.

\subsubsection*{Capstone 12.4 --- APQ Production Gateway [senior]}
\begin{center}\small
\begin{tabularx}{\textwidth}{l c X}
\toprule
\textbf{Criterion} & \textbf{Weight} & \textbf{Full marks when\ldots} \\
\midrule
Allow-list strictness & 30\% & strict mode rejects unregistered documents with a distinct code and zero resolver executions (proven). \\
Handshake fidelity & 25\% & client completes miss $\to$ register $\to$ hash-only, sending full text exactly once. \\
Cost-based limiting & 25\% & complexity-based rate limiting admits small queries while rejecting god-query bursts at equal request count. \\
Introspection policy & 20\% & introspection disabled/governed per environment with tests. \\
\bottomrule
\end{tabularx}
\end{center}
\textbf{Reference approach.} APQ cache + persisted-document registry ahead
of validation; complexity scoring feeds the rate limiter
(Chapter~\ref{ch:rate_limiting}). Stretch: serve registered reads over GET
for CDN cacheability with a documented cache-key design.

\subsubsection*{Capstone 12.5 --- Federated Robotics Graph [staff]}
\begin{center}\small
\begin{tabularx}{\textwidth}{l c X}
\toprule
\textbf{Criterion} & \textbf{Weight} & \textbf{Full marks when\ldots} \\
\midrule
Cross-subgraph correctness & 30\% & one operation spanning catalog + reputation returns correctly stitched data. \\
Batched resolution & 25\% & cross-service entity resolution batched, counts asserted. \\
Failure isolation & 25\% & subgraph outage degrades only its fields (fault-injection test), never the whole response. \\
Scope honesty \& ops & 20\% & written scope note lists unimplemented federation features; gateway has metrics and a documented rollback. \\
\bottomrule
\end{tabularx}
\end{center}
\textbf{Reference approach.} Two subgraph services share the
\texttt{Product} entity; the hand-rolled gateway plans queries, fans out,
batches \texttt{\_entities} calls, and merges with per-field error
containment. Stretch: a subscription carried through the gateway and
per-subgraph cost budgets.

%% ============================================================================
\section{Chapter 13: gRPC \& High-Performance RPC}
\label{sec:sol-ch13}

\subsection{Self-Assessment Model Answers}

\emph{Numbering note: the chapter uses a plain 1--10 list; numbering below
matches.}

\begin{enumerate}[leftmargin=2.2em]
  \item The 5-byte prefix \texttt{00 00 00 00 0A} decodes as compression
    flag 0 (uncompressed) and message length 10 octets; the ten protobuf
    bytes are the message on stream 3. If the stream ends after six payload
    bytes, the message is truncated --- the receiver must fail the RPC
    (INTERNAL / stream error), never deliver a partial message
    \cite{grpcdocs}.
  \item 300 bytes $=$ \texttt{0x12C}, so the prefix is \texttt{00 00 00 01
    2C}. The 32-bit length field can express up to $2^{32}-1$ bytes
    ($\approx$4\,GiB), but grpc-python enforces
    \texttt{grpc.max\_receive\_message\_length} at 4\,MiB by default ---
    the wire format is not the binding constraint.
  \item (a)~Firmware-update progress push: \textbf{server-streaming} --- one
    request, many progress messages. (b)~Uploading a day of buffered sensor
    logs: \textbf{client-streaming} --- many chunks, one summary response.
    (c)~Trader-style command/ack console: \textbf{bidirectional} ---
    independent streams both ways. (d)~One robot's current pose:
    \textbf{unary} --- request/response, simplest possible surface.
  \item The downstream server sees \texttt{grpc-timeout} of $500 - 40 =
    460$\,ms, encoded \texttt{"460m"}. If the proxy \emph{adds} 50\,ms, the
    downstream keeps working after the original client has already
    abandoned the call --- wasted work and misleading success signals;
    deadlines must only shrink as they propagate, never grow
    \cite{grpcdocs}.
  \item \texttt{UploadTelemetry} mutates server state: retrying
    DEADLINE\_EXCEEDED risks applying the upload twice (the first attempt
    may have succeeded after the deadline). \texttt{GetRobotStatus} is a
    read with no side effects, so retrying it is harmless --- retry safety
    is a property of the method's semantics, not the status code
    (Chapter~\ref{ch:idempotency}).
  \item Two-phase migration: (1)~add \texttt{battery\_millivolts} (new tag)
    and dual-publish both fields while consumers migrate; (2)~after usage
    hits zero, mark tag 3 and the old name \texttt{reserved} and stop
    publishing. Never retype tag 3 in place --- wire-type/type changes are
    the incompatibility class of Chapter~\ref{ch:serialization}
    \cite{protobufdocs}.
  \item The interceptor reads the token from metadata, checks a per-token
    bucket, and on exhaustion aborts with RESOURCE\_EXHAUSTED plus a
    RetryInfo-style detail packed in a \texttt{-bin} trailer (binary
    metadata). Streaming handlers differ from unary: the limiter must
    decide whether to charge per call or per message, and must not break
    long-lived streams mid-flight --- four handler shapes, four charging
    decisions.
  \item Hypothesis: the server's connection-recycling option (e.g.\
    \texttt{max\_connection\_age} at 300\,s) forces reconnects every five
    minutes, and reconnect storms surface as UNAVAILABLE spikes. Confirm by
    changing \texttt{grpc.max\_connection\_age\_ms} to a different value ---
    the spike interval follows the setting, proving causation
    \cite{grpcdocs}.
  \item grpc-web: full contract fidelity, browser-safe, but needs a proxy
    and has limited streaming. JSON transcoding: annotation-driven, cheap
    ops, weakest for streams. Bespoke WebSocket shim: maximal flexibility,
    maximal contract drift. For a browser dashboard over
    \texttt{StreamTelemetry}, pick grpc-web when streaming fidelity
    matters, transcoding when the surface is mostly unary --- justify on
    contract fidelity and ops cost, not fashion.
  \item \texttt{grpc-status} lives in trailers because the definitive
    outcome of a streaming RPC is knowable only after the last message; a
    \emph{trailers-only} response is HEADERS carrying the status with no
    DATA --- an error before any message. Trailers let intermediaries
    terminate failures without parsing bodies; conversely gRPC traffic is
    opaque to HTTP caches, and browsers need translation (grpc-web) because
    they cannot expose trailers mid-fetch \cite{rfc9113, grpcdocs}.
\end{enumerate}

\subsection{Capstone Grading Rubrics}

\emph{Numbering note: the chapter labels its capstones ``Capstone 1--5'';
they are cited here as C13.1--C13.5.}

\subsubsection*{Capstone 13.1 --- Unary Status Service with Deadline Discipline [junior]}
\begin{center}\small
\begin{tabularx}{\textwidth}{l c X}
\toprule
\textbf{Criterion} & \textbf{Weight} & \textbf{Full marks when\ldots} \\
\midrule
Path coverage & 30\% & pytest covers happy path, unknown robot, and deadline-exceeded cases. \\
Deadline discipline & 30\% & no call site without \texttt{timeout=}; no handler without a remaining-budget check. \\
Service shape & 20\% & grpcurl/descriptor dump demonstrates the service contract. \\
Error mapping & 20\% & domain errors map to correct status codes, asserted via \texttt{.code()}. \\
\bottomrule
\end{tabularx}
\end{center}
\textbf{Reference approach.} One proto, one servicer, one test module;
handlers check \texttt{context.time\_remaining()} before work. Stretch:
keepalive configuration with channel-state logging and a latency-measuring
client interceptor.

\subsubsection*{Capstone 13.2 --- Server-Streaming Telemetry Monitor [mid]}
\begin{center}\small
\begin{tabularx}{\textwidth}{l c X}
\toprule
\textbf{Criterion} & \textbf{Weight} & \textbf{Full marks when\ldots} \\
\midrule
Cancellation correctness & 30\% & after \texttt{cancel()} the server handler exits within one period (proven). \\
Flow-control bounds & 25\% & server never exceeds \texttt{max\_frames}; client drop counter observable. \\
Deadline enforcement & 25\% & stream deadline enforced and tested. \\
CLI quality & 20\% & bounded rendering; clean shutdown; readable output. \\
\bottomrule
\end{tabularx}
\end{center}
\textbf{Reference approach.} Server loop checks cancellation each period;
client renders into a bounded buffer with an explicit drop counter.
Stretch: heartbeat frames and per-frame stalled-stream detection.

\subsubsection*{Capstone 13.3 --- Bulk Ingestion with Rich Validation Errors [mid]}
\begin{center}\small
\begin{tabularx}{\textwidth}{l c X}
\toprule
\textbf{Criterion} & \textbf{Weight} & \textbf{Full marks when\ldots} \\
\midrule
Error richness & 30\% & mid-stream abort yields INVALID\_ARGUMENT with decoded per-field violations from the \texttt{-bin} trailer and an end-to-end correlation ID. \\
Resume correctness & 25\% & fix-and-resume accepts exactly the valid 95\%. \\
Assertion discipline & 25\% & all errors asserted via \texttt{.code()}, never message strings. \\
Interceptor design & 20\% & error mapping lives in the interceptor, not the servicer. \\
\bottomrule
\end{tabularx}
\end{center}
\textbf{Reference approach.} Client streams rows; the server validates per
message and aborts with packed details; the client decodes the trailer,
fixes the rows, and resumes from the acknowledged offset. Stretch: port to
\texttt{google.rpc.Status} and show wire equivalence.

\subsubsection*{Capstone 13.4 --- Bidirectional Control Plane with Interceptor Stack [senior]}
\begin{center}\small
\begin{tabularx}{\textwidth}{l c X}
\toprule
\textbf{Criterion} & \textbf{Weight} & \textbf{Full marks when\ldots} \\
\midrule
Exactly-once effects & 30\% & 500-command run has exactly-once effects despite forced retries, with ack order preserved (sequence-number idempotency). \\
Interceptor separation & 25\% & auth/logging/metrics interceptors emit per-cardinality counters without touching business code; bad tokens fail before servicer code. \\
Deadline propagation & 25\% & deadline test shows the downstream hop observing a smaller remaining budget. \\
Flow control & 20\% & windowed flow control prevents unbounded buffering under a slow consumer. \\
\bottomrule
\end{tabularx}
\end{center}
\textbf{Reference approach.} Bidi stream with per-command sequence numbers
and a server-side dedup table; interceptors cross-cut auth, logging,
metrics on both peers. Stretch: hedged unary reads with a measured p99
delta vs the retry baseline under deterministic latency injection
(Chapter~\ref{ch:microservices}).

\subsubsection*{Capstone 13.5 --- Edge Gateway: Transcoding + Load-Balancing Postmortem [staff]}
\begin{center}\small
\begin{tabularx}{\textwidth}{l c X}
\toprule
\textbf{Criterion} & \textbf{Weight} & \textbf{Full marks when\ldots} \\
\midrule
Reproduction & 25\% & pinned-phase skew $\geq$10$\times$ reproduced deterministically. \\
Fix efficacy & 30\% & each of the three fixes (max\_connection\_age, round\_robin, least-loaded) reduces the max/median RPC ratio below 1.5$\times$. \\
Transcoding correctness & 20\% & JSON edge returns correct HTTP codes for NOT\_FOUND/DEADLINE\_EXCEEDED/INVALID\_ARGUMENT with details intact. \\
Postmortem quality & 25\% & identifies connections-vs-RPCs as root cause, names the per-pod SLI, and includes detection and rollback notes; \texttt{pytest -q} green via one PowerShell driver. \\
\bottomrule
\end{tabularx}
\end{center}
\textbf{Reference approach.} Recreate the war story: one L4-balanced
connection per client pins load; fix at the server (connection aging), the
client (round\_robin), and the policy layer (least-loaded); write the
incident report with the measured evidence. Stretch: DNS-based resolver
config and \texttt{grpc.health.v1} gating the shim's routing on SERVING.

%% ============================================================================
\section{Chapter 14: Rate Limiting \& Throttling}
\label{sec:sol-ch14}

\subsection{Self-Assessment Model Answers}

\begin{enumerate}[label=\textbf{E14.\arabic*.}, leftmargin=2.8em]
  \item The boundary test shows the sliding-window counter admits fewer than
    $2N$ requests across a window edge (unlike fixed window, which admits
    exactly $2N$ in an infinitesimal span). It still over-admits relative to
    the exact sliding log when the previous window's traffic clustered at
    its end and the query happens mid-window: the linear weight then
    under-credits the true count. Bounding proposition: the estimate's error
    is at most the previous window's count, so admissions stay below $2N$.
  \item With correct parameterization (emission interval $T = 1/r$,
    tolerance $\tau = B/r$), no divergent $(r, B)$ pair exists --- the
    strengthened proof is by induction over the event stream: bucket level
    and $(\tau - (\mathrm{TAT} - t))/T$ are invariantly equal, including
    across multi-second idle gaps (both refill fully). The 5{,}000-event
    seeded sweep is the empirical confirmation of the equivalence theorem.
  \item A sliding log has no single reset instant; the honest, monotonic
    choice is to report the seconds until the \emph{oldest request still
    blocking capacity} expires --- i.e.\ the delay until at least one more
    request would be admitted. This decreases monotonically and never
    promises a cliff-edge reset that does not exist
    \cite{ietf_ratelimit_headers}.
  \item Three concrete failures from swapping the injected \texttt{Clock}
    for \texttt{time.monotonic}: (1)~window-rollover unit tests become
    real-time sleeps (slow, flaky); (2)~seeded simulation/replay harnesses
    (AIMD tuning, incident replay) become impossible --- they need virtual
    time; (3)~boundary-condition tests (burst exactly at a window edge)
    become nondeterministic, so regressions pass or fail by timing luck.
  \item Instrumented check-then-increment admits overshoot that grows with
    thread count: the TOCTOU window between read and write widens under
    contention, so 50-thread storms show systematically more over-admissions
    than 5-thread storms. A Lua script executes read-check-increment
    atomically \emph{inside} Redis, so the entire race class is eliminated
    by construction --- not probabilistically reduced.
  \item Keys: \texttt{quota:\{tenant\}:YYYY-MM} (INCR with month-end EXPIRE)
    and \texttt{rate:\{tenant\}} (token bucket). One Lua script checking
    both atomically is preferred; two scripts invite a quota/rate ordering
    hazard. \texttt{GET /v1/usage} returns
    \texttt{\{quota\_limit, quota\_used, quota\_remaining, resets\_at,
    rate\_limit, rate\_remaining\}}. A client at 99/100 with an empty
    per-second bucket gets 429 with \texttt{Retry-After} from the rate
    layer while the response body/headers still show the quota layer's
    remaining 1 --- the two layers must be distinguishable.
  \item From $L_0=100$ with $a{=}5$ rps per 1\,s evaluation: first contact
    with 300 rps capacity at $(300-100)/5 = 40$\,s. Steady state is a
    sawtooth: reach 300, multiplicative cut to $\beta \cdot 300 = 150$,
    climb back --- amplitude 150 rps peak-to-trough. $\beta < 1$ is
    safety-critical because multiplicative decrease sheds load
    \emph{proportionally} at any level (fast relief when far over
    capacity), while the additive term only probes slowly upward.
  \item A good capture documents: presence/absence of \texttt{Retry-After},
    the \texttt{RateLimit-Limit}/\texttt{-Remaining}/\texttt{-Reset} triple
    and \texttt{RateLimit-Policy}, and critically whether \texttt{Reset} is
    delay-seconds or epoch-seconds (a historical interop split)
    \cite{ietf_ratelimit_headers}. Contract comparison: a provider better
    than the chapter's middleware publishes both header families plus a
    problem body; worse is 429 with no headers at all --- clients must
    guess the retry instant.
\end{enumerate}

\subsection{Capstone Grading Rubrics}

\subsubsection*{Capstone 14.1 --- Prove the Algorithms [junior]}
\begin{center}\small
\begin{tabularx}{\textwidth}{l c X}
\toprule
\textbf{Criterion} & \textbf{Weight} & \textbf{Full marks when\ldots} \\
\midrule
Reimplementation fidelity & 35\% & all 14 ported tests pass on the from-definition reimplementation of the five algorithms. \\
Bug-detection mapping & 30\% & each seeded bug is caught by at least one existing test, with the bug$\to$test mapping named. \\
Extension proof & 20\% & the new non-integer-window ($W{=}1.5$) burst test proves 20 admissions in a sub-window interval. \\
Determinism & 15\% & seeded, offline, no wall-clock assertions. \\
\bottomrule
\end{tabularx}
\end{center}
\textbf{Reference approach.} Implement each algorithm from its mathematical
definition against an injected clock; port the proof sheet first, then
break the implementations deliberately. Stretch: property-test the sliding
log's exactness invariant over random seeded timelines.

\subsubsection*{Capstone 14.2 --- Header-Complete Middleware [mid]}
\begin{center}\small
\begin{tabularx}{\textwidth}{l c X}
\toprule
\textbf{Criterion} & \textbf{Weight} & \textbf{Full marks when\ldots} \\
\midrule
Contract completeness & 30\% & 429 bodies validate against the declared RFC~9457 schema; the full header family emitted \cite{rfc9457, ietf_ratelimit_headers}. \\
Degradation honesty & 30\% & fail-open episodes emit \texttt{X-NR-Limiter-Degraded} and no \texttt{RateLimit-*} fields; fail-closed returns 503 + \texttt{Retry-After}. \\
Drop-in compatibility & 25\% & all existing middleware tests pass unmodified against the Redis wiring. \\
Determinism & 15\% & suite deterministic and offline. \\
\bottomrule
\end{tabularx}
\end{center}
\textbf{Reference approach.} Package the ASGI middleware with a Redis-backed
limiter behind the same interface; make degradation modes explicit in the
response contract. Stretch: per-plan policies selected by API-key tier.

\subsubsection*{Capstone 14.3 --- Cell Architecture Simulator [mid]}
\begin{center}\small
\begin{tabularx}{\textwidth}{l c X}
\toprule
\textbf{Criterion} & \textbf{Weight} & \textbf{Full marks when\ldots} \\
\midrule
Topology bounds & 35\% & global topology never exceeds the limit; no-sync overshoot matches the $R \times$ prediction; sync topology respects the published bound at every interval $\Delta$. \\
Reproducibility & 25\% & all results regenerate from one seed. \\
Modeling fidelity & 25\% & the simulator's event model matches the chapter's cell semantics (local counts + sync). \\
Report quality & 15\% & over-admission quantified per topology with a recommendation. \\
\bottomrule
\end{tabularx}
\end{center}
\textbf{Reference approach.} Discrete-event simulation over cells with
local counters and periodic gossip/sync; sweep topology and $\Delta$;
tabulate overshoot. Stretch: inject a cell failure and show consistent-hash
remapping confines counter loss to $1/C$ of keys.

\subsubsection*{Capstone 14.4 --- Adaptive Limits with AIMD [senior]}
\begin{center}\small
\begin{tabularx}{\textwidth}{l c X}
\toprule
\textbf{Criterion} & \textbf{Weight} & \textbf{Full marks when\ldots} \\
\midrule
Control quality & 30\% & tuned controller never exceeds the moving capacity signal by $>$10\% for more than one evaluation interval; the contractual floor is never violated. \\
Tuning evidence & 30\% & $(a, \beta)$ grid sweep and the identified ``knee'' reproducible from seed. \\
Analysis & 25\% & tuning memo quantifies the multiplicative-decrease asymmetry argument. \\
Determinism & 15\% & no wall-clock dependence; capacity signal replayed from fixture. \\
\bottomrule
\end{tabularx}
\end{center}
\textbf{Reference approach.} AIMD controller per tenant evaluated on a
fixed cadence against a replayed capacity signal; sweep parameters; defend
the chosen $(a,\beta)$ with the oscillation data. Stretch: per-tenant
fairness dividing recovered headroom proportional to recent demand.

\subsubsection*{Capstone 14.5 --- Rate-Limiting Architecture for Northwind at 100k rps [staff]}
\begin{center}\small
\begin{tabularx}{\textwidth}{l c X}
\toprule
\textbf{Criterion} & \textbf{Weight} & \textbf{Full marks when\ldots} \\
\midrule
Coverage \& derivation & 25\% & every endpoint has a tier, topology, failure policy, and a numerically derived limit --- no defaults. \\
Bounds \& monitoring & 20\% & over-admission bounds stated and monitored per approximate tier. \\
Runbooks & 25\% & limiter-cell outage, abusive-tenant drill, and AIMD brownout runbooks each include detection and rollback steps. \\
Cost realism & 15\% & cost model shows limiter infrastructure $<$2\% of serving cost. \\
Red-team honesty & 15\% & $\geq$2 residual weaknesses named without spin. \\
\bottomrule
\end{tabularx}
\end{center}
\textbf{Reference approach.} Tier the surface (strict/exact vs
approximate), assign topologies per tier (centralized, cell, AIMD),
derive limits from capacity data, and write the operational runbooks as
part of the design --- not after it. Stretch: GraphQL cost-scoring and gRPC
interceptor integration so all three protocol surfaces share one budget.

%% ============================================================================
\section{Chapter 15: Microservices \& Resilience}
\label{sec:sol-ch15}

\subsection{Self-Assessment Model Answers}

\emph{Numbering note: the chapter uses a plain 1--8 list; numbering below
matches.}

\begin{enumerate}[leftmargin=2.2em]
  \item Budgets must shrink down the chain with headroom for each hop's own
    work: ledger $\le 300$\,ms (its measured p99.9), payments $\le 600$\,ms,
    orders $\le 1{,}000$\,ms, edge $\le 1{,}300$--$1{,}500$\,ms, satisfying
    the rule that each hop's timeout is strictly smaller than the caller's
    remaining budget. Validate by measuring per-hop p99/p99.9 under load
    and by deadline-propagation tests showing downstream hops abandon work
    the client already gave up on.
  \item Unbudgeted, each of the 4 hops retries twice: attempts multiply as
    $3^4 = 81$ downstream requests per original request at the deepest hop
    (total $3+9+27+81=120$). With a one-token retry budget at the two
    middle hops: $3 \times 1 \times 1 \times 3 = 9$. The budget at the
    \emph{earliest} hop matters most, because amplification compounds ---
    one retry near the client re-executes everything downstream.
  \item Defensible set: error-rate threshold $\theta = 50\%$ over a rolling
    window $W = 30$\,s; minimum call volume $N = 100$ (about 0.05\,s of
    traffic at 2{,}000 rps --- the window is always statistically
    meaningful); open duration $d = 30$\,s; half-open admits $p$ probes
    paced to stay under the 20 rps recovery allowance (e.g.\ 1 probe/s);
    close after $s = 5$ consecutive successes. Every number traces to the
    dependency's stated success rate and recovery budget
    \cite{nygard2018releaseit}.
  \item Add a slow-call dimension: a call whose duration exceeds the
    threshold counts as a failure in the window even on success. The
    manual-clock test injects five slow successes and asserts the window
    error rate reaches $\theta$ and the breaker opens --- slow failures are
    exactly the ones that exhaust pools while looking ``healthy''.
  \item Hedging after the measured p95 fires the second call on the $q$
    tail; the measured total call rate should converge to $1 + q$ within
    derived tolerance as the injected tail probability $q$ varies. Hedging
    is safe only when the operation is idempotent (or key-deduplicated)
    \emph{and} the dependency has capacity headroom for the extra $q$
    fraction of load.
  \item Per-aggregate ordering: outbox rows carry \texttt{(aggregate\_id,
    seq)}; the relay publishes in seq order per aggregate; the consumer
    tracks the last applied seq per aggregate and buffers/refetches on a
    gap. Crash-mid-batch replays then cannot interleave with new events ---
    the consumer enforces monotonic application per aggregate regardless of
    publisher timing.
  \item In the choreographed version, each service subscribes to domain
    events and emits its own (\texttt{inventory.reserved} $\to$ payments
    consumes, etc.); the compensation graph no longer lives in one
    orchestrator --- it is \emph{distributed across the subscriptions} of
    each participant. The catching test: fail one step and assert the
    compensating event is consumed and state reverted; a missing
    compensation subscription fails the test loudly.
  \item mTLS authenticates services, not outcomes; it does nothing for:
    (1)~object-level authorization (BOLA) --- the chapter's
    policy/authorization pattern; (2)~semantic abuse over a legitimate
    identity (injection, mass assignment) --- input validation and contract
    enforcement; (3)~business-flow abuse from a credentialed caller ---
    rate limiting and anomaly detection (Chapter~\ref{ch:rate_limiting}).
\end{enumerate}

\subsection{Capstone Grading Rubrics}

\subsubsection*{Capstone 15.1 --- Breaker Under a Microscope [junior]}
\begin{center}\small
\begin{tabularx}{\textwidth}{l c X}
\toprule
\textbf{Criterion} & \textbf{Weight} & \textbf{Full marks when\ldots} \\
\midrule
Reimplementation fidelity & 35\% & all existing breaker tests pass on the from-definition reimplementation. \\
Bug hunting & 30\% & each seeded bug (no min-call floor, no window eviction, unlimited probes) is caught by a named existing test. \\
New-test quality & 25\% & the probe-threshold test ($s{=}3$, success-success-failure) passes on the correct breaker and fails against the unlimited-probe bug. \\
Determinism & 10\% & manual clock only; no sleeps. \\
\bottomrule
\end{tabularx}
\end{center}
\textbf{Reference approach.} Implement the closed/open/half-open state
machine from the formal definition over an injected clock; run the existing
suite; then mutate. Stretch: the slow-call failure dimension with a
manual-clock proof.

\subsubsection*{Capstone 15.2 --- Latency Chaos and the Price of Hedging [mid]}
\begin{center}\small
\begin{tabularx}{\textwidth}{l c X}
\toprule
\textbf{Criterion} & \textbf{Weight} & \textbf{Full marks when\ldots} \\
\midrule
Benefit evidence & 30\% & hedged p99 materially below unhedged p99 on the same seed under bimodal latency faults. \\
Cost accounting & 30\% & measured hedge rate equals $1 + q$ within the derived tolerance --- the price is quantified, not asserted. \\
Breaker interaction & 25\% & hedging is disabled while the breaker is open, proven by test. \\
Determinism & 15\% & fully seeded/offline; no wall-clock assertions. \\
\bottomrule
\end{tabularx}
\end{center}
\textbf{Reference approach.} Chaos stub with bimodal latency; hedge timer
at measured p95; count issued vs completed calls. Stretch: loser-request
cancellation with measured downstream savings, and hedges spending
retry-budget tokens.

\subsubsection*{Capstone 15.3 --- Production Outbox: Retry, DLQ, Ordering [mid]}
\begin{center}\small
\begin{tabularx}{\textwidth}{l c X}
\toprule
\textbf{Criterion} & \textbf{Weight} & \textbf{Full marks when\ldots} \\
\midrule
Retry discipline & 25\% & publish-failure retries respect the budget cap (counting-publisher proof). \\
Dead-lettering & 25\% & DLQ after exactly $k$ attempts with idempotent requeue. \\
Ordering guarantee & 30\% & per-aggregate ordering holds under crash/replay injection without hurting cross-aggregate throughput. \\
Regression safety & 20\% & the original outbox suite still passes unmodified. \\
\bottomrule
\end{tabularx}
\end{center}
\textbf{Reference approach.} Harden the teaching relay: budget wrapper,
dead-letter table, \texttt{(aggregate\_id, seq)} ordering with consumer-side
gap handling. Stretch: lag metrics (unpublished count, oldest unpublished
age) with alert thresholds.

\subsubsection*{Capstone 15.4 --- Measure the Retry-Amplification Theorem [senior]}
\begin{center}\small
\begin{tabularx}{\textwidth}{l c X}
\toprule
\textbf{Criterion} & \textbf{Weight} & \textbf{Full marks when\ldots} \\
\midrule
Unbudgeted verification & 30\% & per-hop attempt counts match $a^k$ exactly (3, 9, 27; total 39). \\
Budgeted verification & 30\% & closed-form prediction for the budgeted case matches measurement. \\
Gap analysis & 20\% & a written paragraph explains any residual theory-vs-measurement gap. \\
Methodology & 20\% & offline, seeded, counting stubs; no wall-clock assertions. \\
\bottomrule
\end{tabularx}
\end{center}
\textbf{Reference approach.} Four-hop chain of counting stubs, each retrying
per its budget token; count attempts per hop and compare against
$a(a^n - 1)/(a - 1)$. Stretch: swap the stubs for the real ASGI chain and
show the theorem survives real HTTP machinery.

\subsubsection*{Capstone 15.5 --- Resilience Architecture for Fleet Telemetry [staff]}
\begin{center}\small
\begin{tabularx}{\textwidth}{l c X}
\toprule
\textbf{Criterion} & \textbf{Weight} & \textbf{Full marks when\ldots} \\
\midrule
Parameter justification & 25\% & every edge has SLO-derived parameters --- zero defaults anywhere. \\
Deadline correctness & 20\% & deadline budgets satisfy the formal rule on every path (ingestion and dashboard read). \\
Failure-mode coverage & 25\% & table has signal + countermeasure per scenario, including at least one \emph{degradation} (not just failure) case. \\
Operational program & 30\% & chaos-drill program includes abort criteria and rollback; mesh boundary list correctly assigns $\geq$3 concerns to the application. \\
\bottomrule
\end{tabularx}
\end{center}
\textbf{Reference approach.} Walk every edge of ingestion $\to$ normalizer
$\to$ enrichment $\to$ storage (plus the read path); derive timeouts,
retry budgets, and breaker parameters from the SLO arithmetic; tabulate
failure modes with signals and countermeasures \cite{beyer2016sre,
newman2021microservices}. Stretch: cost the design (retry/hedge load,
engineering hours, drill cost) against an incident cost model.

%% ============================================================================
\section{Chapter 16: API Security Engineering}
\label{sec:sol-ch16}

\subsection{Self-Assessment Model Answers}

\begin{enumerate}[label=\textbf{E16.\arabic*.}, leftmargin=2.8em]
  \item The fetcher must re-validate the resolved IP of \emph{every}
    redirect hop, not just the first URL: the two-302 chain to an internal
    host must be blocked at the final hop. With \texttt{max\_redirects=1}
    the same chain dies one hop earlier. Operationally, the first failure
    proves policy correctness end-to-end; the second proves defense in
    depth and leaves a clean audit signal of where the chain was cut
    \cite{owaspapisecurity}.
  \item Policy rows for \texttt{support}: read \texttt{notes} allowed, read
    \texttt{battery\_pct} denied, all writes denied. Tests: allowed read,
    denied field, denied write --- plus the union case: a principal that is
    both \texttt{support} and \texttt{owner} gets the union of permissions
    (and nothing beyond it).
  \item Moving claims mapping ahead of signature verification means attacker
    bytes become trusted claims before authenticity is established. The
    catching test mints an unsigned \texttt{alg=HS256} token with an
    attacker-chosen \texttt{sub} and asserts rejection --- verification
    order is a security invariant, not a style choice \cite{rfc7519}.
  \item \texttt{app.routes} on the vulnerable build lists debug/admin
    routes absent from its own OpenAPI document --- undocumented attack
    surface. This instantiates OWASP API9 (Improper Inventory Management):
    you cannot protect what your inventory does not know about
    \cite{owaspapisecurity}.
  \item Capture robot 1001's JSON per role on the hardened build and diff
    field sets against the read policy --- exact match required. On the
    vulnerable build the same diff counts the excess fields leaked
    (property-level authorization failure, OWASP API3).
\end{enumerate}

\subsubsection*{Deep-Dive Questions}

\begin{enumerate}[label=\textbf{D16.\arabic*.}, leftmargin=2.8em]
  \item \texttt{aud} enforcement is mandatory \emph{now} because the moment
    a second API exists (18 months later), tokens minted for API A replay
    cleanly against API B --- audience confusion. Retrofitting
    \texttt{aud} then is a breaking change for every client; enforcing it
    from day one is free.
  \item Dependent pair: a DTO allowlist and Pydantic \texttt{extra=forbid}
    --- both fail identically when a new sensitive field is added to the
    shared model (one omission defeats both; ceremony without independence).
    Independent pair: DTO allowlist \emph{plus} database-layer column
    permissions --- different mechanisms, different failure modes.
  \item Connecting by pinned IP breaks SNI and hostname verification: the
    certificate's SANs name hosts, not IPs, and a shared IP may serve the
    wrong vhost. Production SSRF guards resolve the tension by resolving
    $\to$ validating the IP $\to$ connecting with the original
    \texttt{Host}/SNI pinned, or by pushing egress through an allowlisting
    proxy that sees the hostname.
  \item Across two hops, \texttt{sub} must remain the original user while
    \texttt{act} accumulates each acting party (delegation chain). If hop
    two impersonates instead of delegates, the audit trail loses the agent
    --- logs can no longer distinguish what the service did on the user's
    behalf from what the user did.
  \item Invisible to an edge WAF: API1 (BOLA), API3 (property-level
    authorization), API6 (unrestricted sensitive business flows), API9
    (inventory) --- everything requiring object-ownership or business
    context, which no payload-inspecting appliance can possess
    \cite{owaspapisecurity}.
  \item A per-process limiter at twelve replicas admits twelve times the
    intended limit --- the limit silently scales with the fleet. The fix is
    Chapter~\ref{ch:rate_limiting}'s distributed architecture: a shared
    store (Redis cell/token bucket) so the limit is a property of the
    service, not of one process.
\end{enumerate}

\subsection{Capstone Grading Rubrics}

\subsubsection*{Capstone 16.1 --- BOLA Hunter [junior]}
\begin{center}\small
\begin{tabularx}{\textwidth}{l c X}
\toprule
\textbf{Criterion} & \textbf{Weight} & \textbf{Full marks when\ldots} \\
\midrule
Exploit demonstration & 30\% & the exploit script retrieves a foreign note on the vulnerable build. \\
Fix completeness & 30\% & the hardened build returns 403 for \emph{every} cross-owner access, proven by test. \\
Bidirectional tests & 25\% & pytest asserts both exploit-succeeds (vulnerable) and fix-holds (hardened), offline. \\
Conceptual understanding & 15\% & README explains why unguessable IDs alone are not the control. \\
\bottomrule
\end{tabularx}
\end{center}
\textbf{Reference approach.} Ownership column on notes; exploit walks
sequential IDs; fix enforces \texttt{owner\_id == caller} in the query
itself. Stretch: ETag conditional writes, and the UUIDv7 demonstration that
enumeration resistance is not authorization.

\subsubsection*{Capstone 16.2 --- Scope Forge \& Token Gauntlet [mid]}
\begin{center}\small
\begin{tabularx}{\textwidth}{l c X}
\toprule
\textbf{Criterion} & \textbf{Weight} & \textbf{Full marks when\ldots} \\
\midrule
Forgery battery & 35\% & all seven forgery tests (expired, wrong \texttt{aud}, wrong \texttt{iss}, \texttt{alg=none}, unknown \texttt{kid}, tampered payload, missing scope) yield 401/403 --- never 200. \\
Verification order & 25\% & provably no claim is read before signature verification. \\
Scope design & 25\% & the scope document justifies each scope's granularity. \\
Test hygiene & 15\% & deterministic offline token minting for tests. \\
\bottomrule
\end{tabularx}
\end{center}
\textbf{Reference approach.} Fixed pipeline: parse header $\to$ resolve key
by \texttt{kid} $\to$ verify signature $\to$ validate
\texttt{exp}/\texttt{iss}/\texttt{aud} $\to$ map scopes; the gauntlet is a
parameterized test. Stretch: token exchange to a second service with
narrowed scope/audience and clock-skew boundary tests \cite{rfc7519}.

\subsubsection*{Capstone 16.3 --- SSRF Bunker [mid]}
\begin{center}\small
\begin{tabularx}{\textwidth}{l c X}
\toprule
\textbf{Criterion} & \textbf{Weight} & \textbf{Full marks when\ldots} \\
\midrule
Exploit/fix symmetry & 30\% & all naive exploits (metadata IP, internal host, redirect chain) succeed pre-fix and are blocked post-fix in one pytest run. \\
Layer coverage & 25\% & DNS-rebinding and redirect re-validation tests included --- all seven layers exercised. \\
Documentation & 25\% & README layer-by-layer blame table maps each exploit to its catching control. \\
Determinism & 20\% & simulated adversary network; fully offline. \\
\bottomrule
\end{tabularx}
\end{center}
\textbf{Reference approach.} Apply the seven-layer gauntlet: scheme
allowlist, host allowlist, DNS resolution with IP validation, pinned
connect, redirect re-validation, response limits, egress policy. Stretch:
IPv6-mapped IPv4 bypass attempts and an independent egress-rule transport
wrapper.

\subsubsection*{Capstone 16.4 --- Threat-Model Pipeline [senior]}
\begin{center}\small
\begin{tabularx}{\textwidth}{l c X}
\toprule
\textbf{Criterion} & \textbf{Weight} & \textbf{Full marks when\ldots} \\
\midrule
Gate enforcement & 30\% & \texttt{--check} exits non-zero on undocumented/unmodeled endpoints and zero when worksheets are current. \\
Generator heuristics & 25\% & $\geq$3 new STRIDE heuristics, each unit-tested. \\
Abuse-case tests & 25\% & two abuse cases expressed as passing security tests. \\
CI integration & 20\% & gate wired into CI with documented exception/expiry handling. \\
\bottomrule
\end{tabularx}
\end{center}
\textbf{Reference approach.} Treat threat models as build artifacts: a
worksheet per endpoint, generated skeletons from the route table, staleness
detection by hashing the route surface. Stretch: attack-tree DOT output and
risk scoring by flag count and data sensitivity.

\subsubsection*{Capstone 16.5 --- Full Red Team / Blue Team Exercise [staff]}
\begin{center}\small
\begin{tabularx}{\textwidth}{l c X}
\toprule
\textbf{Criterion} & \textbf{Weight} & \textbf{Full marks when\ldots} \\
\midrule
Executable findings & 25\% & every red finding exists as an executable contract test --- no prose-only findings. \\
Cross-service exploit & 25\% & audience-confusion exploit demonstrated \emph{and} closed, with regression test. \\
Functional safety & 20\% & the functional suite stays green after all hardening. \\
Retro quality & 30\% & retro maps findings $\to$ OWASP items $\to$ detection signals $\to$ permanent controls; rollback plan and residual risk stated. \\
\bottomrule
\end{tabularx}
\end{center}
\textbf{Reference approach.} Two-sided exercise on the multi-service
topology: red exploits seeded flaws, blue fixes behind contract tests, and
the retro institutionalizes every finding as a permanent control or alert
\cite{owaspapisecurity}. Stretch: mTLS via a local CA showing the
token-confusion exploit's remaining path, and time-to-detect measurement on
a deliberate regression.

%% ============================================================================
\section{Chapter 17: Deployment \& Infrastructure}
\label{sec:sol-ch17}

\subsection{Self-Assessment Model Answers}

\emph{Numbering note: the chapter numbers its exercises ``Exercise
7.1--7.8'' (a template artifact); this appendix preserves that numbering.
The chapter's four deep-dive questions are unnumbered prose; they are
answered here as D17.1--D17.4 in order of appearance.}

\begin{enumerate}[label=\textbf{Exercise 7.\arabic*}, leftmargin=3.4em]
  \item With source \texttt{COPY} before dependency install, a one-line
    source edit invalidates every later layer --- the 200\,MB dependency
    layer rebuilds (and re-pushes) on \emph{every commit}. Correct order:
    copy manifests, install dependencies, then copy source; a source edit
    then rebuilds only the final thin layer. The cost quantified: 200\,MB
    $\times$ every build vs a few KB.
  \item With \texttt{timeout\_graceful\_shutdown=1}, in-flight requests are
    severed after one second: the verdict line counting severed/reset
    requests flips from zero to nonzero. The amputated act of the four-act
    shutdown is the \emph{drain} --- waiting for accepted work to finish.
  \item Boot 75\,s with \texttt{periodSeconds: 5} needs $75/5 = 15$
    failures tolerated; with 25\% headroom, \texttt{failureThreshold} $\ge
    19$ (say 20). Raising the \emph{liveness} failureThreshold instead is
    wrong: liveness exists to detect deadlocks, and a tolerant liveness
    probe will happily restart slow-booting pods mid-boot --- the
    CrashLoopBackOff the startup probe was invented to prevent.
  \item With latency p95 factor 1.20 the regressed fixture is rolled back at
    batch 1 by the latency gate. Setting $\delta = 0.02$ widens the
    intolerance band: fewer false rollbacks, but detection needs more
    evidence, so detection latency grows. Choosing $\delta$ is choosing how
    much error budget you are willing to burn \emph{while deciding} ---
    $\delta$ must stay well below the budget a full rollout would spend.
  \item With an 18-second boot (matching A), Scenario B survives; at 35
    seconds it crash-loops. The exact boundary is derivable from probe
    parameters alone: boot time must not exceed \texttt{initialDelaySeconds
    + periodSeconds $\times$ failureThreshold} of the startup probe ---
    beyond that, the kubelet kills the container before it finishes
    booting, forever.
  \item \texttt{preStop} sleep $\ge$ LB target-removal convergence: choose
    15\,s for a 12\,s bound. Uvicorn graceful timeout $\ge$ p99.9 request
    duration: choose 10--12\,s for a 9\,s bound.
    \texttt{terminationGracePeriodSeconds} $\ge$ preStop + drain + margin:
    choose 30--35\,s. Each number is justified by the mechanism it protects
    --- endpoint propagation, in-flight drain, and the kubelet's hard kill.
  \item Without the PDB, a routine node upgrade evicts \emph{all} replicas
    simultaneously (nothing enforces \texttt{minAvailable}); the drain
    completes before replacements schedule; the HPA sees no CPU signal to
    react to; the API is fully offline until new pods schedule, pull, and
    pass readiness --- an outage caused entirely by a routine, healthy
    operation.
  \item Method: compute the expected per-request log-likelihood-ratio
    increment under each hypothesis (baseline 0.4\% vs alternative 1.4\%
    with $\delta = 1$ pt), then $E[N_{\text{rollback}}] \approx
    \log((1-\beta)/\alpha)\,/\,E[\mathrm{LLR} \mid H_1]$ and
    $E[N_{\text{promote}}] \approx \log((1-\alpha)/\beta)\,/\,
    E[\mathrm{LLR} \mid H_0]$. Verify both against the simulator --- the
    point is that decision latency is a \emph{derived} quantity, not a
    guess.
\end{enumerate}

\subsubsection*{Deep-Dive Questions}

\begin{enumerate}[label=\textbf{D17.\arabic*.}, leftmargin=2.8em]
  \item Digest pinning guarantees every node pulls identical bits, but pull
    time is on the scale-out critical path: during a surge, heterogeneous
    pull times stagger readiness across the node pool exactly when capacity
    is most needed. Mitigations: pre-pulled images on the node image,
    registry mirrors near the pool, and surge headroom that does not depend
    on cold pulls.
  \item Ambient mesh removes sidecars: pod resource requests shrink and
    lose a container-start ordering concern; L4 moves to ztunnel, L7 policy
    to waypoints (changing where mTLS terminates and where authorization
    policy lives); drain and lifecycle hooks no longer coordinate with an
    injected proxy; observability and HPA metrics now come from mesh
    components outside the pod.
  \item In PromQL: compute canary-vs-baseline error-rate (or latency)
    ratios over rolling windows and express the gate as an alert expression
    with \texttt{for:} dwell --- the batch-file SPRT becomes a recording
    rule plus an abort condition the rollout controller watches.
  \item \texttt{PodLifecycleSleepAction}'s graduation history (alpha
    $\to$ beta $\to$ GA over multiple releases) is the cautionary tale:
    alpha features can change shape or vanish between releases, so
    lifecycle-correctness mechanisms (drain!) must never depend on them ---
    wait for GA, or implement the sleep portably.
\end{enumerate}

\subsection{Capstone Grading Rubrics}

\subsubsection*{Capstone 17.1 --- Harden the Image [junior]}
\begin{center}\small
\begin{tabularx}{\textwidth}{l c X}
\toprule
\textbf{Criterion} & \textbf{Weight} & \textbf{Full marks when\ldots} \\
\midrule
Size \& layering & 25\% & final image $\le$20\% of the naive size (or a documented reason); one-line source edit rebuilds only the source layer (build-log evidence). \\
Runtime hardening & 30\% & \texttt{docker inspect} shows a non-root user and exec-form entrypoint. \\
Hygiene & 25\% & no \texttt{.env}, \texttt{.git}, or test artifacts in any layer. \\
Reproducibility & 20\% & pinned bases/dependencies; clean-build instructions work offline. \\
\bottomrule
\end{tabularx}
\end{center}
\textbf{Reference approach.} Multi-stage build: deps from lockfile first,
source last, runtime stage minimal, non-root user, exec entrypoint
\cite{twelvefactor}. Stretch: slim vs distroless variants with the traded
operational capability (shell/debuggability) documented.

\subsubsection*{Capstone 17.2 --- Prove the Drain Under Load [mid]}
\begin{center}\small
\begin{tabularx}{\textwidth}{l c X}
\toprule
\textbf{Criterion} & \textbf{Weight} & \textbf{Full marks when\ldots} \\
\midrule
Verdict accounting & 30\% & three verdict classes reported per run reconcile exactly with attempted requests. \\
Contrast demonstration & 30\% & zero severed requests at production timeout, nonzero at timeout 0 --- both explained mechanistically. \\
CI gating & 25\% & pytest integration exits non-zero on any severed request. \\
Reproducibility & 15\% & seeded, offline, stable across runs. \\
\bottomrule
\end{tabularx}
\end{center}
\textbf{Reference approach.} Drive a concurrent traffic mix through the
shutdown demo; classify every request outcome; run the drill at both
timeout extremes. Stretch: a preStop-equivalent delay showing
refused-after-signal drop to zero.

\subsubsection*{Capstone 17.3 --- Manifest Review Board [mid]}
\begin{center}\small
\begin{tabularx}{\textwidth}{l c X}
\toprule
\textbf{Criterion} & \textbf{Weight} & \textbf{Full marks when\ldots} \\
\midrule
Fault detection & 35\% & the linter catches all seven injected faults with rule-specific messages; clean set passes. \\
Rule traceability & 25\% & every rule cites its chapter section. \\
Review realism & 25\% & minutes give a failure story per fault (blast radius, mechanism) --- not rule restatements. \\
Engineering & 15\% & YAML parsing, deterministic output, documented exit codes. \\
\bottomrule
\end{tabularx}
\end{center}
\textbf{Reference approach.} Encode the chapter's rules (probes, PDB,
resources, non-root, image pinning, \ldots) as linter rules over parsed
manifests; run the formal review; record minutes. Stretch: a second-reviewer
mode diffing two manifest revisions for safety-relevant changes only.

\subsubsection*{Capstone 17.4 --- Gate Tuner [senior]}
\begin{center}\small
\begin{tabularx}{\textwidth}{l c X}
\toprule
\textbf{Criterion} & \textbf{Weight} & \textbf{Full marks when\ldots} \\
\midrule
Sweep evidence & 30\% & the $\delta/\alpha/\beta$ sweep table shows the trade-off: tighter intolerance costs more requests per decision. \\
Promotion safety & 25\% & the $0.5\delta$ canary is promoted (never rolled back) in every sweep cell. \\
Theory-measurement agreement & 25\% & measured detection latency matches the LLR-drift derivation within one batch. \\
Policy artifact & 20\% & the policy page names $\delta$, $\alpha$, $\beta$, batch size, and rollback action with cited evidence. \\
\bottomrule
\end{tabularx}
\end{center}
\textbf{Reference approach.} Parameterize the SPRT gate; sweep; derive
expected drift per request; reconcile measurement with derivation. Stretch:
seasonal-baseline mode and an analysis of seasonality's effect on the SPRT
assumptions.

\subsubsection*{Capstone 17.5 --- Zero-Downtime Reference Architecture [staff]}
\begin{center}\small
\begin{tabularx}{\textwidth}{l c X}
\toprule
\textbf{Criterion} & \textbf{Weight} & \textbf{Full marks when\ldots} \\
\midrule
ADR quality & 25\% & every ADR names rejected alternatives and reversal cost. \\
Compliance decisions & 20\% & the payments path has an explicit TLS-termination decision with compliance rationale --- no plaintext last hop. \\
Gate calibration & 25\% & per-tier gate parameters come with simulator evidence they detect the stated regression sizes. \\
Operational completeness & 30\% & WebSocket drain budget derived from measured connection lifetimes; C17.3's linter finds zero violations in the template; rollback semantics defined per tier. \\
\bottomrule
\end{tabularx}
\end{center}
\textbf{Reference approach.} ADR set covering image supply, rollout,
drain, probes, and traffic policy across the 12-service, two-region
platform; every number traced to measurement or derivation
\cite{twelvefactor, beyer2016sre}. Stretch: the multi-region failover
half --- health-based DNS shifting and where canary gates live when
traffic itself moves.

%% ============================================================================
\section{Chapter 18: Observability \& SRE}
\label{sec:sol-ch18}

\subsection{Self-Assessment Model Answers}

\emph{Numbering note: the chapter uses a plain 1--10 list; numbering below
matches.}

\begin{enumerate}[leftmargin=2.2em]
  \item Availability SLI: fraction of polling URLs that eventually yield a
    completed plan (good = terminal success within a deadline). Latency
    SLI: time from POST to plan-ready, not the 202's return time. Freshness
    SLI: age of the plan's input data at completion. A 202 only means
    \emph{accepted} --- nothing about the finished plan --- so it can never
    be the availability SLI for the async result \cite{beyer2016sre}.
  \item Monthly budget $= 0.0005 \times 43{,}200 = 21.6$ min. Burned 14 by
    day 10: average burn rate $\approx 1.94\times$ (consumed fraction
    $14/21.6$ in one third of the month). Budget left: 7.6 min; at
    sustained $2\times$ burn the budget drains at $2 \times 21.6/30 =
    1.44$ min/day, covering $7.6/1.44 \approx 5.3$ days --- the latest a
    $2\times$ burn can start is about day 24.7 (day 24 to be safe).
  \item A 100\% error burst of $t$ minutes contributes error ratio $t/60$
    to the 1-hour window, so burn rate $= (t/60)/0.001 = 16.67\,t$. The
    14.4$\times$ fast page fires when $16.67\,t \ge 14.4$, i.e.\ $t \ge
    0.864$ min $\approx 52$~s --- the shortest full outage that pages.
    Verify against the chapter simulator.
  \item Choose boundaries so the 300\,ms SLO threshold is bracketed tightly
    --- e.g.\ a bucket boundary \emph{at} 300\,ms plus neighbors on both
    sides ($\ge$2 buckets straddling the threshold) --- because percentile
    estimates interpolate within a bucket. With a boundary exactly at the
    threshold, threshold classification is exact; otherwise the estimate
    error at the SLO is bounded by the containing bucket's relative width.
  \item Series count: routes $\times$ statuses $\times$ 4{,}000 tenants
    $\times$ 30 versions $\approx$ millions of series --- cardinality
    explosion. Keep \texttt{route} and \texttt{status}; reject
    \texttt{tenant\_id} outright and demote \texttt{client\_version} (or
    cap its values); rejected dimensions belong in logs and traces
    (exemplars), where cardinality is free \cite{opentelemetrydocs}.
  \item Inject the trace context into message metadata on publish; extract
    and continue the trace on consume. A consumer batching 50 messages into
    one span must use \emph{span links} to the 50 producers --- a single
    parent would falsely claim causal descent from one message and orphan
    the other 49 \cite{opentelemetrydocs}.
  \item ``Average p99 across instances'' understates pain two ways:
    (1)~instances with the worst tails are diluted by quiet instances;
    (2)~it weights instances equally regardless of request volume, so a hot
    instance serving most users counts once. Correct aggregation: merge the
    per-instance histograms and compute the global percentile over all
    requests --- never average percentiles.
  \item Head sampling keeps a uniform 5\% baseline; tail rules add: keep
    100\% of error traces (they cost money) and 100\% of traces exceeding
    the latency SLO (they cost customers), plus a sampled slice of
    slow-but-successful requests near the threshold. Guaranteed survivors:
    every error trace and every SLO-breaching trace.
  \item A growing gap between SERVER span start and the first child span is
    queueing/middleware time, not work time. Three candidates, each with a
    USE check: (1)~connection/thread-pool exhaustion --- pool utilization
    and queue depth (saturation); (2)~event-loop or CPU blockage --- CPU
    utilization/saturation; (3)~a slow synchronous middleware call (auth,
    config) --- that dependency's latency metric.
  \item Blameless rewrite: replace ``the engineer misconfigured the pool''
    with mechanism questions --- what validation allowed the
    misconfiguration to deploy, what signal would have caught it pre-impact,
    and why the change process permitted an unreviewed pool change. Action
    items: startup-time config validation, canary/announce for config
    changes, and a pool-saturation alert --- fixes aimed at the system that
    made the error easy \cite{beyer2016sre}.
\end{enumerate}

\subsection{Capstone Grading Rubrics}

\subsubsection*{Capstone 18.1 --- Trace Context Conformance Pack [junior]}
\begin{center}\small
\begin{tabularx}{\textwidth}{l c X}
\toprule
\textbf{Criterion} & \textbf{Weight} & \textbf{Full marks when\ldots} \\
\midrule
Corpus conformance & 35\% & all 25 hand-written corpus cases yield the expected verdict. \\
Property testing & 25\% & seeded round-trip property test passes 500/500. \\
Diagnostic quality & 25\% & every rejection message names the offending field. \\
Robustness & 15\% & parser only ever raises \texttt{ValueError} --- no crashes on any input. \\
\bottomrule
\end{tabularx}
\end{center}
\textbf{Reference approach.} Extend the \texttt{traceparent} parser into a
checker with a corpus-driven test table; round-trip encode/decode under a
seed. Stretch: fuzz with random truncations and case flips.

\subsubsection*{Capstone 18.2 --- RED Metrics for a Real Route Set [mid]}
\begin{center}\small
\begin{tabularx}{\textwidth}{l c X}
\toprule
\textbf{Criterion} & \textbf{Weight} & \textbf{Full marks when\ldots} \\
\midrule
Metric accuracy & 35\% & per-route rate/errors/duration table matches the seeded model within jitter. \\
Cardinality discipline & 25\% & series-count assertion passes, fails loudly when a \texttt{user\_id} label is added, and passes again when reverted. \\
Determinism & 25\% & one seed, no network, reproducible table. \\
Instrumentation quality & 15\% & per-route labeling consistent; no high-cardinality leaks. \\
\bottomrule
\end{tabularx}
\end{center}
\textbf{Reference approach.} Extend the instrumentation listing to label
per route; replay a seeded traffic model; assert the RED table. Stretch:
exemplars attaching one trace id per bucket.

\subsubsection*{Capstone 18.3 --- SLO Workbench [mid]}
\begin{center}\small
\begin{tabularx}{\textwidth}{l c X}
\toprule
\textbf{Criterion} & \textbf{Weight} & \textbf{Full marks when\ldots} \\
\midrule
Incident replay & 30\% & three scripted incidents produce reproducible alert timelines. \\
Flapping analysis & 25\% & the flapping incident is shown to reset-and-refire under the stock table, and your modification demonstrably addresses it. \\
Budget math & 25\% & the memo states the budget arithmetic for every policy row. \\
Policy defense & 20\% & one policy change defended with the generated data, not intuition \cite{beyer2016sre}. \\
\bottomrule
\end{tabularx}
\end{center}
\textbf{Reference approach.} Drive the burn-rate simulator with scripted
incident timelines; tabulate alert fire/reset times per policy; modify one
window/threshold and show the delta. Stretch: alert-resolution tracking
with mean-time-to-resolve per policy.

\subsubsection*{Capstone 18.4 --- Queue Propagation End to End [senior]}
\begin{center}\small
\begin{tabularx}{\textwidth}{l c X}
\toprule
\textbf{Criterion} & \textbf{Weight} & \textbf{Full marks when\ldots} \\
\midrule
Trace continuity & 30\% & the exported tree crosses the broker boundary with one trace id and correct parentage. \\
Batch semantics & 25\% & batch consumption uses span links, never a false single parent \cite{opentelemetrydocs}. \\
Legacy detection & 25\% & the legacy-consumer test (drops metadata) detects the orphaned span programmatically. \\
Determinism & 20\% & in-process fake broker; seeded; offline. \\
\bottomrule
\end{tabularx}
\end{center}
\textbf{Reference approach.} Inject context into message metadata on
publish; extract on consume; link batch spans to all producers; the orphan
detector flags spans whose parent never arrives. Stretch: baggage redaction
at the queue boundary with proof no baggage key survives.

\subsubsection*{Capstone 18.5 --- Incident Program in a Box [staff]}
\begin{center}\small
\begin{tabularx}{\textwidth}{l c X}
\toprule
\textbf{Criterion} & \textbf{Weight} & \textbf{Full marks when\ldots} \\
\midrule
Artifact completeness & 25\% & SLOs, runbooks, and dashboards exist for three services and cite one another. \\
Game-day rigor & 30\% & reproducible detection/mitigation/resolution times; a second pair reproduces the timeline within one simulated minute. \\
Postmortem quality & 25\% & blameless postmortem includes an alerting-topology change with a named owner. \\
Operational loop & 20\% & rollback criteria for every mitigation; alert fatigue considered; follow-ups tracked. \\
\bottomrule
\end{tabularx}
\end{center}
\textbf{Reference approach.} Define SLOs from user journeys; derive
burn-rate alerts; write runbooks that link to the dashboards that justify
them; run the scripted game day end to end \cite{beyer2016sre}. Stretch: a
budget-policy gate blocking deploys when $>$75\% of the error budget is
consumed.

%% ============================================================================
\section{Chapter 19: Testing \& Quality Engineering}
\label{sec:sol-ch19}

\subsection{Self-Assessment Model Answers}

\begin{enumerate}[label=\textbf{E-\arabic*.}, leftmargin=2.6em]
  \item Classify each of the fifteen tests into exactly one layer (unit /
    component-service / end-to-end). The service layer is typically
    over-weighted relative to its defect-catching power --- acceptable here
    because an API's contract surface \emph{is} where its defects live; the
    pyramid becomes a honeycomb deliberately, not by drift.
  \item Cheapest catchers: wrong SKU regex $\to$ unit property test over
    generated SKUs; renamed response field $\to$ contract test against the
    pact/schema; mispointed DB connection string $\to$ deployment smoke
    test (readiness against the real target); missing retry jitter $\to$
    unit test of the policy with a fake clock asserting decorrelated
    delays; p99 regression after a serializer change $\to$ microbenchmark
    gate in CI.
  \item The billing pact pins only \texttt{id} and \texttt{price\_cents}.
    Renaming \texttt{name} breaks no consumer (no pact references it);
    renaming \texttt{price\_cents} breaks exactly billing-service. This is
    \emph{usage sensitivity} --- a dimension schema diffing alone cannot
    express, because it cannot see which fields consumers actually read
    \cite{pactdocs}.
  \item Split the seeded provider state into two independent states and
    replay. If two interactions share one app instance, state leaks across
    interactions: execution order changes outcomes, failures misattribute,
    and passes can be false. A fresh app (or transactional reset) per
    interaction restores independence.
  \item Metamorphic relation: creating the same payload twice (minus
    \texttt{id}) must yield records equal up to the \texttt{id} field. It
    is stronger than any example assertion because it quantifies over
    \emph{all} generated payloads --- one counterexample anywhere in the
    generated space fails the property.
  \item In the closed model, clients wait for completion before issuing the
    next request, so offered load self-throttles exactly when the server
    saturates: the queueing that defines the knee is omitted
    (coordinated omission), the measured knee shifts right, and capacity is
    overstated. The open-loop 400 rps sweep lies less because it keeps
    offering load the server cannot absorb.
  \item Steady-state hypothesis: p99 latency and error rate of the shared
    staging environment remain within baseline bands under the latency-fault
    injection. Abort criteria: error budget burn $>$2$\times$, p99 beyond
    the band, or any page. Telemetry proving the experiment stayed in
    bounds: canary-cohort vs control metrics plus blast-radius scoping
    labels on every injected fault.
  \item A defensible 4-minute split: unit tests $\sim$60\,s, contract
    verification $\sim$90\,s, schema-diff/lint $\sim$30\,s, a selected
    service-test slice $\sim$60\,s. Moved to nightly: fuzz/property
    campaigns and load sweeps. The evidence lost: long-run flakiness and
    performance regressions surface a day later --- accepted knowingly,
    with the nightly results gated on release rather than on PR.
\end{enumerate}

\subsection{Capstone Grading Rubrics}

\subsubsection*{Capstone 19.1 --- A Second Consumer [junior]}
\begin{center}\small
\begin{tabularx}{\textwidth}{l c X}
\toprule
\textbf{Criterion} & \textbf{Weight} & \textbf{Full marks when\ldots} \\
\midrule
Pact health & 30\% & both pact files verify green against the healthy app \cite{pactdocs}. \\
Selective sensitivity & 35\% & the \texttt{sku} rename fails exactly one pact while the \texttt{name} rename fails none. \\
Diagnostic quality & 20\% & reports name the failing field with a JSONPath-style location. \\
Reproducibility & 15\% & deterministic, offline, seeded provider states. \\
\bottomrule
\end{tabularx}
\end{center}
\textbf{Reference approach.} Author the billing pact from real usage; run
both renames as experiments and record which pact fails. Stretch: a
\texttt{can-i-deploy(provider\_version)} function consulting a local JSON
``broker''.

\subsubsection*{Capstone 19.2 --- Property Suite for a New Resource [mid]}
\begin{center}\small
\begin{tabularx}{\textwidth}{l c X}
\toprule
\textbf{Criterion} & \textbf{Weight} & \textbf{Full marks when\ldots} \\
\midrule
Bug detection & 30\% & both planted bugs (500 trigger; schema-violating 201) are found and shrunk to minimal reproducers. \\
Invariant coverage & 25\% & zero 5xx and zero non-problem 4xx asserted across the whole campaign. \\
Determinism & 25\% & derandomized (seeded) suites pass repeatably. \\
Resource modeling & 20\% & \texttt{/warehouses} lifecycle fully covered by generators. \\
\bottomrule
\end{tabularx}
\end{center}
\textbf{Reference approach.} Generators for valid and near-valid payloads;
properties over create/read/update; shrinking on failure. Stretch: a
rule-based stateful machine with a local model oracle.

\subsubsection*{Capstone 19.3 --- Record--Replay Cassette Library [mid]}
\begin{center}\small
\begin{tabularx}{\textwidth}{l c X}
\toprule
\textbf{Criterion} & \textbf{Weight} & \textbf{Full marks when\ldots} \\
\midrule
Transparency & 30\% & replay works with zero application-code changes (transport-level interception). \\
Secret hygiene & 30\% & no credential-like header ever reaches disk --- proven by test. \\
Staleness handling & 25\% & stale cassettes detected and reported by date or hash. \\
Matching correctness & 15\% & request matching on method+path+canonical body behaves per spec. \\
\bottomrule
\end{tabularx}
\end{center}
\textbf{Reference approach.} A VCR-style transport for \texttt{httpx}:
record mode writes sanitized interactions; replay mode matches and returns
them; a scrubber runs before persistence. Stretch: strict vs lax matching
modes.

\subsubsection*{Capstone 19.4 --- Capacity Report from a Knee Sweep [senior]}
\begin{center}\small
\begin{tabularx}{\textwidth}{l c X}
\toprule
\textbf{Criterion} & \textbf{Weight} & \textbf{Full marks when\ldots} \\
\midrule
Prediction-measurement agreement & 30\% & predicted and measured knees agree within a factor of two, with the discrepancy explained. \\
Little's Law cross-check & 25\% & concurrency at the knee reconciles with $C=\lambda L$. \\
Recommendation quality & 25\% & explicit headroom recommendation with the assumptions listed. \\
Methodology & 20\% & open-model sweep; no wall-clock assertions anywhere. \\
\bottomrule
\end{tabularx}
\end{center}
\textbf{Reference approach.} Sweep offered load in open model; locate the
knee where latency departs linearity; cross-check with Little's Law; report
with headroom. Stretch: add the coordinated-omission buggy scheduler
variant and quantify how much it overstates capacity.

\subsubsection*{Capstone 19.5 --- Quality-Gate Architecture for a Fleet [staff]}
\begin{center}\small
\begin{tabularx}{\textwidth}{l c X}
\toprule
\textbf{Criterion} & \textbf{Weight} & \textbf{Full marks when\ldots} \\
\midrule
Defect-catcher mapping & 25\% & every defect class has exactly one designated cheapest catcher --- no duplicated or orphaned classes. \\
Budget discipline & 25\% & the PR gate fits four minutes with written justification per cut. \\
Adoption metrics & 25\% & rollout plan includes measurable adoption (pact coverage of production traffic) and a self-check requirement per gate. \\
Operational sustainability & 25\% & gate ownership, exception process, and rollback of a bad gate defined; flaky-gate handling stated. \\
\bottomrule
\end{tabularx}
\end{center}
\textbf{Reference approach.} Build the defect-class inventory first; assign
each class its cheapest layer; budget the PR pipeline by measurement;
instrument the gates themselves. Stretch: scheduled mutation testing of the
gates as a meta-gate.

%% ============================================================================
\section{Chapter 20: Lifecycle, Governance \& Developer Experience}
\label{sec:sol-ch20}

\subsection{Self-Assessment Model Answers}

\emph{Numbering note: the chapter uses a plain 1--10 list; numbering below
matches.}

\begin{enumerate}[leftmargin=2.2em]
  \item (a)~PR adding an endpoint to a published spec: no state transition
    if the spec is still draft; if published, the change requires a new
    minor/major version --- evidence: the spec diff plus the compatibility
    report. (b)~security flag on an unauthenticated endpoint:
    published $\to$ restricted/remediation workflow --- evidence: the
    threat-model worksheet. (c)~brownout revealing an unknown consumer
    three weeks before the kill date: deprecated $\to$ sunset-blocked
    (extension) --- evidence: access logs identifying the consumer.
  \item The rule needs a \emph{function/custom} kind: ``every POST
    operation documents a 201 or 202 response'' is an any-of-keys
    existential that schema-style rule kinds cannot express --- that gap
    (no anyOf-keys kind) is the finding; the rule walks
    \texttt{paths.*.post.responses} and asserts the intersection with
    \texttt{\{201, 202\}} is non-empty.
  \item With 40 APIs, \textbf{error}: the fix cost is low, humans absorb
    the findings, and the standard actually binds. With 400 APIs,
    \textbf{warn} plus a ratchet (error only on new or changed files):
    error-everywhere creates thousands of findings, alert fatigue, and a
    bypass culture that destroys the gate's authority --- governance that
    only says ``no'' gets routed around (Chapter~\ref{ch:lifecycle}).
  \item Hypotheses for the 73.3\% key$\to$first-call conversion with a
    194-minute median gap: (a)~\emph{onboarding friction} --- instrument
    step-level funnel events with timestamps to find the drop-off step;
    (b)~\emph{first-call failure} (docs mismatch, TLS, auth) --- instrument
    first-call outcomes with error reason codes. The instrument, not the
    fix, is the deliverable.
  \item SLO budget: $0.1\%$ allowed, $0.07\%$ consumed (99.93\% measured
    vs 99.9\% SLO) $\to$ 0.03\% remains. SLA budget: $0.5\%$ allowed
    $\to$ 0.43\% remains. A change spending another 0.05\% exceeds the
    remaining \emph{SLO} budget, so the SLO governs: defer the change,
    shrink its blast radius, or take a formal exception --- the stricter
    internal contract exists precisely to protect the looser external one.
  \item The ADR paragraph names: scope (the one legacy partner endpoint),
    reason (contracted consumer, migration cost), expiry (six months, a
    date, not a condition), owner, and removal criteria. It is time-boxed
    and auditable because the lint waiver carries an \texttt{expires} field
    that CI enforces --- the exemption fails the build the day it lapses.
  \item Generated docs derive from the spec, so if the hand-written
    quickstart disagrees, the quickstart is wrong --- unless the spec
    itself is wrong, in which case both are and the spec fix regenerates
    the docs. The process gap: the quickstart is never tested against the
    spec. Preventing lint rule: a docs-snippet conformance check (code
    samples executed or schema-checked in CI).
  \item Standard cadence rows: announce (with date) $\to$
    Deprecation/Sunset headers $\to$ staged brownouts $\to$ kill
    \cite{rfc8594}. For the contracted enterprise consumer the changed
    rows: extended notice per contract, named-contact outreach instead of
    broadcast channels, a documented extension clause, and brownout
    exemptions --- the twelve internal consumers follow the standard
    schedule with owner-assigned migration tasks.
  \item Telemetry query: join gateway access logs (path, method, consumer
    id, volume) against the catalog registry; alert on any path with
    traffic above a threshold that has no registry entry. First-24-hours
    runbook: identify the owning team via deploy logs and request patterns,
    register or block the endpoint deliberately, and assess what data the
    shadow surface exposed.
  \item The honest slide: adoption (teams/APIs onboarded to paved paths),
    TTFC reduction (days $\to$ hours, from funnel telemetry), and incident
    reduction (P1 count and MTTR delta priced at incident cost). Weakest
    assumption: counterfactual attribution --- teams that adopt platform
    tooling were likely the better-instrumented teams anyway (selection
    bias); say so on the slide.
\end{enumerate}

\subsection{Capstone Grading Rubrics}

\subsubsection*{Capstone 20.1 --- Three New Rules for the Fleet [junior]}
\begin{center}\small
\begin{tabularx}{\textwidth}{l c X}
\toprule
\textbf{Criterion} & \textbf{Weight} & \textbf{Full marks when\ldots} \\
\midrule
Rule efficacy & 35\% & compliant specs still exit 0; the extended fixture exits 1 with all three new rule IDs firing. \\
Diagnostic quality & 25\% & every violation carries a JSON path to the offender. \\
Engineering & 25\% & new rule kinds unit-tested, including a bad-ruleset exit-2 test. \\
Convention fit & 15\% & rules follow the existing engine's kinds and style. \\
\bottomrule
\end{tabularx}
\end{center}
\textbf{Reference approach.} Extend the mini-Spectral engine with one rule
per gap; test each rule positive and negative. Stretch:
\texttt{x-lint-ignore} suppression comments honored and logged as warnings.

\subsubsection*{Capstone 20.2 --- Docs-as-Code Pipeline [mid]}
\begin{center}\small
\begin{tabularx}{\textwidth}{l c X}
\toprule
\textbf{Criterion} & \textbf{Weight} & \textbf{Full marks when\ldots} \\
\midrule
Single-command pipeline & 30\% & lint $\to$ docs $\to$ catalog for all specs in one run; non-compliant specs block the run. \\
Staleness gate & 30\% & editing a spec without regenerating docs fails CI. \\
Determinism & 25\% & fully offline, byte-deterministic output. \\
Artifact quality & 15\% & generated docs are navigable and the catalog is complete. \\
\bottomrule
\end{tabularx}
\end{center}
\textbf{Reference approach.} Chain the chapter's linter, doc generator, and
catalog builder in one script; the staleness check hashes spec inputs into
the generated output and compares in CI. Stretch: a stdlib Markdown link
checker over the generated docs.

\subsubsection*{Capstone 20.3 --- TTFC Instrumentation Pack [mid]}
\begin{center}\small
\begin{tabularx}{\textwidth}{l c X}
\toprule
\textbf{Criterion} & \textbf{Weight} & \textbf{Full marks when\ldots} \\
\midrule
Funnel integrity & 30\% & per-channel funnels show conversion deltas with counts --- never a percentage without $n$. \\
Cohort determinism & 25\% & weekly cohort table deterministic and sorted; malformed events counted-and-skipped, never fatal. \\
Decision orientation & 25\% & the experiment report names one metric, one funnel step, one falsifiable threshold. \\
Small-sample honesty & 20\% & small-$n$ cohorts flagged so noise is not over-read. \\
\bottomrule
\end{tabularx}
\end{center}
\textbf{Reference approach.} Extend the funnel analyzer from report to
decision tool: per-channel breakdowns, cohort tables, and an experiment
template. Stretch: the $n < 10$ flag.

\subsubsection*{Capstone 20.4 --- Sunset Campaign for Fleet API v1 [senior]}
\begin{center}\small
\begin{tabularx}{\textwidth}{l c X}
\toprule
\textbf{Criterion} & \textbf{Weight} & \textbf{Full marks when\ldots} \\
\midrule
Projection accuracy & 30\% & the model flags in advance which consumers will miss the kill date under the stated migration velocity. \\
Policy differentiation & 25\% & internal vs external consumers follow the chapter's differentiated policy table. \\
Campaign mechanics & 25\% & headers, brownouts, and communications scheduled; dates are explicit CLI inputs. \\
CI semantics & 20\% & exit codes remain CI-meaningful throughout the campaign. \\
\bottomrule
\end{tabularx}
\end{center}
\textbf{Reference approach.} Replay a heterogeneous consumer log; model
migration velocity per consumer; schedule brownouts against the projection
\cite{rfc8594}. Stretch: one extension request and its effect on the
brownout schedule for non-extended consumers.

\subsubsection*{Capstone 20.5 --- Platform ROI Scorecard \& Governance Operating Model [staff]}
\begin{center}\small
\begin{tabularx}{\textwidth}{l c X}
\toprule
\textbf{Criterion} & \textbf{Weight} & \textbf{Full marks when\ldots} \\
\midrule
Reproducible numbers & 30\% & every scorecard figure traces to a reproducible computation from synthetic inputs. \\
Narrative honesty & 25\% & the ROI narrative states its weakest assumption explicitly. \\
Operating model & 25\% & two pages naming decision rights for every governance event (exceptions, sunsets, standard changes). \\
Quarterly cadence & 20\% & the review is runnable end-to-end as a quarterly ritual with owners and artifacts. \\
\bottomrule
\end{tabularx}
\end{center}
\textbf{Reference approach.} Build the scorecard from the funnel, incident,
and adoption data pipelines; write the operating model as a decision-rights
table; rehearse the quarterly review once. Stretch: a counterfactual
section estimating incident cost avoided, with the selection-bias caveat
addressed head-on.

%% ============================================================================
\section{Chapter 21: Synthesis \& Future Directions}
\label{sec:sol-ch21}

\subsection{Self-Assessment Model Answers}

\subsubsection*{Concept Review}

\begin{enumerate}[leftmargin=2.2em]
  \item Delivery model comes first because it prunes whole protocol
    families: a streaming requirement eliminates plain request/response
    options before consumer type is even examined. Reversed-order failure:
    choosing ``internal consumers $\to$ gRPC'' for a fleet telemetry
    \emph{push} leads to bolting streaming semantics onto the wrong
    conversation --- delivery shape dominates consumer identity.
  \item Any cell is a tendency; a legitimate exception to ``gRPC is for
    internal service-to-service'': several major public platforms expose
    gRPC (or gRPC+transcoding) externally with great success --- the
    tendency encodes tooling/browser friction, not a law
    \cite{grpcdocs}.
  \item The outbox relay reads from the database because a service cannot
    atomically write its database \emph{and} publish to the broker: crash
    between the two writes yields either a lost event (DB committed,
    publish never sent) or a phantom event (publish sent, DB rolled back).
    The outbox makes the database the single atomic commit point and the
    relay a replayable follower --- the dual-write anomaly is designed out
    \cite{kleppmann2017ddia}.
  \item \texttt{ETag} + \texttt{Cache-Control} (conditional GETs and shared
    caching) make a public REST catalog dramatically cheaper: repeat reads
    cost headers only or hit a CDN \cite{rfc9110}. GraphQL endpoints do not
    benefit equally because queries are POSTed bodies to one endpoint ---
    HTTP caches have nothing to key on without persisted-query-over-GET
    gymnastics.
  \item O-002 (a ProblemDetails reference on every 4xx/5xx) prevents the
    \emph{contract/observability} class of failure: unstructured errors
    consumers cannot parse programmatically. It protects consumers (stable
    error handling) and on-call (classifiable failures) alike
    \cite{rfc9457}.
  \item Deterministic tie-breaking makes the advisor a pure function of its
    inputs --- ``either is fine'' is an unstable output that breaks the
    snapshot-style tests asserting byte-identical recommendations per
    answer set; determinism is what makes the tool testable and its output
    diffable in review.
  \item MCP tool schemas generated from the OpenAPI source of truth cannot
    drift from the human-facing contract; hand-maintained tool schemas rot
    silently, and an agent acting on a stale schema produces confident,
    wrong calls --- the same contract-drift failure class as any codegen
    target, now with an autonomous consumer.
  \item The case-study reasoning: the frontiers table prices edge compute
    as immature tooling plus new operational burden (observability,
    deploys, data residency), while the order path's bottleneck was
    measurable elsewhere; the marginal latency win did not pay the
    migration risk for a small team --- being early is a cost, not a
    virtue.
\end{enumerate}

\subsubsection*{Architecture-Design Drills}

\begin{enumerate}[label=\textbf{Drill \arabic*.}, leftmargin=2.8em]
  \item \emph{Hospital device monitoring.} Profiles: internal, regulated,
    no public internet. Surfaces: device ingestion over gRPC streaming
    (efficient, schema-strict); nurse dashboards over SSE (freshness
    $<$2\,s, proxy-friendly, resumable); auditor reads over REST with
    \texttt{Cache-Control: max-age=86400} and ETags. Top failure modes:
    device-data loss on disconnect (buffer at the edge + outbox on
    ingestion) and dashboard staleness (heartbeats, Last-Event-ID resume,
    freshness SLI).
  \item \emph{Marketplace payout API.} Profiles: public, financial,
    exactly-once at the boundary. Surfaces: REST POST with
    \texttt{Idempotency-Key} plus unique constraints (Chapter
    \ref{ch:idempotency}); async bank confirmation via signed webhooks
    with DLQ (Chapter \ref{ch:async_apis}); RFC~9457 errors (Chapter
    \ref{ch:problem_details}). Top failure modes: duplicate payout
    (keys + constraints + replayed stored responses) and slow/flaky
    partner banks (deadlines, circuit breaker, saga with compensation);
    bursts handled by a queue-backed admission tier with rate limits
    (Chapter \ref{ch:rate_limiting}).
  \item \emph{Multi-tenant analytics backend.} Defend GraphQL: forty entity
    types with unknowable, user-configured payload shapes are its home
    territory. Guardrails either way: persisted queries for production
    traffic, depth/complexity limits, per-tenant cost-based rate limiting,
    per-request DataLoaders, and APQ-over-GET for the metered mobile
    variant (Chapter \ref{ch:graphql}).
  \item \emph{Fleet firmware orchestration.} Some networks forbid inbound
    device connections, so server-push alone is disqualified: choose
    \textbf{device-initiated long-polling with ETags/signed cursors} ---
    resumable, NAT-friendly, cacheable; offer SSE where inbound is
    permitted. Commands signed and expiring; jittered poll intervals to
    avoid a post-outage thundering herd of 50{,}000 devices. Failure modes:
    herd after outage (jitter + backoff), stale/replayed commands
    (signatures + expiry + per-device cursors).
  \item \emph{Your own platform.} Apply the six criteria surface by
    surface; where the card disagrees with history, decide whether the
    card missed a constraint actually present (data residency, team
    skills, regulatory) --- extend the card --- or whether the historical
    choice was fashion-driven --- document the drift. The deliverable is
    the annotated disagreement list, not agreement.
\end{enumerate}

\subsection{Capstone Grading Rubrics}

\emph{Numbering note: the chapter labels its capstones ``Project 1--5'';
they are cited here as C21.1--C21.5.}

\subsubsection*{Capstone 21.1 --- Extend the Decision Advisor [junior]}
\begin{center}\small
\begin{tabularx}{\textwidth}{l c X}
\toprule
\textbf{Criterion} & \textbf{Weight} & \textbf{Full marks when\ldots} \\
\midrule
Questionnaire robustness & 25\% & interactive mode runs offline and rejects invalid answers clearly. \\
Regression safety & 25\% & all pre-existing advisor tests pass unmodified. \\
New-criterion correctness & 25\% & data-residency tests cover scoring and \texttt{--explain} mentions residency when it moved a score. \\
Explainability & 25\% & \texttt{--explain} output traces every recommendation to criterion scores. \\
\bottomrule
\end{tabularx}
\end{center}
\textbf{Reference approach.} Add the seventh criterion to the scoring
table, thread it through the rationale trace, and wrap the questionnaire
around the existing pure scoring function. Stretch: JSON score output for
CI and a confidence metric from the winner's margin.

\subsubsection*{Capstone 21.2 --- Governance Toolchain [mid]}
\begin{center}\small
\begin{tabularx}{\textwidth}{l c X}
\toprule
\textbf{Criterion} & \textbf{Weight} & \textbf{Full marks when\ldots} \\
\midrule
Rule coverage & 30\% & all five original rule families plus three new rules fire on failing fixtures and pass on clean ones, with positive/negative tests each. \\
Reporting & 25\% & JSON report parses and round-trips deterministically; fix-hints attached. \\
Changed-only mode & 25\% & CI mode ignores pre-existing violations in untouched files. \\
Engineering & 20\% & offline, deterministic, documented exit codes. \\
\bottomrule
\end{tabularx}
\end{center}
\textbf{Reference approach.} Grow the conformance checker: rule registry,
finding model with fix hints, diff-aware file selection. Stretch: severity
levels with per-rule policy-file overrides and SARIF output.

\subsubsection*{Capstone 21.3 --- Multi-Protocol Read/Write Split [senior]}
\begin{center}\small
\begin{tabularx}{\textwidth}{l c X}
\toprule
\textbf{Criterion} & \textbf{Weight} & \textbf{Full marks when\ldots} \\
\midrule
Write-path discipline & 25\% & REST writes return RFC~9457 errors and honor idempotency keys \cite{rfc9457}. \\
Read-path consistency & 25\% & GraphQL reads never expose uncommitted state. \\
Concurrency proof & 25\% & the no-overselling concurrency test passes deterministically in 100 runs. \\
Contract governance & 25\% & OpenAPI and GraphQL SDL artifacts both pass the C21.2 checker. \\
\bottomrule
\end{tabularx}
\end{center}
\textbf{Reference approach.} One shared store behind three surfaces:
transactional REST writes, GraphQL read models with DataLoaders, gRPC for
internal coordination; consistency proven by the seeded concurrency
harness. Stretch: CDC-style lag metric with a freshness SLO and a
broker-outage lossless-recovery drill.

\subsubsection*{Capstone 21.4 --- Webhook Notification Pipeline with Outbox Replay [senior]}
\begin{center}\small
\begin{tabularx}{\textwidth}{l c X}
\toprule
\textbf{Criterion} & \textbf{Weight} & \textbf{Full marks when\ldots} \\
\midrule
Losslessness & 30\% & no event lost across the simulated outage, proven by test (outbox $\to$ relay $\to$ delivery log). \\
Security & 25\% & signature failures and stale timestamps rejected with 401/422 problem details. \\
Replay fidelity & 25\% & replay re-sends exactly the logged payload --- byte-identical. \\
Contract \& fallback & 20\% & AsyncAPI spec passes the conformance checker; SSE fallback demonstrated \cite{asyncapi}. \\
\bottomrule
\end{tabularx}
\end{center}
\textbf{Reference approach.} Compose the book's reference pieces: outbox
relay, HMAC signing, delivery log, operator replay endpoint, AsyncAPI
documentation, SSE fallback for push-averse consumers. Stretch: DLQ with
operator triage endpoint and a delivery-latency histogram metric.

\subsubsection*{Capstone 21.5 --- The Grand Capstone: Northwind Robotics API Platform [staff]}
\begin{center}\small
\begin{tabularx}{\textwidth}{l c X}
\toprule
\textbf{Criterion} & \textbf{Weight} & \textbf{Full marks when\ldots} \\
\midrule
Surface completeness & 25\% & all five surfaces respond per spec in one offline in-process test run. \\
Governance & 20\% & the conformance checker passes on the full manifests folder. \\
Measured SLOs & 20\% & SLOs measured green against a recorded load test --- asserted, not claimed \cite{beyer2016sre}. \\
Resilience proof & 20\% & the chaos drill (kill inventory mid-order) produces the designed failure behavior with zero data loss. \\
Coherence \& score & 15\% & rubric score $\ge$16/20 with no dimension below 3; rollback and observability designed in, not bolted on. \\
\bottomrule
\end{tabularx}
\end{center}
\textbf{Reference approach.} Build in the four milestones (contracts,
surfaces, resilience, governance), each milestone shippable; reuse every
prior capstone's hardened component rather than rewriting; the final system
is graded as a \emph{platform} --- coherent contracts, measured SLOs,
rehearsed failure, and a working governance loop
\cite{newman2021microservices, kleppmann2017ddia}. Stretch: canary
simulation with traffic splitting, an MCP tool surface generated from the
OpenAPI spec, and a regional-affinity design document for EU expansion.

